\documentclass[11pt]{amsart}
\IfFileExists{lmodern.sty}%
  {\usepackage[T1]{fontenc}\usepackage{lmodern}\usepackage{microtype}}%
  {\usepackage[expansion=false]{microtype}}
\usepackage{amssymb,mathtools}
\usepackage{booktabs,array}
\usepackage{enumitem}
\emergencystretch=2.5em
\usepackage[colorlinks=true,linkcolor=blue,citecolor=blue,urlcolor=blue]{hyperref}

\newcommand{\FormalizationURL}{https://github.com/slitsey/catalan-irrationality-measure}
\newcommand{\FormalizationLink}{\url{\FormalizationURL}}

% Numerical constants: decimal truncations (or outward roundings) of rational numbers of the formal proof.
\newcommand{\Mustar}{269.78696}
\newcommand{\Rho}{278.376313578}
\newcommand{\Ctwo}{3.596893}
\newcommand{\Sigm}{0.33871}
\newcommand{\Cost}{3.935603}
\newcommand{\Psiw}{50.111339578}
\newcommand{\Czero}{33.052}
\newcommand{\Dcirc}{3.2506}
\newcommand{\Utarget}{2.255858836}
\newcommand{\Gain}{24.64}
\newcommand{\Uzero}{2.2631834}
\newcommand{\NormCap}{0.77363}
\newcommand{\SupX}{-0.96450499}
\newcommand{\SupY}{-1.57499400}
\newcommand{\FieldCap}{-2.53949899}
\newcommand{\ResCap}{-0.0631062}
\newcommand{\EtaTot}{54.68745}
\newcommand{\LamCrit}{0.9699101}
\newcommand{\Kcoef}{160}
\newcommand{\Ngeo}{75}
\newcommand{\Ncomplex}{41}
\newcommand{\Mmargin}{2.986\cdot10^{-6}}
\newcommand{\Decay}{6447.4755}
\newcommand{\DecayOA}{6422.8355}
\newcommand{\EtaExact}{54.68853}
\newcommand{\PayL}{53.0000596}
\newcommand{\PayR}{53.0419104}
\newcommand{\AlphaOne}{51/500}
\newcommand{\Chi}{22.06076}
\newcommand{\ChiLean}{22.0607}
\newcommand{\LamDen}{6392.7870}

\newtheorem{theorem}{Theorem}[section]
\newtheorem{lemma}[theorem]{Lemma}
\newtheorem{proposition}[theorem]{Proposition}
\newtheorem{corollary}[theorem]{Corollary}
\theoremstyle{definition}
\newtheorem{definition}[theorem]{Definition}
\theoremstyle{remark}
\newtheorem{remark}[theorem]{Remark}

\numberwithin{equation}{section}
\setcounter{tocdepth}{1}

\newcommand{\Z}{\mathbb Z}
\newcommand{\Q}{\mathbb Q}
\newcommand{\R}{\mathbb R}
\newcommand{\N}{\mathbb N}
\newcommand{\C}{\mathbb C}
\newcommand{\F}{\mathbb F}
\newcommand{\rank}{\operatorname{rank}}
\newcommand{\cV}{\mathcal V}
\newcommand{\cR}{\mathcal R}
\newcommand{\cA}{\mathcal A}
\newcommand{\cL}{\mathcal L}
\newcommand{\cK}{\mathcal K}
\newcommand{\cD}{\mathcal D}
\newcommand{\rev}{\operatorname{rev}}
\newcommand{\lean}[1]{\texttt{#1}}

\begin{document}

\title{An irrationality measure for Catalan's constant}

\author{Sean Litsey}
\email{sean.litsey.math@gmail.com}

\date{November 8, 2026}

\subjclass[2020]{Primary 11J82; Secondary 11J72, 15A15, 31A15, 68V20}

\keywords{Catalan's constant, irrationality measure, determinants, corank modulo $p$, Andr\'eief's identity,
logarithmic energy, Pick matrices, Nevanlinna--Pick interpolation, Gauss quadrature, Lean}

\begin{abstract}
Let $G=\sum_{j\ge0}(-1)^j(2j+1)^{-2}$ be Catalan's constant. We prove that $|G-a/b|>b^{-270}$ for all
integers $a$ and all sufficiently large integers $b$. Hence the irrationality exponent of $G$ is at most
$270$. The proof refines OpenAI's recent proof that $G$ is irrational. That proof compares two bounds for
a large determinant whose entries are linear forms in $1$ and $G$. The lower bound comes from the
denominators of the entries, and the upper bound from an integral representation. We improve both. On the
arithmetic side we show that the matrix of rational parts has large corank modulo $p$ for the primes in
explicit intervals, which lowers the denominators. On the analytic side we bound a determinant that the
energy method treats trivially. A phase inequality for Pick matrices, proved by Nevanlinna--Pick
interpolation, compares it with a determinant of a simpler kind. That determinant is bounded by a finite
Gram inequality and two logarithmic gases. Gauss quadrature at Chebyshev nodes of many short intervals
extends the upper bound to an exponentially small neighbourhood of $G$, as a measure requires. The proof is formalized in the Lean proof
assistant and kernel-checked.
\end{abstract}

\maketitle

\tableofcontents

\section{Introduction}\label{sec:intro}

\subsection{The result}
Catalan's constant
\[
 G=\sum_{j=0}^\infty\frac{(-1)^j}{(2j+1)^2}=L(2,\chi_{-4})=0.91596\,55941\,77219\ldots
\]
is the value at $s=2$ of the Dirichlet $L$-function of the nontrivial character modulo $4$. Values of
$L$-functions at integers where no formula in terms of $\pi$ is available are notoriously hard to study.
Ap\'ery proved $\zeta(3)$ irrational in 1978 \cite{Apery79}, and Calegari, Dimitrov and Tang proved
$L(2,\chi_{-3})$ irrational in 2024 \cite{CDT2024}. For $G$ the classical methods produce linear forms
that do not decay fast enough, and its irrationality remained open. In October 2026 OpenAI
\cite{OpenAI2026} proved that $G$ is irrational, with a new argument based on a large determinant, and
verified the proof in Lean. The present paper turns that argument into an irrationality measure.

The \emph{irrationality exponent} $\mu(\xi)$ of a real number $\xi$ is the infimum of the real numbers
$\mu$ such that $|\xi-a/b|>b^{-\mu}$ for all integers $a$ and all sufficiently large integers $b$. Every
irrational number has $\mu(\xi)\ge2$, and almost every real number has $\mu(\xi)=2$. An
\emph{irrationality measure} is an explicit upper bound for $\mu(\xi)$. The exponent does not change under
$\xi\mapsto r\xi+r'$ with $r,r'\in\Q$ and $r\ne0$.

\begin{theorem}\label{thm:main}
For all integers $a$ and all sufficiently large integers $b$,
\[
 \left|G-\frac ab\right|>b^{-270}.
\]
In particular $\mu(G)\le270$.
\end{theorem}

The proof uses the prime number theorem. The theorem and every result it depends on have been formalized
in the Lean~4 proof assistant on top of OpenAI's formalization, and the proof is kernel-checked. The formal proof has no unproved steps
and uses only the standard axioms (Section~\ref{sec:formal}). The use of AI in this work is disclosed at
the end of the paper.

\subsection{Background}
\emph{Catalan's constant.} Irrationality proofs for values of $L$-functions at integers usually follow
Ap\'ery's proof for $\zeta(3)$. One constructs linear forms $q_n\xi-p_n$ with integer coefficients that
tend to zero faster than the coefficients grow. If $|q_n\xi-p_n|=e^{-\alpha n+o(n)}$ and
$|q_n|=e^{\beta n+o(n)}$ with $\alpha>0$, then $\xi$ is irrational and $\mu(\xi)\le1+\beta/\alpha$.
Beukers \cite{Beukers79} obtained Ap\'ery's forms for $\zeta(2)$ and $\zeta(3)$ as integrals of
Legendre-type polynomials. Hata \cite{Hata93} showed how much is gained by dividing the forms by the
primes that divide all their coefficients, an idea that goes back to Chudnovsky. Rhin and Viola
\cite{RhinViola96,RhinViola01} found finite groups acting on families of such integrals. The invariance
yields a large common factor of the coefficients, and this \emph{arithmetic method} gave the best
measures for $\zeta(2)$ and $\zeta(3)$ for many years. Marcovecchio \cite{Marcovecchio09} applied it to
$\log2$, and Zeilberger and Zudilin \cite{ZeilbergerZudilin20} to $\pi$, using integrals of Salikhov
\cite{Salikhov08}. For Catalan's constant the analogous linear forms exist but do not decay fast enough.
Rivoal and Zudilin \cite{RivoalZudilin2003} proved that one of the numbers $\beta(2),\beta(4),\ldots,\beta(2m)$
is irrational for suitable $m$, where $\beta$ is the Dirichlet beta function. This does not single out
$G=\beta(2)$. Fischler \cite{Fischler20} and Kumar and Tosi \cite{KumarTosi26} obtained lower bounds for the number of
irrational values among $\beta(2),\beta(4),\ldots,\beta(2m)$. These bounds again do not single out $G$.
Zudilin \cite{Zudilin2003} found an Ap\'ery-like recurrence for $G$, whose solutions give
linear forms that tend to zero too slowly. Related constructions and their limits are described in
\cite{Zudilin02Cheb,Rivoal06,KrattRivoal08,KrattZudilin19,Zudilin19Survey}. Nesterenko \cite{Nesterenko16} developed a
general construction of such approximations and compared their quality. Rational approximations to $G$ with
$|G-p_n/q_n|\le q_n^{-0.52}$ follow from these works. Eskandari, Murty and Nemoto \cite{EMN2025} construct a mixed motive
with period $G$ and obtain linear forms in $1$ and $G$ from it. Eskandari \cite{Eskandari2026} raises the exponent $0.52$
to $0.62$. An exponent below $1$ does not imply irrationality.

\emph{Recent results.} Since September 2026 a series of new results in this area has appeared, several of
them obtained with AI systems. OpenAI proved $\mu(\pi)=2$, with a Lean formalization \cite{OpenAIPi2026}.
Liu, Jiang and Zhang proved $\mu(\log r)=2$ for every rational $r>0$ with $r\ne1$ \cite{LiuJiangZhang26}.
Chudnovsky adapted OpenAI's method to periods of elliptic curves and proved $\mu(\Omega)=2$ for the lemniscate period
$\Omega$ \cite{Chudnovsky26}.
Giraudin lowered the bound for $\mu(\pi^2)$ to $5.0193784$ with two hypergeometric constructions of Zudilin
\cite{Zudilin14} at new parameters \cite{Giraudin26}, and the bound for $\mu(\zeta(3))$ to $5.5138800$
\cite{Giraudin26z3}. Leibtag, Barkan, Kalisch, Weinbaum, Shalyt and Kaminer \cite{LBKWSK26} showed that the
classical constructions are one-dimensional trajectories in conservative matrix fields. Searching over other
trajectories, they obtained $\mu(\pi^2)<4.62675$ and $\mu(\zeta(3))<5.4197$, and further Ap\'ery-like
approximations to $G$, the best of which still has a negative approximation exponent. Prellberg
\cite{Prellberg26} proved that $\mathrm{Li}_2(-1/4)$, $\mathrm{Li}_2(1/5)$ and $\mathrm{Li}_2(-1/3)$ are
irrational. Fauzan proved that $\zeta(5)$ is irrational, using Hankel determinants \cite{Fauzan2026}, and Fasten verified the
irrationality statement in Lean \cite{Fasten26}. Fauzan's argument also gives an irrationality measure. Using its results,
Giraudin proved $\mu(\zeta(5))\le171$ \cite{Giraudin26z5}. Lai combined this method with simultaneous approximations to show that $1$, $\zeta(3)$ and
$\zeta(4)$ are linearly independent over $\Q$ \cite{Lai2026}. For $L(2,\chi_{-3})$ the effective holonomy
bounds of Calegari, Dimitrov and Tang give $\mu(L(2,\chi_{-3}))\le24781$ \cite[Thm.~5.1]{CDT2025}.

\emph{Determinant methods.} OpenAI's argument \cite{OpenAI2026} has a different shape. Instead of a
sequence of linear forms it uses one determinant $\Delta_N$ of size $n=48N$, whose entries are rational
linear combinations of $1$ and $G$. The entries are double integrals against the Cauchy kernel
$1/(1-ts)$. If $G$ were rational, the denominators of the entries would force a lower bound
$|\Delta_N|\ge e^{-(\Lambda+o(1))N^2}$ along $N=p$ prime. On the other hand, Andr\'eief's identity
\cite{Forrester2018Andreief} writes $\Delta_N$ as an integral of a product of determinants, and an estimate
for a discrete logarithmic energy gives $|\Delta_N|\le e^{-(D+o(1))N^2}$. OpenAI proves $D>\Lambda$, with a
margin of about $10^{-4}$ per unit of $n^2$, and proves that $\Delta_p\ne0$ by reduction modulo $p$.
Determinants of linear forms had been used for irrationality before, in Zudilin's criterion with Hankel
determinants \cite{Zudilin}.

The first irrationality measure for $G$ was obtained by Jha \cite{Jha2026}, who proved $\mu(G)\le10^8$.
Carreira \cite{Carreira2026} treated the five integer parameters of the construction as free, stated
conditions under which OpenAI's argument goes through, and proved $\mu(G)<2294$. He also expressed the best
constant of the energy estimate as the value of an equilibrium problem for two measures, and in
\cite{Carreira2026b} he records attempts to carry the method to $\beta(4)$. Krisman proved that one of
$\beta(4),\beta(6),\ldots,\beta(36)$ is irrational, with a Lean formalization \cite{Krisman26}. Bojnal \cite{Bojnal2026} and Xu
\cite{Xu2026} obtained $\mu(G)<2402$ and $\mu(G)\le727.03$ at other parameters. Xu's certificate is
numerical and not formalized. Davis \cite{Davis2026} reached $\mu(G)\le528.23288$ with new quadrature nodes, the zeros of $T_K(P(t))$
for a fixed polynomial $P$. The MultiAgency agent swarm regenerated $P$ and reached $\mu(G)\le528.2315$
\cite{MultiAgency2026}. Both results are formalized in Lean. Our nodes are different, but we use two
formal lemmas of \cite{Davis2026} on the weights of nodal polynomials in Section~\ref{sec:real}. The
other constructions are independent of the ideas of the present paper.
A preprint of Sun \cite{Sun2026} claimed the irrationality of $G$, but Wachs
\cite{Wachs2026} showed that its main estimate is not justified.

\subsection{The companion paper}
Catalan's constant is a difference of two values of the trigamma function $\psi'$. With
\[
 \Psi(s,t)=\psi'(s)-\psi'(t)=\sum_{n\ge0}\Bigl(\frac1{(n+s)^2}-\frac1{(n+t)^2}\Bigr)\qquad(0<s<t\le1)
\]
one has $G=\frac1{16}\Psi(\frac14,\frac34)$. The companion paper \cite{Litsey2026b} proves irrationality
measures for $\Psi(s,1)$ with $s=\frac k4,\frac k3,\frac k5,\frac k7,\frac12$, and for
$\Psi(\frac13,\frac23)=9L(2,\chi_{-3})$. For $\Psi(\frac14,1)$ and $\Psi(\frac34,1)$, which are rational multiples
of $G\pm5\pi^2/48$, and for $s=\frac k5,\frac k7$, these are the first proofs of irrationality that we know of. Its determinants are L\"owner matrices of digamma differences,
which are much smaller than OpenAI's. That method also applies to $(\frac14,\frac34)$, but there its
denominators outgrow its decay, so it gives no measure for $G$. The present paper therefore refines
OpenAI's construction. The two papers use the same criterion, the same names for the denominator exponent
$\Lambda$, the decay rate $D$ and the neighbourhood exponent $\rho$, and one formalization.

\subsection{What a measure requires}\label{sec:intro-need}
Write $\Delta_N(X)=\det(\cA^0_N+X\cL_N)$, so that OpenAI's determinant is $\Delta_N(G)$, and let $Q_N$ be
the integer polynomial obtained by clearing denominators. Suppose $|G-a/b|$ is small. Three properties of
$Q_N$ turn OpenAI's comparison into a contradiction.
\begin{enumerate}[label=(\roman*)]
\item \emph{Denominators.} The coefficients of $\Delta_N$ have a common denominator of size
 $e^{(\Lambda+o(1))N^2}$, and $\deg Q_N\le dN$.
\item \emph{Smallness on a neighbourhood.} $|\Delta_N(x)|\le e^{-(D+o(1))N^2}$ for every $x$ with
 $|x-G|\le e^{-\rho N}$, not only at $x=G$.
\item \emph{Nonvanishing at the point.} $Q_p(a/b)\ne0$ for a prime $p$ of size about $\log b$.
\end{enumerate}
OpenAI's proof gives nonvanishing only at $X=G$. We need it at every $a/b$ with $p^2\nmid b$, and we must
find such a prime without losing a power of $b$. Section~\ref{ssec:criterion} shows that the squares of
the larger primes can be recovered from the coefficients of $Q_p$.

\subsection{Outline of the proof}\label{sec:outline}
Suppose that $|G-a/b|<b^{-\mu}$ with $b$ large, and take a prime $p\approx\mu\log b/\rho$. Then
$|a/b-G|\le e^{-\rho p}$, so by (ii) $|\Delta_p(a/b)|\le e^{-(D+o(1))p^2}$. By (i) the number
$b^{dp}Q_p(a/b)$ is an integer, and by (iii) it is nonzero. Hence
\begin{align*}
 1\le\bigl|b^{dp}Q_p(a/b)\bigr|&\le b^{dp}\,e^{(\Lambda-D+o(1))p^2}\\
 &=\exp\bigl(dp\log b-(\eta+o(1))p^2\bigr),\qquad \eta=D-\Lambda .
\end{align*}
The right side is less than $1$ as soon as $\mu>d\rho/\eta$. So $\mu(G)\le d\rho/\eta$
(Proposition~\ref{prop:criterion} gives the precise form).

The three numbers come from different places. The degree $d$ comes from the construction
(Lemma~\ref{lem:degree}). The neighbourhood exponent $\rho$ comes from a quadrature formula for the
coefficient of $G$ (Section~\ref{sec:real}). The difference $\eta$ is a small difference of two large
numbers, $D\approx6447$ and $\Lambda\approx6393$ (Figure~\ref{fig:balance}). For OpenAI's construction with
our parameters, $D$ and $\Lambda$ differ by about $8$. We lower $\Lambda$ by $22.06$ through the corank of
a matrix modulo $p$ (Section~\ref{sec:content}), and raise $D$ by $24.64$ through a bound for a
determinant of row functions (Section~\ref{sec:signed}). The result is $\eta\approx54.69$, which is what
makes $d\rho/\eta$ as small as $270$.

\begin{figure}[ht]
\centering
\setlength{\unitlength}{0.15cm}
\begin{picture}(84,24)(-16,-6)
 \put(0,-1){\line(1,0){70}}
 \multiput(0,-1)(10,0){8}{\line(0,-1){0.8}}
 \put(0,-3.5){\makebox(0,0){$6380$}}
 \put(20,-3.5){\makebox(0,0){$6400$}}
 \put(40,-3.5){\makebox(0,0){$6420$}}
 \put(60,-3.5){\makebox(0,0){$6440$}}
 \qbezier(-1.5,-1.4)(-1,-0.9)(-0.5,-0.4)\qbezier(-0.9,-1.4)(-0.4,-0.9)(0.1,-0.4)
 \put(0,14){\framebox(42.8,3){}}
 \put(42.8,14){\framebox(24.6,3){$+24.6$}}
 \put(-1,15.5){\makebox(0,0)[r]{decay rate $D$}}
 \put(0,6){\framebox(12.8,3){}}
 \put(12.8,6){\dashbox{0.5}(22.1,3){$-22.1$}}
 \put(-1,7.5){\makebox(0,0)[r]{denominators $\Lambda$}}
 \put(12.8,3){\line(0,1){11}}
 \put(67.4,3){\line(0,1){11}}
 \qbezier(12.8,3.6)(40.1,3.6)(67.4,3.6)
 \qbezier(12.8,3.6)(13.4,3.85)(14.0,4.1)\qbezier(12.8,3.6)(13.4,3.35)(14.0,3.1)
 \qbezier(67.4,3.6)(66.8,3.85)(66.2,4.1)\qbezier(67.4,3.6)(66.8,3.35)(66.2,3.1)
 \put(40,1.6){\makebox(0,0){$\eta=D-\Lambda\approx54.7$}}
 \put(42.8,19.2){\makebox(0,0){$D_0=6422.8$}}
 \put(34.9,11.2){\makebox(0,0){$\Lambda_0=6414.8$}}
\end{picture}
\caption{The exponent balance for the parameters $(58,9,3,5,3)$, per unit of $N^2$. The axis is broken at
$6380$. OpenAI's estimates give the decay rate $D_0$ and the denominator exponent $\Lambda_0$, which
differ by about $8$. The bound of Section~\ref{sec:signed} adds $24.6$ to $D_0$, and the corank bounds
of Section~\ref{sec:content} remove $22.1$ from $\Lambda_0$.}\label{fig:balance}
\end{figure}

\subsection{New ideas}\label{sec:intro-ideas}
\begin{enumerate}[label=\arabic*.,wide]
\item \emph{Large corank modulo $p$.} Let $\nu_N(p)$ be the corank modulo $p$ of $p^2$ times the matrix of
rational parts. Each unit of corank lowers the denominator by a factor $p$
(Proposition~\ref{prop:valuation}). We show that the corank is of order $N$ for the primes in
$113$ intervals inside $(\frac{16}7N,\frac{53}2N]$. The kernel vectors are explicit on most of them. They come from
antisymmetric combinations of rows, a reflection formula for Chebyshev polynomials modulo $p$, and the
relation $(1-X)^p=1-X^p$. Together these lower the denominator exponent by $22.06$.

\item \emph{One counting principle.} All these kernels arise in the same way. One takes a space of row
combinations and imposes a few linear conditions that make every entry of the combined row vanish modulo
$p$. Then the corank is at least the dimension minus the number of conditions (Theorem~\ref{thm:count}).
The four closed formulas are instances. On the remaining intervals of primes we apply the principle to
every constant subspace at once (Theorem~\ref{thm:everyU}). Only upper bounds for ranks are then needed,
and these come from forced intersections of row spaces. No value of a moment enters.

\item \emph{A perturbation bound column by column.} For the neighbourhood of $G$ we write the coefficient of
$G$ by a quadrature formula with $53N$ nodes. The nodes are Chebyshev nodes of $106$ short intervals, and the
Lagrange polynomials are bounded by the Green functions of the intervals. The perturbed determinant is then
a sum of determinants in which some columns are replaced by values of the row polynomials at the nodes. We
bound the weight of each replaced column and its generalized Vandermonde factor together. A weighted form of
Hadamard's inequality with a parameter $\zeta=\frac34$ makes the dependence on the column cancel, and one
contour of radius $\frac34$ treats all columns at once. The part of the weight that depends on the node is
carried into the energy bound as an external field.

\item \emph{Trial sequences that move with the configuration.} The energy bound with prescribed columns uses
trial sequences that depend linearly on the fraction of prescribed columns. Their fields are affine in the
parameter, so the suprema are bounded by chords between two end points. A finite cover of the range of the
fraction by cells then reduces every case to exact polynomial inequalities.

\item \emph{A phase inequality for Pick matrices.} OpenAI's integrand contains a determinant
$\det\mathsf L$ of size $n$ that the energy method bounds by $e^{O(n\log n)}$. In the Pick matrix of the nodes the
positive nodes carry the weight $-5$. We show that $\det\mathsf L$ loses only a factor $499^n$ when this
weight is replaced by $0$, as soon as one phase statistic of the nodes is large enough
(Theorem~\ref{thm:window}). The proof interpolates the data of a singular Pick matrix by a finite Blaschke
product and follows its argument along an arc of the unit circle.

\item \emph{A complete log-gas bound.} With weight $0$ the determinant factors through two Gram
determinants. Fischer's inequality and a finite Andr\'eief identity bound it by two logarithmic gases, one
in the plane and one in the Joukowski coordinate. These are bounded by tangent inequalities against fixed
reference measures. This raises the decay rate by $24.64$. Numerical evaluation for small $N$ suggests that
this is most of the true size of $\det\mathsf L$ at the extremal configuration.

\item \emph{Near-extremal configurations.} The phase statistic need not be large on every configuration.
We show that it is large whenever the configuration comes close to the maximum of OpenAI's energy. On all
other configurations OpenAI's bound already suffices.
\end{enumerate}
We use the parameters $(58,9,3,5,3)$. The general corank theorems of Section~\ref{sec:content} hold for all
parameters, and Theorem~\ref{thm:window} does not involve them.

\subsection{The constants}
Table~\ref{tab:constants} lists the constants for the parameters $(58,9,3,5,3)$, per unit of $N$ or $N^2$.
\begin{table}[ht]\centering
\begin{tabular}{@{}llr@{}}\toprule
degree & $d$ & $53$\\
cost of a prescribed column & $c=c_2+\sigma$ & $3.93$\\
weight of a column & $\psi$ & $50.11$\\
neighbourhood exponent & $\rho=58c+\psi$ & $278.37$\\
decay rate (OpenAI's estimate) & $D_0$ & $6422.83$\\
decay rate (with the bound for $\det\mathsf L$) & $D$ & $6447.47$\\
denominator exponent (with the corank gain) & $\Lambda$ & $6392.78$\\
difference & $\eta=D-\Lambda$ & $54.68$\\
\midrule
bound & $d\rho/\eta$ & $269.78\ldots$\\\bottomrule
\end{tabular}
\caption{The constants of the proof, truncated. The formal proof uses exact rational numbers.}
\label{tab:constants}\end{table}
Without the corank gain and the bound for $\det\mathsf L$, $\eta$ would be about $8$, and the method would give an
exponent above $1800$.

\subsection{Organization}
Section~\ref{sec:overview} recalls OpenAI's construction for general parameters and proves the criterion
that turns the three properties of \S\ref{sec:intro-need} into a measure. Section~\ref{sec:setup} bounds the
denominators and proves nonvanishing at every rational point. Section~\ref{sec:content} proves the corank
bounds, which give the arithmetic gain. Sections~\ref{sec:real}--\ref{sec:signed} treat the analytic side.
Section~\ref{sec:real} extends OpenAI's estimate to a neighbourhood of $G$. Section~\ref{sec:energy} recalls
the energy method with prescribed columns, and Section~\ref{sec:signed} proves the bound for the row
determinant. Section~\ref{sec:proof} assembles the proof, and Section~\ref{sec:remarks} discusses the
method. The certified numerical inequalities are collected in the appendices, and
Section~\ref{sec:formal} describes the formalization.

\subsection{Notation}
We use the notation of OpenAI's paper \cite{OpenAI2026} for the construction and refer to it for details.
Citations of the form \cite[Prop.~2.1]{OpenAI2026} point to its numbered statements. We write $\log$ for
the natural logarithm, $v_p$ for the $p$-adic valuation and $\vartheta(x)=\sum_{p\le x}\log p$ for
Chebyshev's function. $T_k$ and $U_k$ are the Chebyshev polynomials of the first and second kind, with
$U_{-1}=0$, $T_{-k}=T_k$ and $U_{-k-1}=-U_{k-1}$. We write $[X^k]F$ or $F[k]$ for the coefficient of $X^k$
in a polynomial $F$, $y_+=\max(y,0)$, and $\operatorname{rev}_EF(X)=X^EF(1/X)$ for the reversal of a
polynomial of degree at most $E$.

\section{The construction and the criterion}\label{sec:overview}

\subsection{Designs}
OpenAI's determinant depends on five integer parameters. Following Carreira \cite{Carreira2026} we call a
quintuple $(n_0,b_0,q_0,g_0,h_0)$ of integers with $n_0,b_0,q_0,h_0\ge1$, $g_0\ge0$ and $2g_0<n_0$ a
\emph{design}. For $N\ge1$ we write
\begin{gather*}
 n=n_0N,\quad b=b_0N,\quad q=q_0N,\quad g=g_0N,\quad h=h_0N,\\
 L=n+b+q,\quad C=L+g,\quad H=C+h,\quad A=b+q+2g ,
\end{gather*}
and $d_0=n_0-g_0$. OpenAI's choice is $(48,7,4,4,2)$. We work throughout with the design $(58,9,3,5,3)$,
for which
\begin{gather*}
 n=58N,\quad b=9N,\quad q=3N,\quad g=5N,\quad h=3N,\\
 L=70N,\quad C=75N,\quad H=78N,\quad A=22N,
\end{gather*}
and $d_0=53$. The search that led to this design is not part of the proof.

\subsection{The determinant}
For $0\le r<n$ the row polynomials are
\begin{equation}\label{eq:rows}
 P_r(t)=(1-t)^ht^{C-1}T_{|r-g|}(1/t),\qquad \deg P_r<H,
\end{equation}
together with companion polynomials $D_r$ built from $U_{r-g-1}$ \cite[(3)]{OpenAI2026}. Thus $P_r$
depends only on $|r-g|$. Put $f(t)=\sqrt{1-t^2}$ and $dm(t)=|t|\,dt/f(t)$ on $(-1,1)$, and
\[
 A_r(t)=P_r(t)-\frac32\,\mathbf 1_{\{t>0\}}\frac{f(t)}{t}D_r(t),\qquad
 \varphi_k(s)=s^{b+k}(1-s)^q .
\]
The determinant studied by OpenAI is that of the $n\times n$ matrix
\begin{equation}\label{eq:matrix}
 \cA_N(r,k)=\int_{-1}^1\int_0^1 A_r(t)\,\frac{\varphi_k(s)}{1-ts}\,ds\,dm(t)
 \qquad(0\le r,k<n).
\end{equation}
The factor $(1-s)^q$ in $\varphi_k$ amounts to a finite difference of order $q$ applied to the columns. The
entries are rational linear combinations of $1$ and $G$. There are rational matrices $\cA^0_N$ and $\cL_N$
with
\begin{equation}\label{eq:L}
 \cA_N=\cA^0_N+G\,\cL_N,\qquad
 \cL_N(r,k)=\sum_{v=0}^{q}(-1)^v\binom qv\sum_{i}[t^i]P_r\cdot4\,c_{(i-b-k-v)/2},
\end{equation}
where $c_l=4^{-l}\binom{2l}{l}$ for integers $l\ge0$ and $c_z=0$ otherwise
\cite[Prop.~2.1 and (20)]{OpenAI2026}. The rational part is
\begin{equation}\label{eq:rational-part}
 \begin{gathered}
 \cA^0_N(r,k)=\sum_{v=0}^q(-1)^v\binom qv R(r,b+k+v),\\
 R(r,j)=\sum_{i<H}P_r[i]\,M(i,j)-\frac32\sum_{i<H}D_r[i]\,Z(i,j)
 \end{gathered}
\end{equation}
with rational numbers $M(i,j)$, the moments of $m$ against the Cauchy kernel, and $Z(i,j)$, where
$Z(i,i)=-\sum_{k\le i}k^{-2}$ and $Z(i,j)=(H_i-H_j)/(i-j)$ for $i\ne j$, $H_i=\sum_{k\le i}1/k$
\cite[Prop.~2.1]{OpenAI2026}. We call the vectors $R(\cdot,j)$, $b\le j<L$, the \emph{moment columns}. The
columns of $\cA^0_N$ are their differences of order $q$. The design enters \eqref{eq:rational-part} only
through the integer coefficients $P_r[i],D_r[i]$ and the range of $j$. Put
$\Delta_N(X)=\det(\cA^0_N+X\cL_N)\in\Q[X]$, so that $\Delta_N(G)=\det\cA_N$.

\subsection{The degree}
\begin{lemma}\label{lem:degree}
$\deg\Delta_N\le n-g$.
\end{lemma}

\begin{proof}
By \eqref{eq:L} the row $r$ of $\cL_N$ depends only on $P_r$, and hence only on $|r-g|$. So the rows $r$
and $2g-r$ of $\cL_N$ agree for $0\le r<g$, and $2g-r\le n-1$ because $2g<n$. Subtracting row $r$ from row
$2g-r$ for each $r<g$ does not change the determinant and makes $g$ rows of the new coefficient of $X$
vanish. A determinant $\det(B+XL')$ in which $L'$ has $n-g$ nonzero rows has degree at most $n-g$ in $X$.
\end{proof}

\subsection{The criterion}\label{ssec:criterion}
The following criterion turns the three properties of \S\ref{sec:intro-need} into a measure. It is the
reason for every later estimate, so we state and prove it first.

\begin{proposition}\label{prop:criterion}
Let $\xi\in\R$, $d\in\N$, and let $\eta,\rho,\lambda,\delta,\omega>0$ satisfy
\begin{gather*}
 \omega<\lambda,\quad \delta<\tfrac13,\quad d\le(\lambda-\omega)\bigl(\eta+2d(1-3\delta)\bigr),\\
 1<2\bigl(\omega-\delta(2\lambda-\omega)\bigr),\quad \eta<2d,\quad d+4d\delta\lambda<\eta\lambda .
\end{gather*}
For each prime $p$ let $Q_p\in\Z[X]$, and suppose that the following hold for all large $p$.
\begin{enumerate}[label=(\alph*)]
\item $\deg Q_p\le dp$.
\item $Q_p(a/b)\ne0$ whenever $a\in\Z$, $b\ge1$ and $p^2\nmid b$.
\item For all primes $p<p'\le2p$ and all $i\le dp$, $p'^{2i}$ divides the $i$-th coefficient of $Q_p$.
\item $|Q_p(x)|\le e^{-\eta p^2}$ for $|x-\xi|\le e^{-\rho p}$.
\end{enumerate}
Then $|\xi-a/b|>b^{-\rho\lambda}$ for all integers $a$ and all sufficiently large $b$.
\end{proposition}

\begin{proof}
Let $a/b$ be a rational number with $b$ large. Put $u_1=\lambda\log b$ and $u_0=(\lambda-\omega)\log b$,
and suppose $|\xi-a/b|\le b^{-\rho\lambda}$. By the prime number theorem,
$(1-\delta)t\le\vartheta(t)\le(1+\delta)t$ for all large $t$, and $u_0\to\infty$ with $b$. If every prime
$p\in(u_0,u_1]$ had $p^2\mid b$, then $2(\vartheta(u_1)-\vartheta(u_0))\le\log b$. But
\begin{align*}
 2\bigl(\vartheta(u_1)-\vartheta(u_0)\bigr)&\ge2\log b\,\bigl((1-\delta)\lambda-(1+\delta)(\lambda-\omega)\bigr)\\
 &=2\log b\,\bigl(\omega-\delta(2\lambda-\omega)\bigr)>\log b
\end{align*}
by the fourth hypothesis. So there are primes in $(u_0,u_1]$ whose squares do not divide $b$. Let $p$ be
the \emph{largest} one. Then every prime $p'\in(p,u_1]$ has $p'^2\mid b$.

Let $T\le\min(u_1,2p)$ with $T\ge p$, and let $\Pi_T$ be the product of the $p'^2$ over the primes
$p'\in(p,T]$. Then $\Pi_T\mid b$ and $\log\Pi_T=2(\vartheta(T)-\vartheta(p))$. Consider the integer
$b^{dp}Q_p(a/b)=\sum_ic_ia^ib^{dp-i}$. Each of its terms is divisible by $p'^{2dp}$ for each of these $p'$.
Indeed $p'^{2i}$ divides $c_i$ by (c), and $p'^2\mid b$ gives $p'^{2(dp-i)}\mid b^{dp-i}$. So $\Pi_T^{dp}$
divides the integer. Since $p^2\nmid b$, $Q_p(a/b)\ne0$ by (b). So $b^{dp}Q_p(a/b)/\Pi_T^{dp}$ is a nonzero
integer, and
\begin{equation}\label{eq:recovered}
 1\le\Bigl(\frac b{\Pi_T}\Bigr)^{dp}\,|Q_p(a/b)| .
\end{equation}
Since $\rho p\le\rho u_1=\rho\lambda\log b$, we have $|a/b-\xi|\le e^{-\rho p}$. So
$|Q_p(a/b)|\le e^{-\eta p^2}$, and \eqref{eq:recovered} gives
\begin{equation}\label{eq:recovered2}
 \eta p\le d\bigl(\log b-2(\vartheta(T)-\vartheta(p))\bigr).
\end{equation}

Suppose first that $u_1\le2p$, and take $T=u_1$. With
$\vartheta(u_1)-\vartheta(p)\ge u_1-p-2\delta u_1$, \eqref{eq:recovered2} reads
\[
 \eta p\le d\bigl(\log b-2u_1+2p+4\delta u_1\bigr).
\]
Hence $(\eta-2d)p\le d\log b-2du_1+4d\delta u_1$. As $\eta<2d$ and $p\le u_1$, the left side is at least
$(\eta-2d)u_1$. So $\eta u_1\le d\log b+4d\delta u_1$, that is, $\eta\lambda\le d+4d\delta\lambda$. This
contradicts the last hypothesis.

Suppose now that $2p<u_1$, and take $T=2p$. Then $\vartheta(2p)-\vartheta(p)\ge(1-3\delta)p$, so
\eqref{eq:recovered2} gives $p\bigl(\eta+2d(1-3\delta)\bigr)\le d\log b$. Since $p>u_0=(\lambda-\omega)\log b$,
this gives $(\lambda-\omega)\bigl(\eta+2d(1-3\delta)\bigr)<d$, which contradicts the third hypothesis.
\end{proof}

The hypotheses have the following roles. The fourth ensures that the interval $(u_0,u_1]$ contains a prime
whose square does not divide $b$. The last makes the final inequality contradictory when all larger primes
of the interval can be recovered. The proof takes the largest prime of the interval whose square does not
divide $b$. The squares of the larger primes then divide $b$, and (c) recovers those below $2p$. Choosing
any prime of the interval would lose a power of $b$. The chosen prime may be so small that $u_1>2p$, and then
(c) recovers only the primes up to $2p$. The third hypothesis treats this case. A small prime $p$ makes the
decay $e^{-\eta p^2}$ weak, but it also makes the height $b^{dp}$ small, and the primes of $(p,2p]$ alone
already win. So $\lambda$ need not exceed $2\omega$, and the bound $\rho\lambda$ may lie below $\rho$.

For Catalan's constant we take $d=53$, $\omega=\frac35$, $\delta=2.9\cdot10^{-7}$ and $\lambda=270/\rho$.
The last hypothesis then reads $270>\mu_*$ with $\mu_*=d\rho/(\eta-4d\delta)=\Mustar\ldots$ The formal proof
checks the six hypotheses with $\rho\lambda$ in place of $270$ for every $\rho\lambda\in[269.787,400]$.

\section{Denominators and nonvanishing}\label{sec:setup}

This section proves the arithmetic properties of $\Delta_N$: the denominators of its coefficients and its
nonvanishing at rational points. Carreira \cite{Carreira2026} checked that OpenAI's argument goes through
for parameters satisfying a few conditions. We need more. The entries must be controlled modulo $p^2$,
since the corank bounds of Section~\ref{sec:content} work there. And nonvanishing is needed at every
rational point $a/b$ with $p^2\nmid b$, not only at $G$.

\subsection{Corank and valuation}
We use repeatedly the following elementary fact.

\begin{lemma}\label{lem:rank-valuation}
Let $M$ be an $n\times n$ matrix with entries in $\Z_{(p)}$ whose reduction modulo $p$ has rank at most
$n-e$. Then $v_p(\det M)\ge e$.
\end{lemma}

\begin{proof}
Factor the reduction through $\F_p^{n-e}$ and lift the two factors to integral matrices $U,V$, so that
$M=UV+pE$ with $E$ integral at $p$. Expand the determinant multilinearly in the columns. Every term that
uses more than $n-e$ columns of $UV$ vanishes, and every other term has at least $e$ factors $p$.
\end{proof}

\subsection{The digit formula}
The reduction of $p^2\cA^0_N$ modulo $p$ is described by the base-$p$ digits of the coefficients of the
row polynomials. This is the one place where the arithmetic of the moments enters, and it does not depend
on the parameters. For an odd prime $p$ and $c\in\F_p^n$ put
\[
 P_c=\sum_rc_rP_r,\qquad D_c=\sum_rc_rD_r,\qquad W=(1-X^2)^{(p-1)/2},
\]
reading the $P_r,D_r$ as polynomials in $X$.

\begin{lemma}[digit formula]\label{lem:digits}
Let $p$ be an odd prime with $H<p^2$, let $0\le j<H$, and write $j=Jp+\ell$ with $0\le\ell<p$. Then
$p^2R(r,j)$ is $p$-integral for every row $r$, and for every $c\in\Z^n$
\[
 p^2\sum_rc_rR(r,j)\equiv\sum_{u\ge0}\Bigl(\alpha_{u,J}\,(XWP_c)[up+\ell]+\beta_{u,J}\,D_c[up+\ell]\Bigr)
 \pmod p ,
\]
where $\alpha_{u,J}=M(u-1,J)$ and $\beta_{u,J}=-\tfrac32Z(u,J)$ are $p$-integral numbers that depend neither
on $c$ nor on $\ell$.
\end{lemma}

\begin{proof}
Two facts about the numbers $M$ and $Z$ are needed, for $0\le i,j<p^2$. Write $j=Jp+\ell$, let
$\varrho\in[0,p)$ be the residue of $j-i-1$ modulo $p$, and put $u=(i+1+\varrho-\ell)/p$.

First, $p^2M(i,j)$ is $p$-integral and $p^2M(i,j)\equiv W[\varrho]\,M(u-1,J)\pmod p$. This is the digit
reduction of the moments in \cite[\S3]{OpenAI2026}. It is obtained from the $p$-adic analysis of their
recursions and does not involve the parameters.

Second, $p^2Z(i,j)$ is $p$-integral. It is congruent to $Z(\lfloor i/p\rfloor,J)$ if $i\equiv j\pmod p$ and
to $0$ otherwise. To see this, write $i'=\lfloor i/p\rfloor$. Since $i<p^2$, we have $H_i=p^{-1}H_{i'}+S_i$
and $H^{(2)}_i=p^{-2}H^{(2)}_{i'}+S'_i$ with $S_i,S'_i$ $p$-integral, which gives the case $i=j$. If
$i\not\equiv j$, then $p^2Z(i,j)=p\cdot pZ(i,j)\equiv0$. If $i\equiv j$ and $i\ne j$, then $i-j=(i'-J)p$. The
integers $k\in(j,i]$ prime to $p$ form $i'-J$ complete systems of residues modulo $p$, so
$\sum_k1/k\equiv(i'-J)\sum_{r=1}^{p-1}1/r\equiv0\pmod p$. Hence $H_i-H_j\equiv(H_{i'}-H_J)/p$ and
$p^2Z(i,j)\equiv Z(i',J)$.

Now use \eqref{eq:rational-part}. By the first fact, the part $\sum_iF[i]\,p^2M(i,j)$ with $F=P_c$ is
$\sum_iF[i]\,W[\varrho(i)]\,M(u(i)-1,J)$. For fixed $u$ the indices with $u(i)=u$ are those for which
$\varrho=up+\ell-1-i$ lies in $[0,p-1]$. Since $W$ has degree $p-1$, the sum over these $i$ is
$(XFW)[up+\ell]$. This gives the $\alpha$-terms. By the second fact the part with $D_c$ is
$\sum_{i\equiv\ell}D_c[i]\,Z(\lfloor i/p\rfloor,J)=\sum_uZ(u,J)D_c[up+\ell]$, which gives the $\beta$-terms.
The $p$-integrality of $M(u-1,J)$ and $Z(u,J)$ for $u,J<p$ is part of the two facts.
\end{proof}

For a row combination $c\in\F_p^n$ and a moment column $j\in[b,L)$ let $R_c(j)$ be the reduction modulo
$p$ of $p^2\sum_rc_rR(r,j)$. We call it the \emph{value of $c$ at the moment column $j$}. The columns of
$\cA^0_N$ are differences of order $q$ of the moment columns, so $c$ lies in the left kernel of the
reduction as soon as the corresponding differences of the $R_c(j)$ vanish.

We use the digit formula in the following form. Let $0\le k<n$, $s=b+k$, $J=\lfloor s/p\rfloor$ and
$\varsigma=(s\bmod p)+q$. If $\varsigma<p$, the $q+1$ moment columns $s,\dots,s+q$ that enter column $k$ have
the same digit $J$ and residues $(s\bmod p)+v<p$. Since
$\sum_{v=0}^q(-1)^v\binom qvF[t+v]=[X^{t+q}]\bigl((X-1)^qF\bigr)$ for every polynomial $F$,
Lemma~\ref{lem:digits} gives
\begin{equation}\label{eq:column-digits}
 p^2\sum_rc_r\cA^0_N(r,k)\equiv\sum_{u\ge0}\Bigl(\alpha_{u,J}\,Q_c[up+\varsigma]+\beta_{u,J}\,\widehat D_c[up+\varsigma]\Bigr),
\end{equation}
where
\begin{equation}\label{eq:Qc}
 Q_c=(X-1)^q\,XWP_c,\qquad \widehat D_c=(X-1)^qD_c .
\end{equation}

\subsection{The valuation at $p$}
For an odd prime $p$ with $H<p^2$ the entries of $p^2\cA^0_N$ are $p$-integral by Lemma~\ref{lem:digits}.
Let $\cA^0_{N,p}$ be its reduction modulo $p$. Let $\nu_N(p)$ be the corank of $\cA^0_{N,p}$, that is, $n$
minus its rank. The corank enters through the following statement.

\begin{proposition}\label{prop:valuation}
Let $p$ be an odd prime with $H<p^2$ and let $z\in\Q$ have denominator prime to $p$. Then
$p^2(\cA^0_N+z\cL_N)$ has $p$-integral entries, its reduction modulo $p$ equals $\cA^0_{N,p}$, and
\[
 v_p\bigl(\Delta_N(z)\bigr)\ge\nu_N(p)-2n .
\]
\end{proposition}

\begin{proof}
The entries of $p^2\cA^0_N$ are $p$-integral, since each column of $\cA^0_N$ is an integer combination of
moment columns with $j<L\le H$. The entries of $\cL_N$ are $p$-integral for odd $p$, since $P_r[i]\in\Z$
and $c_l\in\Z[1/2]$. So $p^2\cL_N\equiv0\pmod{p^2}$. Hence, for $z$ with denominator prime to $p$,
$p^2(\cA^0_N+z\cL_N)$ has $p$-integral entries with the same reduction as $p^2\cA^0_N$. By
Lemma~\ref{lem:rank-valuation} its determinant has valuation at least $\nu_N(p)$, and
$\det\bigl(p^2(\cA^0_N+z\cL_N)\bigr)=p^{2n}\Delta_N(z)$.
\end{proof}

\subsection{Pole counts}
Let $p$ be an odd prime with $\sqrt H<p\le H$. OpenAI's method \cite[\S3]{OpenAI2026} bounds the valuation
of the entries prime by prime. For any choice of the parameters it gives $v_p(\Delta_N(z))\ge-\ell_N(p)$
for every $z$ with denominator prime to $p$. Here $\ell_N(p)$ is the number of columns whose entries can
have a pole at $p$. It is computed by grouping the coefficients of the row polynomials by their base-$p$
digits. All differences below are truncated at $0$. For $2p\le H$,
\[
 \ell_N(p)=2n-\Bigl(\min\bigl(p,\,L-p,\,A\bigr)-\max\bigl(b,\,H-2p\bigr)-q\Bigr).
\]
For $H<2p\le2H$, with $R=(L-b)-(\min(A,L-p)-b)-(\min(p,L)-\max(b,H-p))$, $P'=\min(A,L-p)-b$ and
$E'=\min(b,H-p)-(L-p)$,
\[
 \ell_N(p)=\min\bigl(2n,\ n+R,\ 2R+P'+E'\bigr).
\]
These are the exponents of the valuation theorems of the formalization, which hold for all parameters. For
the primes that interest us they agree with OpenAI's. For our parameters $\ell_N(p)=Nf(p/N)$ up to integer
rounding, with a piecewise linear function $f$ of $\theta=p/N$:
\[
 f(\theta)=\begin{cases}
 116&\theta\le59/2,\\
 175-2\theta&59/2\le \theta\le69/2,\\
 106&69/2\le \theta\le39,\\
 184-2\theta&39\le \theta\le81/2,\\
 265-4\theta&81/2\le \theta\le48,\\
 217-3\theta&48\le \theta\le69,\\
 148-2\theta&69\le \theta\le70,\\
 78-\theta&70\le \theta\le78,
 \end{cases}
\]
and $f=0$ for $\theta\ge78$, since for $p>H$ the matrix $\cA^0_N+z\cL_N$ is $p$-integral. Its integral over
$[1,78]$ is $24269/4$, and $f=116=2n_0$ on $[0,1]$. Since $f(\theta)=2n_0$ for $\theta\le59/2$, every prime of
the range $(\frac{16}7N,\frac{53}2N]$ of the corank bounds of Section~\ref{sec:content} has
$\ell_N(p)=2n$. There Proposition~\ref{prop:valuation} improves the bound $-2n$ to $\nu_N(p)-2n$.

The primes $p\le\sqrt H$ are not covered by these bounds. For them we use the cruder bound $6n\log_pH$ for
the exponent, which contributes $O(\sqrt N\,N\log N)=o(N^2)$ in total.

\subsection{The prime $2$}
At $p=2$ the entries are not integral, and the bound comes from the expansion of the determinant into
minors, as in \cite[\S3]{OpenAI2026}. Its exponent is governed by a quadratic form in the numbers $m,\ell$
of columns of two exceptional types,
\[
 \tfrac32m^2+\tfrac12\ell^2-(H+C-b)\,m-C\,\ell+Cn .
\]
The following identity shows that it is at least $-\Lambda_2$ whenever $m,\ell\ge0$, $m+\ell\le n$ and
$C\ge n$, where
\begin{equation}\label{eq:A2}
 \Lambda_2=\frac{(n+H-b)^2}8-\frac{n^2}2 .
\end{equation}
With $s=n-m-\ell$ and $m_*=(n+H-b)/4$,
\[
 \tfrac32m^2+\tfrac12\ell^2-(H+C-b)m-C\ell+Cn+\Lambda_2
 =2(m-m_*)^2+s\,(C-n+m)+\tfrac12s^2\ge0 .
\]
This gives a $2$-adic bound $v_2(\Delta_N(i))\ge-\Lambda_2-O(N\log N)$, uniform in the integer $i$. For our
parameters $\Lambda_2=(58+78-9)^2/8-58^2/2=2673/8$ in units of $N^2$.

\subsection{The common denominator}
The corank bounds of Section~\ref{sec:content} give a function $\kappa$ of $\theta=p/N$
(Theorem~\ref{thm:profile-D}). We use it here as an input.

\begin{proposition}\label{prop:denominators}
For the parameters $(58,9,3,5,3)$ there are positive integers $\cD_N$ with the following properties.
$\cD_N\Delta_N(i)\in\Z$ for every integer $i$, and $p'^{2n}\mid\cD_N$ for every prime $N<p'\le2N$. For every
$\varepsilon>0$ and all large $N$
\begin{gather*}
 \log\cD_N\le(\Lambda+\varepsilon)N^2,\\
 \Lambda=\frac{2673}8\log2+116+\frac{24269}4-\int_0^\infty\kappa(\theta)\,d\theta=\LamDen\ldots
\end{gather*}
\end{proposition}

Here $\frac{2673}8\log2$ comes from the prime $2$, and $116+\frac{24269}4$ is the area of $f$. The last term
is the gain from the corank.

\begin{proof}
For primes $p\le H$ with $H<p^2$, let $\nu^*_N(p)$ be the largest of the lower bounds for $\nu_N(p)$ of
Section~\ref{sec:content} that apply at $p$, and $0$ if none applies. Put
$e_N(p)=\min\bigl(\ell_N(p),\,2n-\nu^*_N(p)\bigr)$. For the primes with $p^2\le H$ use the crude bound
above, and for $p=2$ the $2$-adic bound. Let $\cD_N$ be the product of $p^{e_N(p)}$ over $p\le H$. By
Proposition~\ref{prop:valuation} and the bounds of this section, $v_p(\Delta_N(i))\ge-e_N(p)$ for every
integer $i$ and every $p$. For $p>H$ the matrix is $p$-integral. So $\cD_N\Delta_N(i)\in\Z$. A prime
$p'\in(N,2N]$ has $\nu^*_N(p')=0$, since no interval of Theorem~\ref{thm:profile-D} contains it, and
$\ell_N(p')=2n$. So $p'^{2n}\mid\cD_N$. Here $N\ge79$ guarantees $p'>\sqrt H$.

For the size of $\cD_N$ we use the prime number theorem in the weighted form
\begin{equation}\label{eq:weighted-pnt}
 \frac1N\sum_{p}F(p/N)\log p\longrightarrow\int_0^\infty F(\theta)\,d\theta\qquad(N\to\infty),
\end{equation}
for bounded, Riemann integrable, compactly supported $F$. By \eqref{eq:weighted-pnt} the sum of
$\ell_N(p)\log p$ over the odd primes is $N^2\int_0^{78}f+o(N^2)$, with $\int_0^{78}f=116+24269/4$. The primes in
the intervals of Theorem~\ref{thm:profile-D} lie in $(16N/7,53N/2]$ and have $\ell_N(p)=116N$. So their
exponent is at most $116N-N\kappa(p/N)+2$. The sum of $(N\kappa(p/N)-2)\log p$ over these primes is
$N^2\int\kappa+o(N^2)$, again by \eqref{eq:weighted-pnt}. The $2$-adic part contributes
$\Lambda_2\log2\,N^2+o(N^2)$.
\end{proof}

The polynomial $\Delta_N$ has degree at most $n-g=d_0N$ and takes values in $\cD_N^{-1}\Z$ at the integers.
An integer-valued polynomial of degree $m$ has coefficients in $\frac1{m!}\Z$, by Newton's interpolation
formula. So
\begin{equation}\label{eq:QN}
 Q_N=(d_0N)!\,\cD_N\,\Delta_N(X)\in\Z[X],\qquad \deg Q_N\le d_0N,
\end{equation}
and $\log(d_0N)!=o(N^2)$. In this way the gain from the corank enters through the integer polynomial
itself, not through a separate common factor of its coefficients.

The criterion needs the squares of the primes in $(N,2N]$ to divide the coefficients of $Q_N$ in a graded
way.

\begin{lemma}\label{lem:squares}
For all large $N$ and every prime $N<p'\le2N$, the $j$-th coefficient of $Q_N$ is divisible by $p'^{2j}$.
\end{lemma}

\begin{proof}
For these primes the matrices $p'^2\cA^0_N$ and $\cL_N$ are $p'$-integral by Lemma~\ref{lem:digits}, and
$p'^{2n}$ divides $\cD_N$. The identity
\[
 Q_N(X/p'^2)=(d_0N)!\,\frac{\cD_N}{p'^{2n}}\det\bigl(p'^2\cA^0_N+X\cL_N\bigr)
\]
shows that $Q_N(X/p'^2)$ has $p'$-integral coefficients. Its $j$-th coefficient is $c_{N,j}/p'^{2j}$, where
$c_{N,j}$ is the $j$-th coefficient of $Q_N$.
\end{proof}

\subsection{Nonvanishing}
\begin{proposition}\label{prop:nonvanishing}
Suppose that the three $n_0\times n_0$ matrices $B_\sigma$, $\sigma\in\{0,1,-1\}$, defined below, are
invertible. Then for all large primes $p$ and every rational $z$ such that $pz$ has denominator prime to
$p$, $\Delta_p(z)\ne0$. In particular $\Delta_p(a/b)\ne0$ whenever $a\in\Z$, $b\ge1$ and $p^2\nmid b$.
\end{proposition}

The matrices are those of OpenAI's reduction modulo $p$ at $N=p$, with general parameters in place of
OpenAI's. Let $B$ be the rational $(n_0+1)\times n_0$ matrix of the entries of $\cA^0_1$, the rational part at
$N=1$, for the rows $0,\dots,n_0$. The row $n_0$ is an auxiliary row. Let $B_\sigma$ have rows
$B_r+\sigma B_{r+1}$, $0\le r<n_0$. For our parameters $B$ has $59$ rows and $58$ columns.

\begin{proof}
Let $p>4H_0$, with $H_0=H/N$, so that $H=H_0p<p^2/4$. Let $p$ also be prime to the numerators and
denominators of the entries of $B$ and of the three determinants. In the construction at $N=p$ write the
row index $r=pr_0+i$ with $0\le r_0<n_0$ and $0\le i<p$. Since $q=q_0p$, $C=C_0p$, $h=h_0p$ and $g=g_0p$,
the relation $(1-t)^{p}\equiv1-t^p$ expresses the row polynomials modulo $p$ through those at $N=1$ and the
successor row $r_0+1$. This is where the auxiliary row enters. It is valid for rows on both sides of $g$,
because the polynomials are defined through $T_{r-g}$ and $U_{r-g-1}$ for signed indices. The column
functions satisfy $s^{b_0p+k}(1-s)^{q_0p}=s^{\ell}(s^p)^{b_0+k_0}(1-s^p)^{q_0}$ in $\F_p[s]$. So the
columns of one residue class use the moment columns $b_0+k_0,\dots,b_0+k_0+q_0$ at $N=1$ and never change
the residue. By \cite[Lemma~3.1 and eqs.~(31), (35), (36)]{OpenAI2026}, which need only $H<p^2/4$, the
reduction of $p^2\cA^0_p$ modulo $p$ is then, after ordering the rows by $(i,r_0)$ and the columns by
$(\ell,k_0)$, the matrix $a^{\mathsf T}\otimes B_0+b^{\mathsf T}\otimes B_1$. Here $a,b$ are the fixed
$p\times p$ matrices of OpenAI's identity (51), with $\det a=2^{-p(p-1)/2}$. The permutation matrix in this
decomposition has the eigenvalue $0$ once and $\pm1$ each $(p-1)/2$ times. Therefore
\[
 \det=(\det a)^{n_0}\,\det B_0\;\det(B_0+B_1)^{(p-1)/2}\det(B_0-B_1)^{(p-1)/2}\ne0
\]
in $\F_p$. Hence $p^{2n}\Delta_p(0)$ is a $p$-adic unit. Let $z$ be such that $pz$ has denominator prime to
$p$. The entries of $\cL_p$ are $p$-integral, so the difference $p^2z\cL_p=p\,(pz)\,\cL_p$ is $p$-integral
with zero reduction. So the reduction of $p^2(\cA^0_p+z\cL_p)$ is the same invertible matrix, and
$v_p(\Delta_p(z))=-2n$. Finally, if $p^2\nmid b$, then $p(a/b)$ has denominator prime to $p$.
\end{proof}

For the parameters $(58,9,3,5,3)$ we verify that $B_0$, $B_0+B_1$ and $B_0-B_1$ are invertible by exhibiting
their inverses modulo $1009$. This prime divides no denominator of $B$, so the three determinants are
nonzero rational numbers. Proposition~\ref{prop:nonvanishing} then applies to all large primes. The proof
uses only $2g<n$ and $H<p^2/4$.
\section{Corank modulo $p$}\label{sec:content}

This section proves the gain in the denominators. The principle is that of the arithmetic method of Hata
and Rhin--Viola, adapted to a determinant. If the reduction modulo $p$ of $p^2\cA^0_N$ has corank $e$, then
$\Delta_N(z)$ gains a factor $p^e$ at every $z$ with denominator prime to $p$
(Proposition~\ref{prop:valuation}). So we need explicit row combinations $c\in\F_p^n$ with
$c^{\mathsf T}\cA^0_{N,p}=0$.

Here is the heuristic. Modulo $p$ the Chebyshev polynomials of index $kp+d$ factor through $X^p$, so rows
placed symmetrically about the centres $g+kp$ are related. Suitable combinations of such rows have row
polynomials with long runs of vanishing base-$p$ digits. By the digit formula their values at most moment
columns then vanish, and the few remaining values are removed by imposing linear conditions. Each surviving
dimension is a unit of corank.

We first prove the counting inequality (Theorem~\ref{thm:count}). It says that a space of row combinations
subject to a few linear conditions lies in the left kernel. We then derive four closed formulas from it,
which give most of the gain. Finally we describe the further families and the resulting function $\kappa$.
Three facts about characteristic $p$ are used throughout: the Chebyshev identities
\begin{equation}\label{eq:cheb-p}
 \begin{gathered}
 T_p=X^p,\qquad U_{p-1}=(X^2-1)^{(p-1)/2},\\
 T_{kp}=T_k(X^p),\qquad U_{kp-1}=(X^2-1)^{(p-1)/2}\,U_{k-1}(X^p)
 \end{gathered}
\end{equation}
in $\F_p[X]$, the relation $(1-X)^p=1-X^p$, and the digit formula \eqref{eq:column-digits}. The last two
identities of \eqref{eq:cheb-p} follow from $T_{kp}=T_k\circ T_p$ and $U_{kp-1}=(U_{k-1}\circ T_p)\,U_{p-1}$.

\subsection{Setting}
Fix parameters, an integer $N\ge1$ and an odd prime $p$ with $H<p^2$. For $0\le r<n$ put $m_r=r-g$, so that
$-g\le m_r\le n-1-g$ and $n-1-g\ge g$. In the notation of \eqref{eq:rows},
\[
 P_r=(1-X)^h\operatorname{rev}_{C-1}T_{m_r},\qquad D_r=(1-X)^h\operatorname{rev}_{C-1}U_{m_r-1},
\]
where $T_{-m}=T_m$ and $U_{-m-1}=-U_{m-1}$, so that $D_g=0$. For $c\in\F_p^n$ write
$S_c=\sum_rc_r\operatorname{rev}_{C-1}T_{m_r}$ and $U_c=\sum_rc_r\operatorname{rev}_{C-1}U_{m_r-1}$, so that
$P_c=(1-X)^hS_c$ and $D_c=(1-X)^hU_c$. The lowest degree of $\operatorname{rev}_{C-1}T_m$ is at least
$C-1-|m|\ge A=b+q+2g$, and that of $\operatorname{rev}_{C-1}U_{m-1}$ is at least $A+1$. Put $\chi=g+q+h$. Then
the polynomials $Q_c$ and $\widehat D_c$ of \eqref{eq:Qc} have their coefficients in $[A+1,\,L+\chi+p-1]$
and $[A+1,\,L+\chi-1]$ respectively.

A vector $c$ lies in the left kernel of $\cA^0_{N,p}$ as soon as the right side of
\eqref{eq:column-digits} vanishes for every column $k$. We only use columns whose $q+1$ moment columns
share the same base-$p$ digit.

\subsection{The counting inequality}
All the kernels of this section arise in the same way. A space of row combinations is chosen, a few linear
conditions are imposed, and the conditions make every value vanish. The dimension of the space minus the
number of conditions is then a lower bound for the corank. The following theorem states this once. It also
allows $q\ge p$.

\begin{theorem}\label{thm:count}
Let $N\ge1$ and let $p$ be an odd prime with $H<p^2$. Write $q=ap+r$ with $0\le r<p$. Let $\cV$ be a
finite-dimensional vector space over $\F_p$, $\iota\colon\cV\to\F_p^n$ an injective linear map and $\Phi$ a
linear map on $\cV$. Suppose that for every $x\in\ker\Phi$ there is a polynomial $g_x\in\F_p[X]$ of degree
less than $r$ with
\[
 \sum_{v=0}^a(-1)^{a-v}\binom av\,R_{\iota(x)}(j+vp)=g_x(j)\qquad(b\le j<L-ap).
\]
Then $\nu_N(p)\ge\dim\cV-\operatorname{rank}\Phi$.
\end{theorem}

For $q<p$ the hypothesis says that the conditions $\Phi$ force all values at the moment columns to vanish.

\begin{proof}
By \eqref{eq:rational-part}, the entry of $c^{\mathsf T}\cA^0_{N,p}$ at column $k$ is
$\sum_{v=0}^q(-1)^v\binom qvR_c(b+k+v)$. This is the polynomial $(1-T)^q$ applied to $R_c$, where $T$ acts as
the shift. Modulo $p$ we have $(1-T)^p=1-T^p$, hence
\[
 (1-T)^q=(1-T)^r(1-T^p)^a .
\]
This is Lucas's congruence for the binomial coefficients $\binom qv$. So the entry at $k$ equals
$\sum_{u=0}^r(-1)^u\binom ru\,\Sigma_c(b+k+u)$, where
$\Sigma_c(j)=\sum_{v=0}^a(-1)^v\binom av R_c(j+vp)$. The arguments satisfy $b\le j$ and
$j+ap\le b+n-1+q=L-1$. For $c=\iota(x)$ with $x\in\ker\Phi$ we have $\Sigma_c(j)=\pm g_x(j)$. The difference
of order $r$ of a polynomial of degree less than $r$ vanishes identically. So $\iota(\ker\Phi)$ lies in the
left kernel of $\cA^0_{N,p}$. Since $\iota$ is injective, $\nu_N(p)\ge\dim\ker\Phi=\dim\cV-\operatorname{rank}\Phi$.
\end{proof}

Suppose instead that the conditions only force the values on a set $\mathcal C$ of columns to vanish. The
same argument, with the remaining $n-|\mathcal C|$ entries added to the conditions, gives
$\nu_N(p)\ge\dim\cV-\operatorname{rank}\Phi-(n-|\mathcal C|)$. The proofs of the four closed formulas below
are of this form. In Theorems~\ref{thm:T} and~\ref{thm:U} the set $\mathcal C$ consists of all columns, and in
Theorems~\ref{thm:U6} and~\ref{thm:S} it is a block of columns.

\subsection{Four closed formulas}
Each closed formula exhibits explicit vectors in the left kernel. The rows used are symmetric about the
centres $g+kp$, and the identities \eqref{eq:cheb-p} make their combinations vanish modulo $p$. Put
$e=(p+1)/2-q-h$.

\begin{theorem}\label{thm:content}
Let $N\ge1$ and let $p$ be an odd prime with $H<p^2$. Then $\nu_N(p)$ is at least each of the following
quantities whenever the stated conditions hold. Negative values are replaced by $0$.
\begin{enumerate}[label=(\alph*)]
\item (Theorem~\ref{thm:T}) If $2g<p$, $H\le4p$ and $k_T=\min(g,\,n-1-g-2p,\,C-3p)\ge1$:
\[
 k_T-(n+q+g+h-3p)_+ .
\]
\item (Theorem~\ref{thm:U}) If $2h\le p+1$, $L\le3p\le H$, $n+q+g+h\le3p$ and $2p\le n+q$:
\[
 \min(n,\,2p+2g+1)-g-2p-1 .
\]
\item (Theorem~\ref{thm:U6}) If $e\ge0$, $n\le4p+2g+1$, $L\le6p$ and $b+q\le p+\varsigma_0$, with
$\varsigma_0=\max(q,H+q-6p)$:
\[
 \min\Bigl(\frac{p-1}2-g,\ n+b+q+h-\frac{11p+1}2\Bigr).
\]
\item (Theorem~\ref{thm:S}) If $e\ge0$ and $n\le4p+2g+1$: $m_1-m_0-e$, with
\[
 m_0=\max(b+q,\,L+\chi-4p),\qquad m_1=\min(p,\,L-3p,\,A+1).
\]
\end{enumerate}
\end{theorem}

For our parameters these bounds apply for $12N<p\le26N$, where they give most of the gain
(Section~\ref{ssec:kappa}). We prove the four parts in turn.

\subsection{Antisymmetric combinations of rows}
\begin{theorem}\label{thm:T}
Let $2g<p$, $H\le4p$ and $k_T=\min(g,\,n-1-g-2p,\,C-3p)\ge1$. Then
$\nu_N(p)\ge k_T-(n+q+g+h-3p)_+$.
\end{theorem}

The rows used are the differences of rows placed symmetrically about the three centres $g$, $g+p$ and
$g+2p$.

\begin{proof}
For $1\le d\le k_T$ and $k=0,1,2$ let $\epsilon_k(d)=e_{g+kp+d}-e_{g+kp-d}\in\F_p^n$. These rows exist
because $d\le g$ and $g+2p+d\le n-1$. The $3k_T$ vectors $\epsilon_k(d)$ are linearly independent, since their
supports are disjoint. Indeed the supports lie in the intervals $[kp,kp+2g]$ of row indices, which are
disjoint because $2g<p$. The row $g+kp\pm d$ has $m=kp\pm d$. By $T_{a+b}-T_{a-b}=2(X^2-1)U_{a-1}U_{b-1}$
and $U_{a+b-1}-U_{a-b-1}=2T_aU_{b-1}$ together with \eqref{eq:cheb-p},
\begin{align*}
 T_{kp+d}-T_{kp-d}&=2(X^2-1)^{(p+1)/2}\,U_{k-1}(X^p)\,U_{d-1},\\
 U_{kp+d-1}-U_{kp-d-1}&=2\,T_k(X^p)\,U_{d-1}.
\end{align*}
Put $R_d=2(1-X)^hX^{C-d}\operatorname{rev}_{d-1}U_{d-1}$, the value of $D_{\epsilon_0(d)}$. Reverse at degree
$C-1$ and use $\operatorname{rev}_{kp}(T_k(X^p))=(\operatorname{rev}_kT_k)(X^p)$, with $T_0=1$, $T_1=X$ and
$T_2=2X^2-1$. One finds
\[
 D_{\epsilon_1(d)}=X^{-p}R_d,\qquad D_{\epsilon_2(d)}=2X^{-2p}R_d-R_d,
\]
\begin{gather*}
 XWP_{\epsilon_0(d)}=0,\qquad XWP_{\epsilon_1(d)}=X^{-p}R_d-X^pR_d,\\
 XWP_{\epsilon_2(d)}=2X^{-2p}R_d-2R_d .
\end{gather*}
The polynomial $R_d$ is supported in degrees $[C-d,H-1]$, which lies in $[3p,4p)$ because $d\le C-3p$ and
$H\le4p$. So at each residue $\ell$ it has a single base-$p$ digit $\varrho_\ell(d)$, in position $u=3$. For
$x=(x_0,x_1,x_2)\in\F_p^3$ the value of $\sum_kx_k\epsilon_k(d)$ at the moment column $Jp+\ell$ is therefore,
by Lemma~\ref{lem:digits}, $\varrho_\ell(d)\,f_J(x)$ with
\[
 f_J(x)=x_0\beta_{3,J}+x_1(\alpha_{2,J}-\alpha_{4,J}+\beta_{2,J})+x_2(2\alpha_{1,J}-2\alpha_{3,J}+2\beta_{1,J}-\beta_{3,J}).
\]
This depends neither on $d$ nor on $\ell$. The moment columns with $J\ge3$ do not meet the support, since
$3p+\ell\ge C-d\ge C-g=L$. Choose $x^*\ne0$ with $f_1(x^*)=f_2(x^*)=0$, which is two linear equations in three
unknowns. For $c=\sum_ds_d\sum_kx^*_k\epsilon_k(d)$ the values at the columns with $J=1,2$ vanish. The columns
with $J=0$ have $\ell\ge b$, and the value there is $f_0(x^*)\sum_ds_d\varrho_\ell(d)$, the coefficient of
$X^{3p+\ell}$ in $\sum_ds_dR_d$. These are at most $(H-3p-b)_+=(n+q+g+h-3p)_+$ linear conditions on
$s=(s_d)$. The map $s\mapsto c$ is injective, so the solutions form a space of dimension at least
$k_T-(n+q+g+h-3p)_+$ in the left kernel.
\end{proof}

\subsection{A reflection of Chebyshev polynomials}
The next bound rests on a reflection formula modulo $p$.

\begin{lemma}\label{lem:u3-reflect}
Let $p$ be an odd prime and $\varpi=(p+1)/2$. In $\F_p[X]$, for every integer $m$,
\[
 T_{p-m}=X^pT_m-(X^2-1)^{\varpi}\,U_{m-1}.
\]
\end{lemma}

\begin{proof}
Both sides satisfy the three-term recurrence in $m$, so it suffices to check $m=0$ and $m=1$. For $m=0$ it
is $T_p=X^p$, by \eqref{eq:cheb-p}. For $m=1$ use $T_{p+1}=XT_p+(X^2-1)U_{p-1}$ and $T_{p-1}=2XT_p-T_{p+1}$.
\end{proof}

\begin{theorem}\label{thm:U}
Let $2h\le p+1$, $L\le3p\le H$, $n+q+g+h\le3p$ and $2p\le n+q$, and put $n^*=\min(n,2p+2g+1)$. Then
$\nu_N(p)\ge n^*-g-2p-1$.
\end{theorem}

\begin{proof}
We restrict $c$ to the rows $r<n^*$, so that $m_r\le n^*-1-g\le2p+g$. Apply Lemma~\ref{lem:u3-reflect} with
$m=m_r$ and reflect at degree $p+C-1$. For every such row this gives
\[
 (1-X^2)^\varpi\operatorname{rev}_{C-1}U_{m-1}-X\operatorname{rev}_{C-1}T_m=-X^{p+1}\operatorname{rev}_{C-1}T_{p-m}.
\]
The right side is divisible by $X^{p+C-|p-m|}$, hence by $X^L$ since $|p-m|\le p+g=p+C-L$. Multiplying by
$(1-X)^h$ and summing, we get for every $c$ supported on these rows
\begin{equation}\label{eq:identityK}
 X^L\mid(1-X^2)^\varpi D_c-XP_c .
\end{equation}

Let $\cV$ be the space of such $c$ with $(1-X^2)^\varpi\mid P_c$. Since $(1-X)^h$ divides every $P_r$, this is
$\varpi-h$ conditions at $X=1$ and $\varpi$ conditions at $X=-1$, that is, at most $p+1-h$ linear
conditions. For $c\in\cV$ write $P_c=(1-X^2)^\varpi Y$ and $\widetilde Q=XY$. Since $(1-X^2)^p=1-X^{2p}$,
\[
 XWP_c=(1-X^{2p})\,\widetilde Q .
\]
The polynomial $P_c$ has no coefficients below $X^{A^*}$, $A^*=C-n^*+g$, and none from $X^H$ on. So
$\widetilde Q$ is supported in $[A^*+1,\,H-p-1]$. Dividing \eqref{eq:identityK} by the unit $(1-X^2)^\varpi$
modulo $X^L$ gives $D_c[k]=\widetilde Q[k]$ for $k<L$.

Now let $H-3p\le\ell<L-2p$. The inequality $L-2p\le A^*+1$ is equivalent to $n^*\le2p+2g+1$, and
$L-2p\le p$ since $L\le3p$. By the supports, $\widetilde Q[\ell]=0$, $\widetilde Q[2p+\ell]=0$ and
$\widetilde Q[3p+\ell]=0$, and $D_c[3p+\ell]=0$ because $3p+\ell\ge H$. So with
$\varkappa=\widetilde Q[p+\ell]$ the base-$p$ digits of $XWP_c$ at the residue $\ell$ are
$(0,\varkappa,0,-\varkappa,0,\dots)$, and those of $D_c$ are $(0,\varkappa,0,0,\dots)$. Hence, by
Lemma~\ref{lem:digits}, the value at every moment column $Jp+\ell$ with such an $\ell$ is $\varkappa$ times a
constant depending on $J$ only. On $\cV$ these columns, which have $J\in\{0,1,2\}$, therefore span a space
of dimension at most one for each residue $\ell$.

Put $\ell_0=H-3p$ and $\ell_1=L-2p$. The hypotheses give $\ell_0\le b\le\ell_1$. For $\ell_0\le\ell<b$ the
moment columns $j\in[b,L)$ with this residue have $J\in\{1,2\}$, and for $b\le\ell<\ell_1$ they have
$J\in\{0,1,2\}$. So the saving is $2\ell_1-\ell_0-b$. The number of linear conditions needed to make all
values on the $n+q=L-b$ moment columns vanish on $\cV$ is therefore
$(p+1-h)+(n+q)-(2\ell_1-\ell_0-b)$. Substituting $\ell_1=L-2p$ and $\ell_0=H-3p$, this is $2p+1+g$. So the
space of solutions in the span of $n^*$ rows has dimension at least $n^*-g-2p-1$.
\end{proof}

\subsection{A block of columns: the digits $1$ to $5$}
\begin{theorem}\label{thm:U6}
Let $e=(p+1)/2-q-h\ge0$, $n\le4p+2g+1$ and $L\le6p$, and put $\varsigma_0=\max(q,H+q-6p)$ and
$\varsigma_1=L-5p$. Assume $b+q\le p+\varsigma_0$. Then $\nu_N(p)\ge\varsigma_1-\varsigma_0-e$, that is,
\[
 \nu_N(p)\ge\min\Bigl(\frac{p-1}2-g,\ n+b+q+h-\frac{11p+1}2\Bigr).
\]
The first term is the minimum if $6p\le H$ and the second if $6p\ge H$.
\end{theorem}

In Theorems~\ref{thm:T} and~\ref{thm:U} the conditions are imposed on the rows. Here the argument concerns a
block of columns. We first need an identity.

\begin{lemma}\label{lem:K6}
In $\F_p[X]$, for every integer $k$, $2X^p(X^2-1)^{(p-1)/2}T_k-(2X^{2p}-1)U_{k-1}=U_{2p-k-1}$. If
$n\le4p+2g+1$, then for every $c$, $X^L\mid 2XWS_c-(2-X^{2p})U_c$, and consequently
\begin{equation}\label{eq:K6}
 2\,Q_c[i]=2\,\widehat D_c[i]-\widehat D_c[i-2p]\qquad(i<L).
\end{equation}
\end{lemma}

\begin{proof}
For all integers $m,k$ one has $U_{m-1}T_k-T_mU_{k-1}=U_{m-k-1}$, since both sides satisfy the same
recurrence in $k$ and agree for $k=0,1$. Take $m=2p$ and use \eqref{eq:cheb-p}: $T_{2p}=2X^{2p}-1$ and
$U_{2p-1}=2T_pU_{p-1}=2X^p(X^2-1)^{(p-1)/2}$. Reflect at degree $2p+C-1$. The first term becomes
$2XW\operatorname{rev}_{C-1}T_k$, the second $(2-X^{2p})\operatorname{rev}_{C-1}U_{k-1}$, and the right side
$X^{2p+C-1}U_{2p-k-1}(1/X)$. For $k=m_r\ge-g$ with $2p-k-1\ge0$ the right side is divisible by $X^{C+k}$, hence
by $X^L$. If $2p-k-1<0$, it has degree $k-2p-1$ and is divisible by $X^{4p+C-k}$, hence by $X^L$ when
$k\le4p+g$. This holds because $k\le n-1-g\le4p+g$. Summing with weights $c_r$ and multiplying by
$(X-1)^q(1-X)^h$ gives \eqref{eq:K6}.
\end{proof}

\begin{proof}[Proof of Theorem~\ref{thm:U6}]
Let $\mathcal M=[\varsigma_0,\varsigma_1)$ and $w=\varsigma_1-\varsigma_0$. We may assume $w>e$. The block consists of the columns
$k(J,\varsigma)=Jp+\varsigma-(b+q)$ for $1\le J\le5$ and $\varsigma\in \mathcal M$. These indices are nonnegative by the hypothesis
$b+q\le p+\varsigma_0$, and smaller than $n$ because $5p+\varsigma_1-1=L-1$. There are $5w$ of them. As $\varsigma<p$, formula
\eqref{eq:column-digits} applies, and the value of $c$ at $k(J,\varsigma)$ is
$\sum_u\alpha_{u,J}Q_c[up+\varsigma]+\beta_{u,J}\widehat D_c[up+\varsigma]$.

\emph{Supports.} For $\varsigma\in \mathcal M$ we have $p+\varsigma\le p+\varsigma_1-1=L-4p-1\le A$ by $n\le4p+2g+1$. So the digits $u=0,1$
of $Q_c$ and $\widehat D_c$ vanish. From $\varsigma_0\ge H+q-6p=L+\chi-6p$ we get that the digits $u\ge6$ of
$\widehat D_c$ and $u\ge7$ of $Q_c$ vanish. For $u\le5$ we have $up+\varsigma\le L-1$, so \eqref{eq:K6} applies.

\emph{From $Q$ to $\widehat D$.} For $u=2,\dots,5$, \eqref{eq:K6} gives
$\widehat D_c[up+\varsigma]=Q_c[up+\varsigma]+\tfrac12\widehat D_c[(u-2)p+\varsigma]$, and for $u=2,3$ the last index is at most $A$,
so that term vanishes. Hence if $Q_c[up+\varsigma]=0$ for $u=2,\dots,5$ and all $\varsigma\in \mathcal M$, then by induction
$\widehat D_c[up+\varsigma]=0$ for $u=2,\dots,5$, and all digits of $\widehat D_c$ in $\mathcal M$ vanish.

\emph{A relation from the Frobenius map.} We have
$Q_c=(-1)^qX(1-X)^{q+h+(p-1)/2}(1+X)^{(p-1)/2}S_c$, and $q+h+(p-1)/2+e=p$. Since $(1-X)^p=1-X^p$,
\[
 (1-X)^e\,Q_c=(1-X^p)\,Y_c,\qquad Y_c=(-1)^qX(1+X)^{(p-1)/2}S_c,
\]
and $Y_c$ is supported in $[A+1,\,C+(p-1)/2]$. Compare the coefficients at $up+\varsigma$, subtract the
coefficient at $(u-1)p+\varsigma$, and sum over $u=2,\dots,6$. The right side telescopes to
$Y_c[6p+\varsigma]-Y_c[p+\varsigma]$, which vanishes for $\varsigma\ge\varsigma_0+e$. The second term vanishes because $p+\varsigma\le A$,
and the first because $6p+\varsigma\ge H+q+e=C+(p+1)/2>\deg Y_c$. So, if $Q_c[up+\varsigma']=0$ for $u=2,\dots,5$ and all
$\varsigma'\in \mathcal M$, then
\[
 \sum_{t=0}^e(-1)^t\binom et\,Q_c[6p+\varsigma-t]=0\qquad(\varsigma_0+e\le\varsigma<\varsigma_1).
\]
This is a recurrence with unit leading coefficient. So if $Q_c[6p+\varsigma]=0$ for the first $e$ values
$\varsigma=\varsigma_0,\dots,\varsigma_0+e-1$, then it vanishes for all $\varsigma\in \mathcal M$.

\emph{Counting.} Impose on $c$ that the values at the $n-5w$ columns outside the block vanish, that
$Q_c[up+\varsigma]=0$ for $2\le u\le5$ and $\varsigma\in \mathcal M$ ($4w$ conditions), and that $Q_c[6p+\varsigma]=0$ for the first $e$
values of $\varsigma$. By the three steps above all the values at the block columns vanish. These are at most
$n-5w+4w+e=n-(w-e)$ linear conditions, so the left kernel has dimension at least $w-e$. Finally, $w-e$
equals $\varsigma_1-q-e=n+b+q+h-(11p+1)/2$ if $\varsigma_0=q$, and $(p-1)/2-g$ if $\varsigma_0=H+q-6p$.
\end{proof}

\subsection{The digits $1$ to $4$}
\begin{theorem}\label{thm:S}
Let $e=(p+1)/2-q-h\ge0$ and $n\le4p+2g+1$, and put
\[
 m_0=\max(b+q,\,L+\chi-4p),\qquad m_1=\min(p,\,L-3p,\,A+1).
\]
Then $\nu_N(p)\ge m_1-m_0-e$.
\end{theorem}

\begin{proof}
Let $\mathcal M=[m_0,m_1)$ and $w=m_1-m_0>e$. For $\varsigma\in \mathcal M$ we have $\varsigma\le A$, so the digit $u=0$ of $Q_c$ and
$\widehat D_c$ vanishes. The inequality $\varsigma\ge L+\chi-4p$ gives that the digits $u\ge4$ of $\widehat D_c$ and
$u\ge5$ of $Q_c$ vanish. Finally $3p+\varsigma<L$. So \eqref{eq:K6} gives, with $D_i=\widehat D_c[ip+\varsigma]$ and
$Q_i=Q_c[ip+\varsigma]$,
\[
 Q_1=D_1,\qquad Q_2=D_2,\qquad Q_3=D_3-\tfrac12D_1 .
\]
So every value of $c$ at a column with residue $\varsigma\in \mathcal M$ is a combination of the four functionals
$D_1(\varsigma),D_2(\varsigma),D_3(\varsigma),Q_4(\varsigma)$ of $c$. For $\varsigma\in \mathcal M$ the columns $k=\varsigma-(b+q)$, $p+\varsigma-(b+q)$,
$2p+\varsigma-(b+q)$ and $3p+\varsigma-(b+q)$ exist, since $b+q\le\varsigma<p$ and $3p+\varsigma<L$. So there are at least $4w$
block columns, and their entries lie in the span $\mathcal F$ of the $4w$ functionals.

We find $w-e$ linear relations among these functionals. As in the proof of Theorem~\ref{thm:U6},
$(1-X)^eQ_c=(1-X^p)Y_c$ with $Y_c$ a polynomial. Put $\sigma_c(\ell)=\sum_{u\ge0}Q_c[up+\ell]$ and sum the
coefficients of this identity over the residue class of $\varsigma\ge e$ modulo $p$. The right side gives
$\sum_u\bigl(Y_c[up+\varsigma]-Y_c[(u-1)p+\varsigma]\bigr)=0$, since $Y_c$ is a polynomial. The left side gives
$\sum_{t=0}^e(-1)^t\binom et\sigma_c(\varsigma-t)$. For $\varsigma\in \mathcal M$ with $\varsigma-e\ge m_0$ all $\varsigma-t$ lie in $\mathcal M$, and there
$\sigma_c=Q_1+Q_2+Q_3+Q_4=D_1+D_2+D_3-\tfrac12D_1+Q_4$. So
\[
 \sum_{t=0}^e(-1)^t\binom et\bigl(\tfrac12D_1+D_2+D_3+Q_4\bigr)(\varsigma-t)=0\qquad(m_0+e\le\varsigma<m_1).
\]
These $w-e$ relations are independent, since $Q_4(\varsigma)$ occurs with coefficient $1$ in the relation for $\varsigma$
and only with smaller arguments in the others. Hence $\dim \mathcal F\le4w-(w-e)$, and the rank of the matrix is at
most $(n-4w)+3w+e=n-(w-e)$.
\end{proof}

The proof uses neither a determinant nor a lower bound for $N$. Unlike the proof of Theorem~\ref{thm:U6}, the
telescoping here runs over all digits, so no vanishing of $Y_c$ at the ends is needed.

\subsection{Decompositions modulo $p$}
The further families organize the row combinations as follows.

\begin{definition}\label{def:slot}
A \emph{decomposition modulo $p$} of a family of row combinations consists of
\begin{enumerate}[label=(\alph*)]
\item a linear map $x\mapsto c(x)$ from $(\F_p^K)^t$ to $\F_p^n$, injective on a subspace of admissible
 $x=(x_1,\dots,x_t)$;
\item polynomials $\pi_1,\dots,\pi_t\in\F_p[X]$ of degree less than $e_{\max}$;
\item linear maps $\mathcal T,\mathcal G\colon\F_p^K\to\F_p[Y]$,
\end{enumerate}
such that for admissible $x$
\begin{gather*}
 (X-1)^hP_{c(x)}=\sum_d\pi_d(X)\,\mathcal T(x_d)(X^p),\\
 (X-1)^h(X^2-1)^{\frac{p-1}2}D_{c(x)}=\sum_d\pi_d(X)\,\mathcal G(x_d)(X^p).
\end{gather*}
\end{definition}

The identities \eqref{eq:cheb-p} and $(1-X)^p=1-X^p$ provide such decompositions. By the digit formula, the
value of $c(x)$ at the moment column $Jp+\ell$ is a linear form in the $x_d$:
\[
 \sum_{0\le e'<e_{\max},\ p\mid H-1-\ell-e'}\ \Theta_{J,\,(H-1-\ell-e')/p}\Bigl(\sum_d\pi_d[e']\,x_d\Bigr).
\]
Here $\pi_d[e']$ is the coefficient of $X^{e'}$, and each $\Theta_{J,s}$ is a linear form in $J$ and the digit
$s$, determined by $\mathcal T$ and $\mathcal G$.

\subsection{Further families}\label{ssec:families}
The following lower bounds for $\nu_N(p)$ are proved for all parameters, by explicit choices of the space,
the conditions and the columns in Theorem~\ref{thm:count} or its variant. Each is a piecewise linear function
of $p$ and the parameters, positive on intervals described by linear inequalities. We describe them in
words. The statements are those of the formalization.
\begin{enumerate}[label=(F\arabic*),wide]
\item \emph{Frobenius products.} The row combinations have polynomials that are the product of a fixed
factor, a free polynomial of bounded degree $d_P$ in the Chebyshev variable, and a polynomial $\Xi(X^p)$ with
exponents in an interval of $K$ consecutive values. The conditions require $\Xi$ to annihilate the
functionals seen by the columns. If $\#\text{seen}$ of these functionals are independent, the bound is
$(d_P+1)(K-\#\text{seen})$. For our parameters this gives $53N-2p-1$ on $(26N,\frac{53}2N)$, where $2p+2\le53N$.
\item \emph{A truncated form of Theorem~\ref{thm:U6}.} Without the hypothesis $n\le4p+2g+1$, the bound
$\varsigma_1-\varsigma_0-e_+-(n-4p-2g-1)_+$ holds whenever $L\le6p$ and $b+q\le p+\varsigma_0$. For our parameters it gives
$5p-59N+1$ on $(\frac{59}5N,12N]$.
\item \emph{A tail of the antisymmetric family.} With the four centres $g+kp$ ($0\le k\le3$) and
$k_T'=n-1-g-3p$, the bound is $p-1-(2g+q+h)$, under $4p\le L$, $b\le p$, $k_T'\le g$, $k_T'\le(p-1)/2$,
$3p+1+g\le n$ and $4p+b\le H$. For our parameters this is $p-16N-1$ on $(16N,\frac{118}7N]$.
\item \emph{Two families of decompositions and their sum.} The first fixes a digit $s\le H/p$ and a range of
blocks of rows, with polynomials $\pi_d$ divisible by a power of $X-1$. The second has polynomials
$(X^p-1)\pi$ with $\pi$ supported in an interval of degrees. Under a separation condition their kernels are
independent when both apply, at the same digit $s$ or at the digits $s$ and $s-1$. Then the sum of the two
bounds is a lower bound. This is the one place where two bounds are added.
\item \emph{Families defined by rational equations.} In five further families the coefficients of the
decomposition satisfy a fixed system of linear equations with rational coefficients. A solution modulo $p$
exists for every $p$ prime to a fixed integer, and the bound is the dimension of the solution space minus
the number of conditions. In two of these families the entries depend only on the residue of the column
modulo $p$, so the difference of order $q$ imposes $(p-q)_+$ conditions. When $h\ge p$ one writes
$(X-1)^h=(X^p-1)(X-1)^{h-p}$, which imposes $h-p$ more conditions. These families give the bounds on
$(\frac{58}{17}N,\frac{59}{17}N]$, $(\frac{29}8N,\frac{11}3N]$, $(\frac{16}7N,\frac52N]$ and on six intervals in
$(\frac{341}{25}N,\frac{809}{50}N]$.
\item \emph{Block decompositions.} Write a row index as $kp+d$ and put $y=X^p$. For $h\ge p$ the identity
\[
 (X-1)^hWT_{kp+d}=\bigl((X-1)T_k\bigr)(y)\,Q_1+\bigl((X-1)(X^2-1)U_{k-1}\bigr)(y)\,Q_2
\]
and its analogue for $U_{kp+d-1}$ hold in $\F_p[X]$, with $Q_1=(X-1)^{h-p}WT_d$ and $Q_2=(X-1)^{h-p}U_{d-1}$.
At a residue with a single nonzero digit the values of all rows of a block are proportional, and the bound
counts the blocks. For $h<p$ a difference of order $p-h$ of the coefficients plays the same role. For our
parameters these give bounds on $(\frac{59}{21}N,3N]$, where $h\ge p$, and on $(3N,\frac{16}5N]$.
\end{enumerate}

For the intervals of primes that these families do not cover, we apply the counting inequality to every
constant subspace at once. The space of row combinations is the tensor product of a space of polynomials
with a subspace $U$ of a fixed finite-dimensional profile space, and the conditions are Taylor coefficients
composed with fixed linear maps. The corank is then at least $c\dim U$ minus a weighted sum of ranks on $U$, for an integer $c$ (Theorem~\ref{thm:everyU}). We never exhibit $U$. Forced intersections of row spaces show that this quantity is large for some subspace, and only upper bounds for ranks are needed. Appendix~\ref{app:windows} states this form and describes how its
hypotheses are verified.

\subsection{The corank profile for our parameters}\label{ssec:kappa}
The bounds combine into the following statement.

\begin{theorem}\label{thm:profile-D}
For the parameters $(58,9,3,5,3)$ there are $113$ intervals $(a_i,b_i]$, with
$\frac{16}7=a_1<b_1\le a_2<\dots<b_{113}=\frac{53}2$, and a function $\kappa$, linear on each interval and $0$
outside their union, with the following property. For all large $N$ and every prime $p$ with
$a_iN<p\le b_iN$ for some $i$,
\[
 \nu_N(p)\ge N\kappa(p/N)-2 .
\]
On each interval the bound is given by one of the families of this section or by
Theorem~\ref{thm:everyU}. The area of $\kappa$ is a rational number,
\begin{multline*}
 \int_0^\infty\kappa(\theta)\,d\theta=\frac{1002889739024620759352746227576740574893}{45460332302422237200000000000000000000}\\
 =\Chi\ldots\ \ge\ \ChiLean .
\end{multline*}
\end{theorem}

The contributions by ranges of $\theta=p/N$ are as follows. The third column counts the intervals on which
the bound comes from Theorem~\ref{thm:everyU}.
\begin{center}
\begin{tabular}{@{}lrrl@{}}
\toprule
range of $\theta$ & intervals & by Theorem~\ref{thm:everyU} & area \\
\midrule
$(16/7,\,3]$ & $11$ & $0$ & $0.06259$ \\
$(3,\,6]$ & $54$ & $24$ & $0.25354$ \\
$(6,\,12]$ & $31$ & $12$ & $0.71870$ \\
$(12,\,53/2]$ & $17$ & $1$ & $21.02592$ \\
\midrule
total & $113$ & $37$ & $\Chi$ \\
\bottomrule
\end{tabular}
\end{center}
Almost all of the area lies in $(12,\frac{53}2]$, where $\kappa$ is as follows.
\begin{center}
\begin{tabular}{@{}lll@{}}
\toprule
interval of $\theta$ & $\kappa(\theta)$ & bound \\
\midrule
$(12,13]$ & $\theta/2-5$ & Thm.~\ref{thm:content}(c), $(p-1)/2-g$ \\
$(13,\,146/11]$ & $73-11\theta/2$ & Thm.~\ref{thm:content}(c), second term \\
$(341/25,\,69/5]$ & $11\theta/2-75$ & (F5) \\
$(69/5,\,14]$ & $\theta/2-6$ & (F5) \\
$(14,\,711/50]$ & $64-9\theta/2$ & (F5) \\
$(30287/2000,\,78/5]$ & $7\theta/2-53$ & (F5) \\
$(78/5,\,16]$ & $25-3\theta/2$ & (F5) \\
$(16,\,809/50]$ & $73-9\theta/2$ & (F5) \\
$(809/50+\tau,\,1619/100-\tau]$ & $72.999999-9\theta/2$ & Thm.~\ref{thm:everyU} \\
$(1619/100,\,118/7]$ & $\theta-16$ & (F3) \\
$(118/7,\,69/4]$ & $9\theta/2-75$ & Thm.~\ref{thm:content}(d) \\
$(69/4,\,35/2]$ & $\theta/2-6$ & Thm.~\ref{thm:content}(d) \\
$(35/2,\,128/7]$ & $64-7\theta/2$ & Thm.~\ref{thm:content}(d) \\
$(64/3,\,23]$ & $3\theta-64$ & Thm.~\ref{thm:content}(a) \\
$(23,\,24]$ & $5$ & Thm.~\ref{thm:content}(a), then (b) \\
$(24,\,26]$ & $53-2\theta$ & Thm.~\ref{thm:content}(b) \\
$(26,\,53/2]$ & $53-2\theta$ & (F1) \\
\bottomrule
\end{tabular}
\end{center}
In the rows of Theorem~\ref{thm:content} the bound is $3p-(n+q+h)$ for part (a), which applies up to $\theta=70/3$ and is followed by part (b), with the bound
$n-g-2p$ from $\theta=24$. Part (d) gives $(9p-2C)/2$, $(p-1)/2+h-b$ and $(2(n+q+h)-7p)/2$ on its three intervals. Here $\tau=10^{-10}$. The intervals that come from Theorem~\ref{thm:everyU} are shortened by $\tau$ at both ends
and their bounds are lowered by $10^{-6}$, which absorbs the rounding of the integer orders in that theorem.
The areas of the rows of this table add up to $21.02592\ldots$, and the $96$ intervals below $12$
contribute $1.03484\ldots$, so the total is $\Chi\ldots$. The exact sum is computed in the formalization
from the list of intervals. Figure~\ref{fig:corank} shows $\kappa$ on $(12,\frac{53}2]$.

\begin{figure}[ht]
\centering
\setlength{\unitlength}{1cm}
\begin{picture}(12.4,4.6)(-0.8,-0.9)
\put(0,0){\line(1,0){11.8}}
\put(0,0){\line(0,1){3.8}}
\put(0.000,0){\line(0,-1){0.12}}
\put(0.000,-0.4){\makebox(0,0){$12$}}
\put(1.600,0){\line(0,-1){0.12}}
\put(1.600,-0.4){\makebox(0,0){$14$}}
\put(3.200,0){\line(0,-1){0.12}}
\put(3.200,-0.4){\makebox(0,0){$16$}}
\put(4.800,0){\line(0,-1){0.12}}
\put(4.800,-0.4){\makebox(0,0){$18$}}
\put(6.400,0){\line(0,-1){0.12}}
\put(6.400,-0.4){\makebox(0,0){$20$}}
\put(8.000,0){\line(0,-1){0.12}}
\put(8.000,-0.4){\makebox(0,0){$22$}}
\put(9.600,0){\line(0,-1){0.12}}
\put(9.600,-0.4){\makebox(0,0){$24$}}
\put(11.200,0){\line(0,-1){0.12}}
\put(11.200,-0.4){\makebox(0,0){$26$}}
\put(0,0.000){\line(-1,0){0.12}}
\put(-0.35,0.000){\makebox(0,0)[r]{$0$}}
\put(0,0.700){\line(-1,0){0.12}}
\put(-0.35,0.700){\makebox(0,0)[r]{$1$}}
\put(0,1.400){\line(-1,0){0.12}}
\put(-0.35,1.400){\makebox(0,0)[r]{$2$}}
\put(0,2.100){\line(-1,0){0.12}}
\put(-0.35,2.100){\makebox(0,0)[r]{$3$}}
\put(0,2.800){\line(-1,0){0.12}}
\put(-0.35,2.800){\makebox(0,0)[r]{$4$}}
\put(0,3.500){\line(-1,0){0.12}}
\put(-0.35,3.500){\makebox(0,0)[r]{$5$}}
\put(11.700,-0.4){\makebox(0,0)[l]{$\theta$}}
\put(0.15,3.85){\makebox(0,0)[l]{$\kappa(\theta)$}}
\qbezier(0.000,0.700)(0.400,0.875)(0.800,1.050)
\qbezier(0.800,1.050)(0.909,0.525)(1.018,0.000)
\qbezier(1.018,0.000)(1.165,0.000)(1.312,0.000)
\qbezier(1.312,0.014)(1.376,0.322)(1.440,0.630)
\qbezier(1.440,0.630)(1.520,0.665)(1.600,0.700)
\qbezier(1.600,0.700)(1.688,0.353)(1.776,0.007)
\qbezier(1.776,0.000)(2.145,0.000)(2.515,0.000)
\qbezier(2.515,0.002)(2.697,0.561)(2.880,1.120)
\qbezier(2.880,1.120)(3.040,0.910)(3.200,0.700)
\qbezier(3.200,0.700)(3.276,0.401)(3.352,0.102)
\qbezier(3.352,0.133)(3.619,0.367)(3.886,0.600)
\qbezier(3.886,0.600)(4.043,1.219)(4.200,1.837)
\qbezier(4.200,1.837)(4.300,1.881)(4.400,1.925)
\qbezier(4.400,1.925)(4.714,0.962)(5.029,0.000)
\qbezier(5.029,0.000)(6.248,0.000)(7.467,0.000)
\qbezier(7.467,0.000)(8.133,1.750)(8.800,3.500)
\qbezier(8.800,3.500)(9.200,3.500)(9.600,3.500)
\qbezier(9.600,3.500)(10.600,1.750)(11.600,0.000)
\end{picture}
\caption{The corank profile $\kappa(\theta)$ on $(12,\frac{53}2]$, where $\nu_N(p)\ge N\kappa(p/N)-2$ for
$p=\theta N$. The bound on the pole count is $f(\theta)=116$ on this whole range, so $\kappa$ is the gain over
OpenAI's denominators.}\label{fig:corank}
\end{figure}

On each interval we use a single bound. The families come from different constructions, and we have not
shown that their kernels are independent, except for the pair (F4). So the bounds are not added at a prime
where several apply.

For Theorem~\ref{thm:content} one checks the hypotheses on the intervals above. The hypotheses $e\ge0$ and
$n\le4p+2g+1$ both read $\theta\ge12$, and $L\le6p$ reads $\theta\ge35/3$. So Theorems~\ref{thm:U6}
and~\ref{thm:S} apply for $\theta\ge12$. For $12\le\theta\le13$ one has $\varsigma_0=H+q-6p=81N-6p$, so the
first branch of Theorem~\ref{thm:U6} applies. For $\theta\ge13$ one has $\varsigma_0=3N$, and the second
branch applies, which is positive until $\theta=146/11$. In Theorem~\ref{thm:S}, $m_1=p$ for $\theta\le35/2$ and
$m_1=70N-3p$ above, and $m_0=81N-4p$ for $\theta\le69/4$ and $12N$ above. The value $m_1-m_0-e$ is
$9\theta/2-75$, $\theta/2-6$ and $64-7\theta/2$ on the three intervals $(50/3,69/4]$, $(69/4,35/2]$ and
$(35/2,128/7]$, and it is nonpositive outside $(50/3,128/7)$. Theorem~\ref{thm:T} applies for $\theta>10$ and
$\theta\ge39/2$, with $k_T=\min(5N,52N-1-2p,75N-3p)$. It gives $3p-(n+q+h)=3\theta-64$ up to $\theta=23$ and
the bound $g=5N$ up to $\theta=70/3$. From $\theta=70/3$ to $\theta=26$ the hypotheses of Theorem~\ref{thm:U}
hold. They read $3p\ge L$, $3p\le H$, $2p\le n+q$ and $2h\le p+1$. The bound is $\min(n,2p+2g+1)-g-2p-1$,
which is $5N$ for $\theta\le24$ and $53N-2p-1$ above. In each case the value of the theorem is the value in
the table minus at most one unit, from the integer rounding. The same kind of check, with linear
inequalities in $N$, $p$ and the parameters, gives the bounds of Section~\ref{ssec:families} on their
intervals, and Appendix~\ref{app:windows} gives the others. This proves Theorem~\ref{thm:profile-D}. The loss
of $2$ in its statement covers the integer rounding of all the bounds.
\section{The neighbourhood of $G$}\label{sec:real}

This section and the next two prove that $\Delta_N$ is small on a whole neighbourhood of $G$. The idea is to
write the perturbation $\Delta_N(x)-\Delta_N(G)$ as a sum of determinants in which some columns of $\cA_N$
are replaced by values of the row functions at prescribed nodes. Each such determinant is bounded by
OpenAI's energy method, extended to configurations with prescribed points (Section~\ref{sec:energy}). The
loss caused by each replaced column, in the exponent, determines $\rho$. OpenAI's method has two cases. In
the first it already gives the improved decay rate. In the second it falls short when few columns are
replaced, and Section~\ref{sec:signed} closes that gap.

\subsection{The statement}
OpenAI's energy method bounds $|\Delta_N(G)|$ by $\exp\bigl(-(n_0^2(U-\frac12\log2)+o(1))N^2\bigr)$, where $U$
is an upper bound for a discrete logarithmic energy (Section~\ref{sec:energy}). For our parameters we
certify $U=\Utarget$, so OpenAI's estimate gives the decay rate $D_0=n_0^2(U-\frac12\log2)=\DecayOA\ldots$
Let
\begin{gather*}
 \Gamma_0=\Gain,\qquad U_0=U+\Gamma_0/n_0^2=\Uzero\ldots,\\
 D=n_0^2(U_0-\tfrac12\log2)=\Decay\ldots
\end{gather*}
The number $\Gamma_0$ is the improvement of the decay rate proved in Section~\ref{sec:signed}. Put
\begin{equation}\label{eq:rho-def}
 \rho=n_0c+\psi,\qquad c=c_2+\sigma,
\end{equation}
with $c_2=\Ctwo$, $\sigma=\Sigm$ and $\psi=\Psiw$,
so that $\rho=\Rho$. Here $c$ is the increase of the energy bound per prescribed column
(Section~\ref{sec:energy}), and $\psi$ bounds the logarithmic size of the weight and the Vandermonde factor
of one column, per unit of $N$ (\S\ref{ssec:weights}).

\begin{theorem}\label{thm:real}
For our parameters, uniformly for real $x$ with $|x-G|\le e^{-\rho N}$,
\[
 |\Delta_N(x)|\le e^{-(D+o(1))N^2}\qquad(N\to\infty).
\]
\end{theorem}

We use the following abbreviations for general parameters:
\begin{gather*}
 \alpha_A=\frac An,\quad \beta_h=\frac{2h_0}{n_0},\quad \gamma_H=\frac Hn,\quad \gamma_d=\frac{d_0}{n_0},\\
 a_0=1-\frac{b_0+q_0+g_0}{n_0},\quad a_1=3-\frac{b_0+q_0+2g_0}{n_0}.
\end{gather*}
For our parameters these are $\alpha_A=11/29$, $\beta_h=3/29$, $\gamma_H=39/29$, $\gamma_d=53/58$,
$a_0=41/58$ and $a_1=76/29$.

\subsection{A quadrature formula for the coefficient of $G$}\label{ssec:quadrature}
Let $\varepsilon=x-G$. Then $\Delta_N(x)=\det(\cA_N+\varepsilon\cL_N)$. The question is how large the columns
of $\cL_N$ are compared with those of $\cA_N$. By \eqref{eq:L} the entries of $\cL_N$ are linear
functionals of the coefficients of the row polynomials. These polynomials have bounded degree, so the
functionals can be written exactly as finite combinations of point evaluations. Let $t_1,\ldots,t_J$ be
distinct points, with $J$ at least the number of degrees of freedom of the rows. Lagrange interpolation
gives
\begin{equation}\label{eq:quad}
 \cL_N(r,k)=\sum_{a=1}^JP_r(t_a)\,\beta_{a,k},
\end{equation}
where the weight $\beta_{a,k}$ is the value of the same functional on the $a$-th Lagrange basis polynomial
$\ell_a$. This is a quadrature formula of Gauss type. It makes the columns of $\cL_N$ combinations of point
values of the polynomials $P_r$. So the perturbed determinant can be treated by the energy
method. The cost is the size of the weights.

We take $J=d_0N=53N$ nodes, which are Chebyshev nodes of short intervals. Let $[a_c,b_c]$,
$1\le c\le106$, be fixed intervals with dyadic end points, $48$ in $[-1,0)$ and $58$ in $(0,1)$, adjacent
within each side. Their union is $[-0.99931266\ldots,-0.37573294\ldots]\cup[0.19559072\ldots,0.99999983\ldots]$. In
the coordinate $x$ with $t=2x/(1+x^2)$ it is contained in
\[
 E_x=[-\tfrac{2409}{2500},-\tfrac{19501}{100000}]\cup[\tfrac{98749}{1000000},1].
\]
So the nodes avoid a neighbourhood of $x=0$, where the fields of the energy method are largest. We pair the
intervals $(2i-1,2i)$ within each side, which gives $24+29=53$ pairs. The two members of a pair receive
$K_c=\lceil N/2\rceil$ and $\lfloor N/2\rfloor$ nodes, so that $\sum_cK_c=53N$. Let
$\ell_c(z)=(2z-a_c-b_c)/(b_c-a_c)$ and $w_c=b_c-a_c$. The nodes in $[a_c,b_c]$ are the zeros of
$T_{K_c}\circ\ell_c$. They lie in the open interval, so all nodes are distinct.

Let $G(w)=\log|w+\sqrt{w^2-1}|$, with the branch of modulus at least $1$, be the Green function of
$\C\setminus[-1,1]$. Put $G_c=G\circ\ell_c$ and
\begin{equation}\label{eq:Ldef}
 L(z)=\tfrac12\sum_c\bigl(G_c(z)+\log(w_c/4)\bigr).
\end{equation}

\begin{lemma}\label{lem:lagrange}
There is a constant $C$ such that for all $N$, every node $t_a$ and every $z\in\C$ with $|z|\le1$, the
Lagrange polynomial $\ell_a$ of $t_a$ satisfies $|\ell_a(z)|\le CN\,e^{N(L(z)-L(t_a))}$.
\end{lemma}

\begin{proof}
We use three facts about $T_K$. First, $|T_K(w)|\le e^{KG(w)}$, since $2T_K(w)=\phi^K+\phi^{-K}$ with
$\phi=w+\sqrt{w^2-1}$. Second, $|T_K(v)|=\cosh(KG(v))\ge\frac12e^{KG(v)}$ for real $|v|\ge1$. Third, the
Lagrange polynomial $\lambda_j$ of the $j$-th of the $K$ Chebyshev nodes satisfies
$|\lambda_j(w)|\le Ke^{(K-1)G(w)}$ for all $w\in\C$. For the third, write $w=\frac12(u+u^{-1})$ with
$|u|=e^{G(w)}\ge1$ and the node as $\frac12(v+v^{-1})$ with $|v|=1$. Then
\[
 \frac{T_K(w)}{w-\frac12(v+v^{-1})}=\Bigl(\sum_{i<K}u^iv^{K-1-i}\Bigr)\Bigl(\sum_{i<K}(uv)^{-i}\Bigr),
\]
which has modulus at most $K^2|u|^{K-1}$. The product of the differences of the $j$-th node from the others
has modulus at least $K/2^{K-1}$. This gives the third fact.

Let $t_a$ be the $j$-th node of the interval $c$. Grouping the factors of $\ell_a$ by intervals gives
\[
 \ell_a(z)=\lambda_j(\ell_c(z))\prod_{c'\ne c}\frac{T_{K_{c'}}(\ell_{c'}(z))}{T_{K_{c'}}(\ell_{c'}(t_a))}.
\]
For $c'\ne c$ the node lies outside the open interval $(a_{c'},b_{c'})$, so $|\ell_{c'}(t_a)|\ge1$. The
three facts and $G_c(t_a)=0$ give $|\ell_a(z)|\le K_c2^{105}e^{\Lambda(z)-\Lambda(t_a)}$ with
$\Lambda=\sum_cK_cG_c$. Since $|K_c-N/2|\le\frac12$ and $\sum_cG_c$ is bounded on the closed unit disc,
$\Lambda(z)-\Lambda(t_a)\le\frac N2\sum_c\bigl(G_c(z)-G_c(t_a)\bigr)+O(1)$. The constants $\log(w_c/4)$ cancel
in $L(z)-L(t_a)$, and $K_c\le N$.
\end{proof}

The term $-L(t_a)$ depends on the node. We bound it by a piecewise linear function. On each interval
$[a_d,b_d]$ we have $G_d=0$, and for $c\ne d$ the function $|\ell_c|$ is affine with values at least $1$.
Since $\operatorname{arccosh}$ is concave and increasing on $[1,\infty)$, the function $L$ is concave on each
interval. Let $k_0<\dots<k_{865}$ be fixed dyadic knots that include all end points $a_c$, $b_c$, and let
$v_i$ be rational numbers with $v_i\ge-L(k_i)-C_0$, where $C_0=\Czero$. Let $\hat\delta$ be the piecewise
linear function through the points $(k_i,v_i)$. By concavity,
\begin{equation}\label{eq:chord}
 -L(t)\le C_0+\hat\delta(t)
\end{equation}
for every $t$ in the union of the intervals, and so at every node. One value $L(k_i)$ needs no
logarithm per term: $2L(k_i)$ is the logarithm of a product of $106$ factors, each an algebraic number with
one square root.

The weight of column $k$ is a
coefficient functional of $\ell_a$ multiplied by the column function. After removal of the common factor of
the rows it equals
\[
 \beta'_k=4\,[t^0]\,F_k(t)\,(1-t^{-2})^{-1/2},\qquad F_k(t)=t^{\,A-b-q-k}(1-t)^h(t-1)^q\ell_a(t),
\]
where $(1-t^{-2})^{-1/2}=\sum_lc_lt^{-2l}$ with $c_l=\binom{2l}l4^{-l}$. The next lemma bounds it along a contour.

\begin{lemma}[keyhole bound]\label{lem:keyhole}
Let $P$ be a real polynomial, $E\ge0$, and $F(t)=t^{-E}P(t)$. Let $0<r\le1$, and suppose that $|F|\le M$ on
the circle $|t|=r$ and $|F|\le M'$ on the segments $r\le|t|\le1$ of the real axis. Then
\[
 \Bigl|4\sum_{l\ge0}c_l\,[t^{E+2l}]P\Bigr|\le4(1+E)\max(M,M').
\]
\end{lemma}

\begin{proof}
Write $F=\Pi+R$ with $\Pi$ a polynomial and $R=\sum_{i=1}^E\pi_it^{-i}$ the principal part at $0$. The
functional sees only $\Pi$. Since $c_l=\frac1\pi\int_{-1}^1t^{2l}(1-t^2)^{-1/2}dt$, its value is
$\frac4\pi\int_{-1}^1\Pi(t)(1-t^2)^{-1/2}dt$, which is at most $4\max_{[-1,1]}|\Pi|$. By Cauchy's estimate
$|\pi_i|\le Mr^i$, so $|R(t)|\le EM$ for $r\le|t|\le1$, and $|\Pi|\le M'+EM$ there. For $|t|<r$ the maximum
principle gives $|\Pi(t)|\le\max_{|t'|=r}|\Pi|\le(1+E)M$.
\end{proof}

The formal proof replaces the integral by the finite Gauss--Chebyshev identity
$\frac1M\sum_{j<M}\cos^{2l}\frac{(2j+1)\pi}{2M}=c_l$ for $l<M$, and Cauchy's estimate by the finite Cauchy formula
at roots of unity. So no integration is needed. For $0\le y\le1$ and $z\in\C$ let
\[
 \Phi_y(z)=(10-58y)\log|z|+6\log|1-z|+L(z).
\]
Here $10=(A-b-q)/N$ and $6=(q+h)/N$. With $y=k/n$, Lemmas~\ref{lem:lagrange} and~\ref{lem:keyhole} give
\begin{equation}\label{eq:wk}
 |\beta'_k|\le e^{(w(y)-L(t_a)+o(1))N},\qquad w(y)=\max_{z\in\Gamma}\Phi_y(z),
\end{equation}
where $\Gamma$ is the circle $|z|=\frac34$ together with the segments $\frac34\le|z|\le1$ of the real axis. One
radius serves all columns and all nodes. The bound for this maximum is in \S\ref{ssec:weights}.

\subsection{Column expansion and Andr\'eief's identity}\label{ssec:andreief}
Write $\mathcal P(r,a)=P_r(t_a)$. Expanding $\det(\cA_N+\varepsilon\mathcal P\beta)$ multilinearly in the
columns gives a sum over maps $f$ from the columns to $\{\varnothing,1,\ldots,J\}$ of
$\prod_{f(k)\ne\varnothing}\varepsilon\beta_{f(k),k}\det M_f$. Here $M_f$ is the matrix whose $k$-th column is
the $k$-th column of $\cA_N$ if $f(k)=\varnothing$, and the vector $(P_r(t_{f(k)}))_r$ otherwise.

\begin{lemma}\label{lem:expansion}
Suppose that $|\det M_f|\le M_0\prod_{f(k)\ne\varnothing}C_k$ for every $f$, and that $|\beta_{a,k}|\le B_k$.
Then
\[
 \bigl|\det(\cA_N+\varepsilon\mathcal P\beta)-\det\cA_N\bigr|\le M_0\Bigl(\prod_k(1+J|\varepsilon|B_kC_k)-1\Bigr).
\]
\end{lemma}

The lemma allows a different constant for each column. Each $M_f$ has $n'=n-m$ columns that are integrals
and $m$ columns that are point values. Andr\'eief's identity extends to this situation.

\begin{lemma}\label{lem:andreief}
Let $m_1,m_2$ be finite measures, let $F_r$ $(r<n'+m)$, $F'_i$ $(i<n')$ and $k$ be integrable functions, and let
$w_j$ $(j<m)$ be vectors indexed by $r<n'+m$. Let $B$ be the matrix with
$B(r,i)=\iint F_r(t)k(t,s)F'_i(s)\,dm_1(t)\,dm_2(s)$ for $i<n'$ and $B(r,n'+j)=w_j(r)$ for $j<m$. For
$t=(t_0,\ldots,t_{n'-1})$ let $\mathcal W(t)$ be the matrix whose column $c<n'$ is $(F_r(t_c))_r$ and whose column $n'+j$ is
$w_j$. Then
\begin{multline*}
 (n'!)^2\det B=\int\!\!\int\det\mathcal W(t)\\
 \cdot\det[k(t_i,s_j)]_{i,j<n'}\det[F'_i(s_j)]_{j,i<n'}\,dm_1^{n'}(t)\,dm_2^{n'}(s).
\end{multline*}
\end{lemma}

\begin{proof}
For $m=0$ this is Andr\'eief's identity applied twice \cite{Forrester2018Andreief}. In general, expand
$\det B$ multilinearly in its first $n'$ columns. This introduces one pair of variables $(t_i,s_i)$ for each
of them. Then antisymmetrize, first in the $t$ variables and then in the $s$ variables. The last $m$ columns
are constant throughout.
\end{proof}

We apply the lemma to $M_f$ with $F_r=A_r$ and with $w_j=(P_r(t_{a_j}))_r$ at the prescribed nodes $t_{a_j}$.
With $k(t,s)=1/(1-ts)$ the middle determinant is a Cauchy determinant. The third is a generalized
Vandermonde determinant $\prod_js_j^b(1-s_j)^q\det[s_j^{E_i}]$. Here the exponents $E_i<n$ are those of the
columns that remain integrals, and the missing exponents are the replaced columns. We bound the generalized
Vandermonde determinant by a product over the missing exponents.

\begin{lemma}\label{lem:binomvand}
Let $0<\zeta<1$, let $s_1,\ldots,s_{n'}\ge0$, let $E_1<\dots<E_{n'}<n$ be integers, and let $\{k_1,\ldots,k_m\}$ be
the complement of $\{E_i\}$ in $\{0,\ldots,n-1\}$. Then
\[
 \bigl|\det[s_l^{E_i}]\bigr|\le\prod_l(1+\zeta s_l)^m\prod_{j=1}^m\mathcal B_\zeta(k_j)\prod_{i<l}|s_i-s_l|,\qquad
 \mathcal B_\zeta(k)=\sum_{i}\zeta^{-(n'-k+i)},
\]
where the sum runs over the $i<m$ with $0\le n'-k+i\le n'$. Moreover $\mathcal B_\zeta(k)\le\zeta^{-(n-1-k)}/(1-\zeta)$.
\end{lemma}

\begin{proof}
Let $P(X)=\prod_l(X-s_l)=\sum_r(-1)^re_r(s)X^{n'-r}$. Multiplying the matrix by a unitriangular matrix built
from $P$ gives $|\det[s_l^{E_i}]|=\prod_{i<l}|s_i-s_l|\cdot|\det Q|$, where $Q$ is the $m\times m$ matrix with
entries $Q_{j,i}=[X^{k_j-i}]P$. So $|Q_{j,i}|=e_{n'-k_j+i}(s)$, and this is $0$ unless $0\le n'-k_j+i\le n'$.
Since the $s_l$ are nonnegative, $\prod_l(1+\zeta s_l)=\sum_r\zeta^re_r(s)$ is a sum of nonnegative terms. So
$e_r(s)\le\zeta^{-r}\prod_l(1+\zeta s_l)$, and the row $j$ of $Q$ has $\ell^1$ norm at most
$\prod_l(1+\zeta s_l)\,\mathcal B_\zeta(k_j)$. Hadamard's inequality, applied to the rows of $Q$ in the $\ell^1$ norm,
gives the first claim. The exponents $n'-k+i$ in $\mathcal B_\zeta(k)$ are distinct and at most $n-1-k$, which gives the
second.
\end{proof}

\subsection{The weight of one column}\label{ssec:weights}
The bound of Lemma~\ref{lem:binomvand} has two parts. The factor $\prod_l(1+\zeta s_l)^m$ depends on the free points.
With $m=\alpha n$ it equals $\exp\bigl(n\sum_l\alpha\log(1+\zeta s_l)\bigr)$. So it adds the field
$\alpha\log(1+\zeta s)$ to the column field of the energy, and we treat it in Section~\ref{sec:energy}. The other factor
belongs to the replaced columns. With $y=k/n$ we have $\zeta^{-(n-k)}=e^{58N(1-y)\log(1/\zeta)}$. By \eqref{eq:wk}, the
replaced column $k$ at the node $t_a$ therefore contributes at most
\[
 e^{(w(y)+58(1-y)\log(1/\zeta)-L(t_a)+o(1))N}.
\]
We take $\zeta=\frac34$, the radius of the contour $\Gamma$. Two inequalities bound $w$, with $D_\Gamma=\Dcirc$:
\begin{gather}
 6\log|1-z|+L(z)\le D_\Gamma\qquad(|z|=\tfrac34),\label{eq:circlecap}\\
 6\log|1-x|+L(x)-10\log\tfrac34\le D_\Gamma\qquad(x\in\R,\ \tfrac34\le|x|<1).\label{eq:segcap}
\end{gather}
Appendix~\ref{app:field} proves them. On the circle, \eqref{eq:circlecap} gives
$\Phi_y(z)\le(10-58y)\log\frac34+D_\Gamma$. On a real segment put $l=\log\frac34$ and $u=\log|x|$, so that $l\le u\le0$.
Then $(10-58y)(u-l)\le10(u-l)\le-10l$ for $0\le y\le1$, and \eqref{eq:segcap} gives the same bound. The end points
$|x|=1$ follow by continuity. Hence
\[
 w(y)+58(1-y)\log\tfrac43\le(10-58y)\log\tfrac34+D_\Gamma-58(1-y)\log\tfrac34=D_\Gamma-48\log\tfrac34
\]
for every $y$. The dependence on the column cancels exactly. This is the reason for choosing the parameter of
Lemma~\ref{lem:binomvand} equal to the radius of the contour. With \eqref{eq:chord}, the column $k$ at the node $t_a$
contributes at most $e^{(\psi+\hat\delta(t_a)+o(1))N}$ as soon as
\[
 \psi>C_0+D_\Gamma-48\log\tfrac34=50.1113394\ldots
\]
We take $\psi=\Psiw$, which leaves a strict margin for the terms $o(1)$. The factor $e^{N\hat\delta(t_a)}$ depends on the
node. We do not bound it here. It is carried into the energy bound, where it raises the field of each prescribed point
$y$ by $\hat\delta(\mathsf t(y))/58$, with $\mathsf t(y)=2y/(1+y^2)$.

The remaining part of $\rho$ is $n_0c$, the increase of the energy bound caused by the prescribed
columns. It comes from Section~\ref{sec:energy}, and with it $\rho=58\cdot\Cost+\Psiw=\Rho$.
\section{The energy method with prescribed columns}\label{sec:energy}

This section recalls OpenAI's energy method, extends it to determinants with prescribed columns, and
proves Theorem~\ref{thm:real} assuming the bound of Section~\ref{sec:signed}. Throughout, $\alpha_1=\AlphaOne$.

\subsection{OpenAI's energy}\label{ssec:energy}
By Andr\'eief's identity, $\det\cA_N$ is an integral over $n$ points $x_i\in(-1,1)$ and $n$ points
$s_j\in(0,1)$. OpenAI bounds the integrand by $e^{n^2E}$, where $E$ is a discrete logarithmic energy of the
points, written in the coordinate $x$ with $t=2x/(1+x^2)$ \cite[\S5]{OpenAI2026}. It has the form
\begin{multline*}
 n^2E=\sum_iV_1(x_i)+\sum_jV_2(s_j)+\sum_{i<i'}\bigl(\log|x_i-x_{i'}|+\log(1-x_ix_{i'})\bigr)\\
 +\sum_{j<j'}\log|s_j-s_{j'}|-\sum_{i,j}\log(1-2x_is_j+x_i^2),
\end{multline*}
schematically, with external fields $V_1,V_2$ of order $n$. The first two sums are the
external fields. The next two are the logarithmic interactions among points of the same kind, and the last
is the interaction between the two kinds. There are two cases, according to which of OpenAI's two bounds for
the determinant of the row functions is used \cite[Props.~5.2, 5.3]{OpenAI2026}. The case with doubled interaction is the one in which OpenAI's condition
$\sum_i(1-x_i^2)/(1+x_i^2)\le n-1-2g$ \cite[Eq.~(67)]{OpenAI2026} holds, and the constrained case is its complement. In the \emph{constrained
case} the row field acquires a term with a multiplier $\lambda_1$ for the constraint that defines the case.
In the \emph{case with doubled interaction} the interaction among the $x_i$ has weight $2$ instead of $1$. This case also bounds a determinant
$\det\mathsf L$ of size $n$ by $e^{O(n\log n)}$, which is the factor that Section~\ref{sec:signed} improves.

The energy is bounded from above by duality. Given trial sequences $u=(u_k)$ and $v=(v_k)$, one writes the
interaction terms as quadratic forms in the moments $\frac1n\sum_iT_k(x_i)$ and $\frac1n\sum_js_j^k$. Then one
completes the squares against $u$ and $v$. What remains is the sum of the norms $\|u\|^2=\sum_ku_k^2/k$ and
$\|v\|^2$ and, for each point, a function of one variable. With $\mathsf T(u,x)=\sum_ku_kT_k(x)/k$ and
$\mathsf S(v,x)=\sum_kv_kx^k/k$, these functions are, in the case with doubled interaction,
\begin{align*}
 X(x)&=\alpha_A\log|x|+\beta_h\log(1-x)-\gamma_H\log(1+x^2)-4\mathsf T(u,x)-\mathsf S(v,x),\\
 Y(s)&=\tfrac{b_0}{n_0}\log s+\tfrac{q_0}{n_0}\log(1-s)+2\mathsf T(v,s).
\end{align*}
In the constrained case, with a trial sequence $u'$ in place of $u$, they are
\begin{multline*}
 X_1(x)=\alpha_A\log|x|+\beta_h\log(1-x)-(\gamma_H-\tfrac12)\log(1+x^2)\\
 +\lambda_1\bigl(\tfrac{2g_0}{n_0}-\tfrac{2x^2}{1+x^2}\bigr)-2\mathsf T(u',x)-\mathsf S(v,x),
\end{multline*}
with the same $Y$. If $X\le M_X$ and $Y\le M_Y$ on the whole domains, the energy of every configuration is at
most a constant determined by $M_X$, $M_Y$ and the norms. In the language of potential theory the extremal
configuration is described by an equilibrium problem for two interacting measures in an external field
\cite{SaffTotik}. The best trial sequences are the moments of its solution, as Carreira observed
\cite{Carreira2026}. Any trial sequences give a valid upper bound.

\subsection{Prescribed points}\label{ssec:prescribed}
When $m$ of the $n$ columns are point values, only $n-m$ points of each kind are integrated over, and the
energy is still normalized by $n$. The prescribed points enter through the row determinant, now evaluated
at the free points together with the prescribed ones. There are two kinds of prescribed points: the
quadrature nodes, and points $y\in(0,1)$. In the coordinate $x$ the nodes lie in the set $E_W$ of the $x$ with
$|x|\le1$ and $2x/(1+x^2)$ in one of the intervals of Section~\ref{sec:real}. We have $E_W\subset E_x$. Let
$\alpha=m/n$ be the fraction of prescribed columns and $\alpha_2$ the fraction of prescribed points of the
second kind. Each missing free point changes the energy bound by a bounded amount. The total change is linear
in $m$, and its slope $c=c_2+\sigma$ enters $\rho$.

Two terms from Section~\ref{sec:real} are added. The factor $\prod_l(1+\zeta s_l)^m$ of Lemma~\ref{lem:binomvand}
adds the field $\alpha g(s)$, with $g(s)=\log(1+\frac34s)$, to the column field. Since $g>0$ on $(0,1]$, the
value with $\alpha g$ is at most the value with $\gamma g$ for every $\gamma\ge\alpha$. The factor
$e^{N\hat\delta}$ of the weights adds $\hat\delta(\mathsf t(x))/58$ to the field of every prescribed point $x$,
where $\mathsf t(x)=2x/(1+x^2)$. Indeed $n^2\cdot\frac1n\sum_j\hat\delta(\mathsf t(y_j))/58=N\sum_j\hat\delta(\mathsf t(y_j))$.

Every step of OpenAI's argument remains valid uniformly in the number of free points. The regularization of
the logarithmic interactions, the diagonal corrections and the completion of squares only involve sums over
the points present. For the prescribed points of the first kind the relevant field is
\[
 F_1(x)=-\gamma_d\log(1+x^2)-2\mathsf T(u,x)+\tfrac1{58}\hat\delta(\mathsf t(x)),\qquad x\in E_W .
\]
For those of the second kind it is
\[
 F_2(y)=-\gamma_d\log(1+y^2)+\frac{n_0-2g_0-1}{n_0}\log y-u_1y+\tfrac1{58}\hat\delta(\mathsf t(y))
\]
for $y\in E_W$ with $0<y<1$, where $u_1$ is the first coefficient of $u$. The coefficient $(n_0-2g_0-1)/n_0$ is a lower bound, valid for all
$N\ge1$, of $(n-1-2g)/n$. The exact coefficient of $-\log(1+x^2)$ in both fields is $(n-g-1)/n=\gamma_d-1/n$. Since
$0\le\log(1+x^2)\le\log2$ for $|x|\le1$, the exact fields exceed $F_1$ and $F_2$ by at most $\log2/n$, which tends to zero. By Lemma~\ref{lem:binomvand} the generalized Vandermonde determinant contributes
the factor $\prod_{i<l}|s_i-s_l|$ and the column field $\alpha g$ to the energy, and the factors $\mathcal B_\zeta$
are assigned to the columns. The energy bound then reduces to a quadratic inequality in $\alpha$ and $\alpha_2$.
The maximum over $\alpha_2$ has a closed form.

\begin{lemma}[energy inequality with prescribed points]\label{lem:energy-ineq}
Let $\ell=\log2$ and let $\mathcal N,\mathcal F,\mathcal R,\mathcal F_2$ be real numbers. Put
\begin{multline*}
 V(\alpha)=\bigl(-(a_0-\tfrac12)+(a_1-1)\alpha-\alpha^2\bigr)\ell+\mathcal N+(1-\alpha)\mathcal F+\alpha\mathcal R\\
 +\frac{(\mathcal F_2+\ell-\mathcal R)^2}{4\ell}+U_0-\tfrac12\ell-c_2\alpha .
\end{multline*}
Then for all real $\alpha$ and all $\alpha_2\ge0$
\begin{multline*}
 \bigl(-(a_0-\tfrac12)+(a_1-1)\alpha-\alpha^2+\alpha_2-\alpha_2^2\bigr)\ell+\mathcal N+(1-\alpha)\mathcal F\\
 +(\alpha-\alpha_2)\mathcal R+\alpha_2\mathcal F_2\le-U_0+\tfrac12\ell+c_2\alpha+V(\alpha) .
\end{multline*}
\end{lemma}

\begin{proof}
The difference of the two sides is $\alpha_2(\mathcal F_2+\ell-\mathcal R)-\ell\alpha_2^2$, and
$tb-\ell t^2\le b^2/(4\ell)$ for all real $t$ and $b$.
\end{proof}

So the energy bound holds at the slope $c=c_2+\sigma$ on every range of $\alpha$ on which $V(\alpha)\le\sigma\alpha$.
The lemma is applied with $\mathcal N$ a bound for $2\|u\|^2+\frac12\|v\|^2$, $\mathcal F$ a bound for $M_X+M_Y$, with the
column field $\gamma g$ included in $Y$, $\mathcal R$ a bound for the supremum of $F_1$ and $\mathcal F_2$ a bound for the
supremum of $F_2$. Keeping $\mathcal R$ and $\mathcal F_2$ apart is much better than one cap for both, because the moved
trial sequences of \S\ref{ssec:trials} lower the first field and raise the second. The analogue for the constrained case has no prescribed
points of the second kind. Its inequality is
\[
 \bigl(-(a_0-\tfrac12)+(a_1-1)\alpha-\alpha^2\bigr)\ell+n_1+(1-\alpha)\mathcal F_1+\alpha\mathcal M_1
 \le-U_0+\tfrac12\ell+c\alpha,
\]
with $n_1$ a bound for the norms of the trial sequences, $\mathcal F_1=M_{X_1}+M_Y$, and $\mathcal M_1$ a bound for
$\lambda_1\bigl(\frac{2g_0}{n_0}-\frac{2x^2}{1+x^2}\bigr)+\hat\delta(\mathsf t(x))/58$ on $E_W$. Since
$1-\frac{2x^2}{1+x^2}=\sqrt{1-\mathsf t(x)^2}$, this bound is a maximum of affine functions of $\lambda_1$ over the
segments of the chord table, with $\hat\delta$ replaced by its larger end value. So it is convex in $\lambda_1$.

\emph{The squares that remain.} Section~\ref{sec:signed} uses what the completion of squares leaves, so we record it for the case with doubled
interaction. Let $\mu_F$, $\mu_M$ and $\mu_S$ be the normalized counting measures of the free points, the prescribed points of the first kind and the
prescribed points of the second kind. Their total masses are $1-\alpha$, $\alpha-\alpha_2$ and $\alpha_2$. Put $W=\mu_F+\mu_M$,
$A_k(\mu)=\int T_k\,d\mu$ and $\langle a,b\rangle=\sum_ka_kb_k/k$, so that $\|a\|^2=\langle a,a\rangle$. Three products carry the interactions
in the variable $x$. The Vandermonde product of the free points comes from the Cauchy determinant. The row determinant gives the Vandermonde product of the
free points and the prescribed points of the first kind, and the Vandermonde product of the points $y$ of the second kind. These last points
enter the row determinant through $1/y$, so it also gives the cross factors $|w-1/y|=(1-wy)/y$, for $w$ a point of $W$.
Since $-\log|x-x'|=\log2+2\sum_kT_k(x)T_k(x')/k$ and $\log(1-wy)\le-wy$, the part of $E$ quadratic in the Chebyshev moments is, apart from the
diagonal terms and the multiples of $\log2$, at most
\[
 -\|A(\mu_F)\|^2-\|A(W)\|^2-\|A(\mu_S)\|^2-A_1(W)A_1(\mu_S).
\]
For any measure $\mu$ we have $-\|A(\mu)\|^2=\|u\|^2-2\langle u,A(\mu)\rangle-\|A(\mu)-u\|^2$. Using this for $\mu_F$ and $W$, and writing
$A_1(W)=u_1+(A_1(W)-u_1)$, we find that the display equals
\[
 2\|u\|^2-2\langle u,A(2\mu_F+\mu_M)\rangle-u_1A_1(\mu_S)-\mathcal Q,
\]
where
\[
 \mathcal Q=\|A(\mu_F)-u\|^2+\|A(W)-u\|^2+\|A(\mu_S)\|^2+(A_1(W)-u_1)A_1(\mu_S).
\]
The first three terms are the term $2\|u\|^2$ of $\mathcal N$ and the terms of $X$, $F_1$ and $F_2$ that contain $u$. So the energy bound obtained from
Lemma~\ref{lem:energy-ineq} improves by $\mathcal Q$. To bound $\mathcal Q$ from below, let $e_1$ be the first
unit vector and put $a=A(\mu_F)-u$ and $b=A(W)-u+\frac12A_1(\mu_S)e_1$. Since $\|e_1\|=1$,
\[
 \|b\|^2=\|A(W)-u\|^2+(A_1(W)-u_1)A_1(\mu_S)+\tfrac14A_1(\mu_S)^2 .
\]
Also $\frac14A_1(\mu_S)^2\le\|A(\mu_S)\|^2$, so $\mathcal Q\ge\|a\|^2+\|b\|^2\ge2\|\frac12(a+b)\|^2$. The average of $a$ and $b$ is
$A(\mu_F+\frac12\mu_M)-u+\frac14A_1(\mu_S)e_1$. This gives
\begin{equation}\label{eq:squares}
 \mathcal Q\ge\Delta,\qquad\Delta=2\sum_k\frac1k\bigl(A_k(\mu_F+\tfrac12\mu_M)+\tfrac14A_1(\mu_S)\delta_{k1}-u_k\bigr)^2 .
\end{equation}
The weight $\frac12$ of $\mu_M$ records that these points occur in one of the two Vandermonde products of the $x$ variables, and the weight
$\frac14$ is half of the $\frac12$ in $b$. The diagonal terms are treated by OpenAI's regularization, which replaces the moments by damped moments
with an error that tends to zero, as in the proof of Lemma~\ref{lem:freetangent}. We suppress this.

\subsection{The trial sequences}\label{ssec:trials}
The trial sequences have the form
\begin{equation}\label{eq:trial}
 u_k=c_k+\operatorname{Re}\sum_jr_jz_j^k\qquad(k\ge1;\ c_k=0\text{ for }k>\Kcoef),
\end{equation}
with $\Kcoef$ Chebyshev coefficients and finitely many geometric terms. The geometric terms have a simple
meaning. Since $\sum_{k\ge1}z^kT_k(x)/k=-\frac12\log(1-2xz+z^2)$, a term $(z,r)$ contributes
$-\frac12\operatorname{Re}\,r\log(1-2xz+z^2)$ to the potential $\mathsf T(u,x)$. If
$x_0+i\epsilon_0=\frac12(z+z^{-1})$, then $1-2xz+z^2=2z(x_0+i\epsilon_0-x)$. So the term is, up to a constant,
the logarithmic potential of a point charge of complex strength $r$ at $x_0+i\epsilon_0$ and its conjugate. The
extremal measure has a gap in its support around $x=0$. Its potential has a logarithmic singularity at $x=0$
and square-root behaviour at the edges of the gap. A Chebyshev expansion converges slowly near such points,
while a charge at a small distance $\epsilon_0$ reproduces the local behaviour with one term. This is also why
the quadrature nodes of Section~\ref{sec:real} avoid a neighbourhood of $x=0$.

\emph{Moving trial sequences.} The best trial sequence depends on $\alpha$. We use families
\[
 u(\tau)=u+\tau\,\mathrm du,\qquad v(\tau)=v+\tau\,\mathrm dv\qquad(0\le\tau\le1),
\]
with fixed finitely supported increments and with $\tau=\tau(\alpha)$ an explicit piecewise linear function. The fields
$X$, $Y$, $F_1$ and $F_2$ are affine in $\tau$. So the supremum of a field at $\tau$ is at most the chord between the
suprema at $\tau=0$ and $\tau=1$, or, for the kept base fields, the supremum at $\tau=0$ plus $\tau$ times the supremum of the
increment. The norms are quadratic in $\tau$. All bounds then become polynomials in $\alpha$ on each piece of
$\tau(\alpha)$. Appendix~\ref{app:field} gives the charges, the increments and the certified bounds.

\subsection{The two cases at the improved rate}\label{ssec:cases}
Theorem~\ref{thm:real} asks for the rate $U_0=U+\Gamma_0/n_0^2$ instead of $U$, at the slope $c$. In both cases this
reduces to finitely many polynomial inequalities in $\alpha$, each of the form $V(\alpha)\le\sigma\alpha$ on a closed
interval with polynomial caps. Each is proved by a cover of the interval by subintervals on which the leading
term dominates, in exact rational arithmetic. The bound at an interior $\alpha$ of a cell uses the caps at the right
end point of the cell, since $g>0$. A finite cover by cells and the uniform continuity of the caps then give the
bound for all $\alpha$, with an error that tends to zero.

In the constrained case we use the increment with $24$ Chebyshev coefficients and $\tau=\min(1,\frac{50}{11}\alpha)$. The
inequality holds on $[0,1]$, which is covered by the $6$ pieces with end points $0,\frac{11}{50},\frac3{10},\frac25,\frac35,\frac45,1$.
In the case with doubled interaction we use the increment with $40$ Chebyshev coefficients and $\tau=1$. The inequality
of Lemma~\ref{lem:energy-ineq} holds on $[\alpha_1,1]$, covered by the $6$ pieces with end points
$\alpha_1,\frac15,\frac3{10},\frac25,\frac35,\frac45,1$. For $\alpha<\alpha_1$ it fails, because the extra term
$U_0-U$ is absorbed by the term $c\alpha$ only once $\alpha$ is not small. So a new argument is needed only in the case
with doubled interaction with fewer than $\alpha_1n$ prescribed columns. Section~\ref{sec:signed} supplies it, with the
same slope $c$. Precisely, it proves the following.

\begin{proposition}\label{prop:signed-goal}
In the case with doubled interaction, for every configuration with $m\le\alpha_1n$ prescribed columns at the nodes
$t_{a_1},\ldots,t_{a_m}$, the integrand of Lemma~\ref{lem:andreief}, with the generalized Vandermonde determinant replaced by
the bound of Lemma~\ref{lem:binomvand}, is at most
\[
 \exp\Bigl(n^2\bigl(-U_0+\tfrac12\log2+c\alpha+o(1)\bigr)-N\sum_j\hat\delta(t_{a_j})\Bigr)\,C_*^m\prod_j\mathcal B_\zeta(k_j),\qquad \alpha=m/n,
\]
with an absolute constant $C_*$.
\end{proposition}

\begin{proof}[Proof of Theorem~\ref{thm:real}, given Proposition~\ref{prop:signed-goal}]
By Proposition~\ref{prop:signed-goal} and the discussion above, in both cases a determinant $M_f$ with $m$
prescribed columns $k_1,\ldots,k_m$ at the nodes $t_{a_j}$ is at most
\[
 C_*^m\prod_j\mathcal B_\zeta(k_j)\,e^{n_0cNm-N\sum_j\hat\delta(t_{a_j})}\exp\bigl(n_0^2N^2(-U_0+\tfrac12\log2+o(1))\bigr).
\]
The factor $C_*^m$ comes from the expansion of the row functions at the prescribed points. By \eqref{eq:wk} and
\S\ref{ssec:weights}, the weight of the column $k_j$ at the node $t_{a_j}$, times $\mathcal B_\zeta(k_j)$, is at most
$e^{(\psi+\hat\delta(t_{a_j})+o(1))N}$. The terms $\hat\delta$ cancel. So for every column $k$ and every node $a$
the weight and the factor of $M_f$ that belong to the pair $(k,a)$ have the product
\[
 \log(B_kC_k)\le\bigl(n_0c+\psi+o(1)\bigr)N=(\rho+o(1))N ,
\]
and Lemma~\ref{lem:expansion} applies with this bound for each pair.
For $|\varepsilon|\le e^{-\rho N}$ Lemma~\ref{lem:expansion} gives a factor $(1+O(N^5))^n=e^{o(N^2)}$. With
$D=n_0^2(U_0-\frac12\log2)$ this is the claim.
\end{proof}
\section{A bound for the row determinant}\label{sec:signed}

This section proves Proposition~\ref{prop:signed-goal}, which completes the proof of Theorem~\ref{thm:real}.
In the case with doubled interaction OpenAI's integrand contains a determinant $\det\mathsf L$ of size $n$
that the energy method bounds by $e^{O(n\log n)}$. We replace this determinant by a simpler one, with a loss of only
$e^{O(N)}$, and we bound the simpler one by a Gram determinant. The argument has three steps.

The first step is a comparison. Let $F(w)$ be the Pick matrix of the nodes with weight $w$ on the positive
nodes. Then $\det\mathsf L$ is $\det F(-5)$ up to a known factor, and a positivity argument bounds it by
$499^n\det F(0)$ as soon as $F(-\frac1{100})\ge0$ (\S\ref{ssec:rowdet}). The second step shows that
$F(-\frac1{100})\ge0$ whenever a phase inequality holds for the nodes (\S\ref{ssec:window},
\S\ref{ssec:windowproof}), and that this inequality holds on every configuration for which the energy bound
is not already sufficient (\S\ref{ssec:alternative}). The third step bounds $\det F(0)$, which is a Vandermonde
determinant of a Laurent system. It splits into a Vandermonde determinant of the positive nodes and a
determinant of folded functions at the negative nodes. Fischer's inequality and Andr\'eief's identity turn
the latter into two log-gas integrals (\S\ref{ssec:gram}). Reference measures then bound these integrals,
and a numerical inequality closes the proof (\S\ref{ssec:budget}).

\subsection{The row determinant and the comparison}\label{ssec:rowdet}
In the case with doubled interaction OpenAI's integrand contains the determinant of the row functions at
points $z_1,\ldots,z_n\in(-1,1)\setminus\{0\}$, which we call the nodes \cite[Prop.~5.2]{OpenAI2026}. They are
the free points $x_i$ and the prescribed points $y_j$ together. We write $\det\varphi(z)$ for this determinant. With $\mathsf e=n-1-2g$ it factors as
\begin{equation}\label{eq:detphi}
 \det\varphi(z)=C_{n,g}\,V(z)\prod_a|z_a|^{-(n-1-g)}\det\mathsf L(z),
\end{equation}
where $C_{n,g}$ depends only on $n$, $g$ and the number of nodes of each kind, and $V(z)=\prod_{a<b}(z_b-z_a)$. To describe $\mathsf L$ let
\[
 B(z)=\prod_a\frac{z-z_a}{1-z_az},\qquad \cK_{ab}=\frac1{1-z_az_b},\qquad D_a=B'(z_a)z_a^{\mathsf e},
\]
so that $B$ is the Blaschke product of the nodes and $\cK$ is the Pick matrix of the Szeg\H{o} kernels. Let
$E_-$ and $E_+$ be the diagonal blocks of $\operatorname{diag}(D_a)$ at the negative and the positive nodes.
For $w\in\R$ let
\[
 F(w)=\cK+\operatorname{diag}(E_-,wE_+).
\]
The computation of \cite[Prop.~5.2]{OpenAI2026} gives, for every real vector $(w_a)$, the identity
\begin{equation}\label{eq:laurent}
 \det\bigl(z_a^{g-r}+w_az_a^{r-g}\bigr)_{a,\,0\le r<n}=c(z)\det\bigl(\cK+\operatorname{diag}(w_aD_a)\bigr),
\end{equation}
where $c(z)$ does not depend on $w$. If all nodes are free, the rows of OpenAI's determinant are
$\frac12(z^{g-r}+z^{r-g})$ at a negative node and $-\frac14z^{g-r}+\frac54z^{r-g}$ at a positive node. So
$\det\mathsf L=\det F(-5)/\det\cK$. The value $-5=1-6$ comes from the factor $\frac32$ in the row functions.
OpenAI's argument uses only $|\det\mathsf L|\le e^{O(n\log n)}$, which is negligible at the scale $e^{N^2}$. Numerical
evaluation suggests that $\log|\det\mathsf L|$ is of order $-N^2$ near the extremal configuration
(Remark~\ref{rem:L}).

A prescribed point has the row $\frac12(z^{g-r}+z^{r-g})$ at either sign. At a positive prescribed point this row
is $-2$ times the row of a free positive node plus $3$ times the row $(z^{r-g})_r$. Let $\Xi$ be the matrix of
all rows at the nodes. Expanding its determinant in these rows gives
\begin{equation}\label{eq:expandS}
 |\det\Xi|\le\sum_{S}2^{m_+-|S|}3^{|S|}|\det\Xi_S|,
\end{equation}
where $m_+$ is the number of positive prescribed points and $S$ runs over the subsets of these points. In
$\Xi_S$ a node of $S$ has the row $(z^{r-g})_r$, a negative node has the row $\frac12(z^{g-r}+z^{r-g})$,
and every other positive node has the row of a free positive node. The total weight in \eqref{eq:expandS} is at
most $5^m$. A node of $S$ is what Section~\ref{sec:energy} calls a prescribed point of the second kind. Let
$\Xi^0_S$ be the same matrix with the rows of the positive nodes outside $S$ replaced by $(z^{g-r})_r$.

Both $\Xi_S$ and $\Xi^0_S$ are obtained from \eqref{eq:laurent}. The left side of \eqref{eq:laurent} is
affine in each $w_a$. Taking the coefficient of $\prod_{a\in S}w_a$ gives
\begin{equation}\label{eq:minors}
 \frac{|\det\Xi_S|}{|\det\Xi^0_S|}=4^{-|P\setminus S|}\frac{|\det F(-5)_{S^c}|}{\det F(0)_{S^c}},
\end{equation}
where $P$ is the set of positive nodes and the index $S^c$ denotes the principal submatrix on the complement of
$S$. The case $w=1$ is a Pick matrix.

\begin{lemma}\label{lem:pick}
Let $H=\operatorname{diag}(B'(z_a))$ and let $X$ be a real diagonal matrix. If $\cK-X\cK X\ge0$, then $\cK+HX\ge0$. In
particular $F(1)\ge0$ for every configuration.
\end{lemma}

\begin{proof}
First we show that $\cK H^{-1}\cK=H$. For a real vector $c$ put $\phi_c(z)=B(z)\sum_ac_a/(z-z_a)$. It has node values
$\phi_c(z_a)=c_aB'(z_a)$. Partial fractions give $\phi_c=\sum_ad_ak_a$ with $k_a(z)=1/(1-z_az)$, for a real vector $d$, and comparing node values
gives $\cK d=Hc$. On the unit circle $|\phi_c(z)|=|\sum_ac_a/(z-z_a)|=|\sum_ac_ak_a(\bar z)|$. So $\|\phi_c\|_{H^2}^2=c^{\mathsf T}\cK c$, and also
$\|\phi_c\|_{H^2}^2=d^{\mathsf T}\cK d=c^{\mathsf T}H\cK^{-1}Hc$. Hence $H\cK^{-1}H=\cK$, which is the claim.

Let $T=H^{-1}\cK X$ and $A=I+T$. Then $\cK T=HX$, $T^{\mathsf T}\cK=XH$ and $T^{\mathsf T}\cK T=X\cK X$. So
\[
 0\le A^{\mathsf T}\cK A=\cK+2HX+X\cK X .
\]
Adding $\cK-X\cK X\ge0$ gives $2(\cK+HX)\ge0$. For the second statement take $X=\operatorname{diag}(z_a^{\mathsf e})$. The entries of
$\cK-X\cK X$ are $\sum_{j<\mathsf e}(z_az_b)^j$, which is a Gram matrix.
\end{proof}

\begin{lemma}\label{lem:compare}
If $F(-\frac1{100})\ge0$, then $|\det F(-5)_T|\le499^{|T|}\det F(0)_T$ for every set $T$ of nodes. Hence
$|\det\Xi_S|\le499^n|\det\Xi^0_S|$ for every $S$.
\end{lemma}

\begin{proof}
Put $Y=F(0)$ and $Q=\operatorname{diag}(0,E_+)$. By hypothesis and Lemma~\ref{lem:pick}, $Y+Q\ge0$ and
$Y-\frac1{100}Q\ge0$. A convex combination of the two gives $Y\ge0$. Let $Z=F(-5)=Y-5Q$. Then
\[
 499Y-Z=5(Y+Q)+493Y\ge0,\qquad 499Y+Z=500\bigl(Y-\tfrac1{100}Q\bigr)\ge0 .
\]
So $-499Y\le Z\le499Y$, and the same holds for principal submatrices. If $Y_T>0$, all eigenvalues of
$Y_T^{-1/2}Z_TY_T^{-1/2}$ have modulus at most $499$, which is the claim. The singular case follows by replacing
$Y_T$ by $Y_T+tI$ and letting $t\to0$. The second statement follows from \eqref{eq:minors}.
\end{proof}

So everything reduces to two questions: when is $F(-\frac1{100})\ge0$, and how large is $\det\Xi^0_S$?

\subsection{A phase criterion for positivity}\label{ssec:window}
The matrix $F(-c)$ is a Pick matrix with a small negative weight on the positive nodes. The following theorem
gives a criterion for its positivity that depends only on the positions of the positive nodes. It holds for
all configurations and has no relation to our parameters.

\begin{theorem}\label{thm:window}
Let $z_1,\ldots,z_n\in(-1,1)\setminus\{0\}$ be distinct, with $k_+$ positive and $k_-$ negative entries. Let
$t\ge0$ be an integer, let $c=\frac1{100}$, and let $D_a=B'(z_a)z_a^t$. For a positive node $p$ put
$\rho_p=(1-p)/(1+p)$. If
\begin{equation}\label{eq:window}
 t\arctan4-\sum_{p>0}\arctan(4\rho_p)>\frac{\pi}{2}(k_-+1)+\frac{17}4\,\frac{c^2}{1-c^2}\,k_+,
\end{equation}
then $\cK+\operatorname{diag}(E_-,-cE_+)\ge0$.
\end{theorem}

The number $\arctan4$ is the half-angle of the arc of the unit circle around $z=-1$ with end point
$z_0=(4i-1)/(4i+1)$. The proof compares two phases on this arc. We apply the theorem with $t=\mathsf e$.

\subsection{Proof of Theorem~\ref{thm:window}}\label{ssec:windowproof}
Let $f_a(\gamma)=\gamma z_a^t$ at a negative node and $f_a(\gamma)=-cz_a^t$ at a positive node. Let
$W(\gamma)_{ab}=(1-f_a(\gamma)f_b(\gamma))\cK_{ab}$ be the Pick matrix of these data. By Lemma~\ref{lem:pick} it suffices to prove
$W(1)\ge0$. With $X=\operatorname{diag}(z_a^t)$ we have
\[
 W(-c)=(1-c^2)\cK+c^2(\cK-X\cK X),
\]
and $\cK-X\cK X\ge0$ as in Lemma~\ref{lem:pick}. So $W(-c)>0$. If $W(1)$ is not positive semidefinite, then there is
$\gamma\in(-c,1)$ with $W(\gamma)\ge0$ singular. Fix such a $\gamma$ and write $W=W(\gamma)$, $f_a=f_a(\gamma)$. Note that $|f_a|<1$.

Let $\mathcal T$ be a minimal set of nodes such that the principal submatrix $W_{\mathcal T}$ is singular. Since the diagonal entries of $W$ are
positive, $|\mathcal T|\ge2$. By minimality $W_{\mathcal T}$ has a kernel vector $v$ with all entries nonzero, and every proper principal
submatrix is positive definite. From now on only the nodes of $\mathcal T$ are used, and $k_+$ and $k_-$ count these nodes.
Inequality \eqref{eq:window} still holds, because removing a node lowers the right side or raises the left side. If all nodes of $\mathcal T$ have
the same sign, the data are a constant of modulus less than $1$ times $z^t$, and $W_{\mathcal T}>0$ as above. So $k_+\ge1$ and $k_-\ge1$.

We construct a Blaschke product that interpolates the data. Put
\[
 \varphi(\zeta)=\sum_{a\in\mathcal T}\frac{v_a}{1-z_a\zeta},\qquad \psi(\zeta)=\sum_{a\in\mathcal T}\frac{v_af_a}{1-z_a\zeta}.
\]
The relation $Wv=0$ says that $\varphi(z_b)=f_b\psi(z_b)$ for $b\in\mathcal T$. For $|\zeta|<1$ let $\beta_a=(z_a-\zeta)/(1-z_a\bar\zeta)$. From
$(1-z_a\bar\zeta)(1-z_b\zeta)-(z_a-\zeta)(z_b-\bar\zeta)=(1-z_az_b)(1-|\zeta|^2)$ we get
\[
 \frac{1-|\zeta|^2}{(1-z_a\bar\zeta)(1-z_b\zeta)}=\frac{1-\beta_a\bar\beta_b}{1-z_az_b}.
\]
Multiplying by $v_av_b(1-f_af_b)$ and summing gives
\begin{equation}\label{eq:phipsi}
 (1-|\zeta|^2)\bigl(|\psi(\zeta)|^2-|\varphi(\zeta)|^2\bigr)=u^*W_{\mathcal T}u,\qquad u_a=v_a\bar\beta_a,
\end{equation}
because $v^{\mathsf T}W_{\mathcal T}v=0$. So $|\varphi|\le|\psi|$ in the disc. For $|\zeta|=1$ one has
\[
 \frac1{(1-z_a\bar\zeta)(1-z_b\zeta)}=\frac1{1-z_az_b}\Bigl(\frac1{1-z_a\bar\zeta}+\frac1{1-z_b\zeta}-1\Bigr),
\]
and summation with $W_{\mathcal T}v=0$ gives $|\varphi|=|\psi|$ on the circle. Since the functions $1/(1-z_a\zeta)$ are independent, $\varphi\ne0$ and so $\psi\ne0$.
Hence $h=\varphi/\psi$ is a rational function with real coefficients, with $|h|\le1$ in the disc and $|h|=1$ on the circle. So it is a finite
Blaschke product. Both $\varphi$ and $\psi$ are polynomials of degree at most $|\mathcal T|-1$ divided by $\prod_a(1-z_a\zeta)$, so the degree of $h$ is at most
$|\mathcal T|-1$. Thus
\[
 h=\varepsilon z^s\frac{P}{P^*},\qquad \varepsilon=\pm1,\quad P^*(z)=z^{\deg P}P(1/z),
\]
where $P$ is a real polynomial with all roots in the open disc and $P(0)\ne0$. Put $d=s+\deg P$. Then $d+1\le k_++k_-$.

We have $h(z_b)=f_b$ for $b\in\mathcal T$. This is clear if $\psi(z_b)\ne0$. If $\psi(z_b)=0$, then $\varphi(z_b)=0$, and \eqref{eq:phipsi} at $\zeta=z_b$ gives $u^*W_{\mathcal T}u=0$
with $u_b=0$ and $u\ne0$. This contradicts the positive definiteness of $W_{\mathcal T\setminus\{b\}}$. We show that $h$ cannot exist.

\emph{The case $t\le s$.} The real polynomial $\varepsilon z^{s-t}P+cP^*$ is nonzero, because $|P|=|P^*|$ on the
unit circle and $c<1$. It vanishes at the $k_+$ positive nodes and has degree at most $d-t$. So $k_++t\le d\le
k_++k_--1$, that is, $t<k_-$. This contradicts \eqref{eq:window}, which gives $t\arctan4>\frac\pi2k_-$.

\emph{The case $s<t$.} Put $T=t-s\ge1$ and $r=h/z^t=\varepsilon P/(z^TP^*)$. Then $r=\gamma$ at the $k_-$ negative
nodes and $r=-c$ at the $k_+$ positive ones. On the unit circle $zh'/h=t+\Omega$, with $\Omega=zr'/r$ real.

\begin{lemma}\label{lem:contact}
$\Omega\ge0$ on the unit circle.
\end{lemma}

\begin{proof}
Let $v=\deg P$. The numerator of $zr'/r$ is the self-reciprocal real polynomial $zP'P^*-zP^{*\prime}P-TPP^*$ of
degree at most $2v$. So it equals $z^vR(z+1/z)$ with $\deg R\le v\le k_++k_--1$, and $\Omega=R(2\cos\theta)/|P|^2$
at $z=e^{i\theta}$. Rolle's theorem gives a critical point of $r$ between consecutive nodes of the same sign,
hence $k_++k_--2$ distinct roots of $R$ outside $[-2,2]$. So $R$ has at most one further root. Suppose that $R<0$
somewhere in $[-2,2]$. If $R(2)\ge0$ and $R(-2)\ge0$, then $R$ has two further roots in $[-2,2]$. If $R(2)<0$, then
$|r|$ decreases at $z=1$ on the real axis, so $|r|>1$ just to the left of $1$. Since $|r|<1$ at the largest node,
$r$ takes the value $r(1)=\pm1$ on the right of that node. Rolle's theorem gives one further critical point,
and a further root of $R$. The same holds at $-2$. If exactly one of $R(\pm2)$ is negative, this root comes with a sign
change of $R$ in $[-2,2]$. If both are negative, there are two such roots. In every case $R$ has two further roots,
which contradicts the degree bound.
\end{proof}

Let $p_1,\ldots,p_m$ be the positive nodes, so $m=k_+$ in the rest of this proof, and let $k=k_-$. We compare the
phase of $h$ with the phase of an explicit interpolant of degree $m$. Put $\lambda=c^2$ and $w_a=p_a^t$. Put
$\mathsf P_{ab}=(1-\lambda w_aw_b)/(1-p_ap_b)$ and $M=\mathsf P^{-1}$, and write $\langle a,b\rangle=a^{\mathsf T}Mb$. Let
$\xi_z(a)=1/(1-zp_a)$, $g_1=1$ and $g_2=-cw$. With $F_i(z)=g_i\xi_z$, $C_j=g_j\xi_{-1}$ and $J=\operatorname{diag}(1,-1)$ define
\[
 \Theta_{ij}(z)=\delta_{ij}-(1+z)J_j\langle F_i(z),C_j\rangle .
\]
So $\Theta(-1)=I$. Direct expansion gives the kernel identity
\begin{equation}\label{eq:kernelid}
 \sum_k\Theta_{ik}(z)J_k\Theta_{jk}(u)=\delta_{ij}J_i-(1-zu)\langle F_i(z),F_j(u)\rangle .
\end{equation}
The vector $F_1(p_a)-f_aF_2(p_a)$ is row $a$ of $\mathsf P$, where now $f_a=-cw_a$. This gives the node relation
$\Theta_{1j}(p_a)=f_a\Theta_{2j}(p_a)$. Let $D(z)=\prod_a(1-p_az)$, $b(z)=\prod_a(z-p_a)$ and
$\kappa=(-1)^m$. The functions $e_{ij}=D\Theta_{ij}$ are real polynomials of degree at most $m$. Their
determinant vanishes at the $p_a$ by the node relation and at the $1/p_a$ because $(e_{ij}(1/p_a))$ has rank one.
So it equals $\kappa bD$. Taking $u=1/z$ in \eqref{eq:kernelid} gives
\[
 e_{11}(z)=\kappa z^me_{22}(1/z),\qquad e_{12}(z)=\kappa z^me_{21}(1/z).
\]
On the unit circle \eqref{eq:kernelid} with $i=j=2$ and $u=\bar z$ gives $|\Theta_{22}|^2=1+|\Theta_{21}|^2$. Let
$\sigma=h(-1)=\pm1$ and $e=\sigma e_{21}+e_{22}$. Then $W=e/D=\Theta_{22}(1+\sigma\Theta_{21}/\Theta_{22})$ does not
vanish on the circle, and $W(-1)=1$.

\emph{The size of the argument of $\Theta_{22}$.} Let $T_l$ be $\Theta_{22}$ with $\lambda$ replaced by $l$, so
$T_0=1$. Let $\cK^{(p)}$ be the matrix with entries $1/(1-p_ap_b)$, let $Z=\operatorname{diag}(w_a)$, and let
$M_l=(\cK^{(p)}-lZ\cK^{(p)}Z)^{-1}$. For $0\le l<l'<1$ and real vectors $a,b$ the resolvent identity gives
\[
 l'\,a^{\mathsf T}ZM_{l'}Zb-l\,a^{\mathsf T}ZM_lZb=(l'-l)(M_{l'}Za)^{\mathsf T}\cK^{(p)}(M_lZb) .
\]
Since $\cK^{(p)}-Z\cK^{(p)}Z\ge0$, the right side has modulus at most $(l'-l)\|a\|\|b\|/((1-l)(1-l'))$, where
$\|a\|^2=a^{\mathsf T}(\cK^{(p)})^{-1}a$. For $c=0$ the identity \eqref{eq:kernelid} at $z=\zeta$ and $u=\bar\zeta$ gives
$(1-|\zeta|^2)\|\xi_\zeta\|^2=1-|b(\zeta)/D(\zeta)|^2$. Since $1-\prod t_a\le\sum(1-t_a)$ for $t_a\in[0,1]$, this yields
\[
 \|\xi_\zeta\|^2\le\sum_a\frac{1-p_a^2}{|1-p_a\zeta|^2}\qquad(|\zeta|\le1).
\]
In particular $\|\xi_{-1}\|^2\le m$. For $z$ on the arc from $z_0$ to $-1$ one has $\operatorname{Re}z\le\frac{15}{17}$,
and a quadratic inequality in $p$ gives $|1+z|^2(1-p^2)/|1-pz|^2\le(\frac{17}4)^2$. So
$|1+z|\,\|\xi_z\|\,\|\xi_{-1}\|\le\frac{17}4m$, and
$|T_{l'}(z)-T_l(z)|\le\frac{l'-l}{(1-l)(1-l')}\cdot\frac{17}4m$ on the arc. Since $|T_l|\ge1$, a telescoping product
over a fine grid of $[0,\lambda]$ shows that the continuous argument of $\Theta_{22}$ on the arc, normalized to vanish
at $-1$, has modulus at most $\mathcal X=\frac{17}4m\lambda/(1-\lambda)$. The factor $1+\sigma\Theta_{21}/\Theta_{22}$ has
positive real part. So the continuous argument $\ell(\theta)$ of $W(e^{i\theta})$, with $\ell(\pi)=0$, satisfies
$|\ell(\theta_0)|\le\frac\pi2+\mathcal X$, where $z_0=e^{i\theta_0}$ and $\pi-\theta_0=2\arctan4$.

\emph{The phases.} By Lemma~\ref{lem:contact}, $h(e^{i\theta})=\sigma e^{-i\phi(\theta)}$ with
$\phi(\theta)=\int_\theta^\pi(t+\Omega)\ge t(\pi-\theta)$. Let
$\psi(\theta)=m(\pi-\theta)+2\sum_a\arg(1-p_ae^{i\theta})$. Then $(-z)^m\bar e/e=\exp(-i(\psi+2\ell))$. At $\theta_0$,
$1-pz_0=(1+p)(1+4i\rho_p)/(1+4i)$. So $\psi(\theta_0)=2\sum_p\arctan(4\rho_p)$.

\emph{The root count.} Let $\hat e(z)=z^me(1/z)$ and
\[
 \mathcal N=\varepsilon z^sPe-\sigma\kappa P^*\hat e,
\]
a real polynomial of degree at most $d+m$. Put $a=\phi-\psi-2\ell$, so $a(\pi)=0$. At every $\theta$ with
$a(\theta)\in2\pi\Z$ the polynomial $\mathcal N$ vanishes at $e^{\pm i\theta}$. It also vanishes at every $p_a$ by the
node relation and the reversal formulas, at every $1/p_a$ because $z^{d+m}\mathcal N(1/z)=\pm\mathcal N(z)$, and at $-1$. Suppose
that $a(\theta_0)>\pi k$ and let $L=\lfloor k/2\rfloor$. If $\mathcal N\ne0$, the intermediate value theorem gives $L$
conjugate pairs of roots on the open arc. So $2L+2m+1\le d+m\le2m+k-1$, which contradicts $2L+2>k$. If
$\mathcal N=0$, then $e^{ia}=1$ on the arc, which contradicts the intermediate value theorem at the value $\pi$.
Hence $a(\theta_0)\le\pi k$, and
\[
 2t\arctan4\le\phi(\theta_0)\le\pi k+2\sum_p\arctan(4\rho_p)+\pi+2\mathcal X .
\]
This contradicts \eqref{eq:window}, and Theorem~\ref{thm:window} is proved.\qed

\subsection{The alternative on near-extremal configurations}\label{ssec:alternative}
We now show that \eqref{eq:window} holds with $t=\mathsf e$ and $c=\frac1{100}$ on every configuration for which the
energy bound of Section~\ref{sec:energy} does not already suffice. The idea is that such a configuration has
Chebyshev moments close to the trial sequence $u$, and the left side of \eqref{eq:window} is then computed by a
fixed test function.

Let $\phi_H(w)=\frac\pi2-\arctan\bigl(4(1-w)/(1+w)\bigr)$ for $0\le w<1$. Then $\arctan(4\rho_p)=\frac\pi2-\phi_H(p)$.
With $\mathsf e=48N-1$, $n=58N$ and $k_-=n-k_+$, the inequality \eqref{eq:window} follows for large $N$ from
\begin{equation}\label{eq:phaseneed}
 \frac1n\sum_{p>0}\phi_H(p)\ge\arctan\tfrac14+\tfrac5{29}\arctan4+\tfrac{17}4\,\frac{c^2}{1-c^2}+\xi,\qquad \xi=\tfrac1{1000},
\end{equation}
whose right side is $0.474992\ldots$. Let $f=\frac{F_0}2+\sum_{k\ge1}F_kT_k$ be a polynomial of degree $32$ with
dyadic coefficients and denominator $2^{31}$, where $F_0=992223195/2^{30}$. 
It satisfies $f\le\phi_H$ on $(0,1)$, $f\le0$ on $[-1,0]$, and
\begin{equation}\label{eq:testS}
 \tfrac14F_0+\tfrac14|F_1|y\le\phi_H(y)\qquad(0<y<1).
\end{equation}
Also $\sum_kkF_k^2\le7877/10000$. On this range of $\alpha$ the case with doubled interaction uses the moving trial
sequence $u(\tau)=u+\tau\,\mathrm du$ of \S\ref{ssec:trials}, with the increment of $40$ Chebyshev coefficients and
$\tau=\alpha/\alpha_1$. We have $\sum_kF_ku_k\ge879/10000$ and $\sum_kF_k\mathrm du_k\ge-0.010526$.

Fix a configuration and a set $S$. Let $\mu_F$, $\mu_M$ and $\mu_S$ be the measures of \S\ref{ssec:prescribed}, with the points of $S$ as the
prescribed points of the second kind, and let $\beta$ be the mass of $\mu_S$. By \eqref{eq:squares}, the bound of Lemma~\ref{lem:energy-ineq} with
$\alpha_2=\beta$ improves by the number $\Delta$ of \eqref{eq:squares}. Divide $[0,\alpha_1]$ into $64$ equal cells. On the cell $j$ the
left side of the inequality of Lemma~\ref{lem:energy-ineq}, minus the right side at the slope $c$, is at most a rational number $\Delta_j$.
These numbers lie between $10^{-6}$ and $0.008365$. So either $\Delta\le\Delta_j$, or the energy bound gives the conclusion of
Proposition~\ref{prop:signed-goal} for $\Xi_S$ directly.

Suppose that $\Delta\le\Delta_j$. Since $\phi_H\ge0$ on $(0,1)$ and $f\le\phi_H$, while $f\le0$ on the negative
half, we have $\frac12f\le\phi_H$ on positive prescribed points and $\frac12f\le0$ on negative ones. With
\eqref{eq:testS},
\begin{align*}
 \frac1n\sum_{p>0}\phi_H(p)&\ge\int f\,d(\mu_F+\tfrac12\mu_M)+\tfrac14F_0\beta+\tfrac14|F_1|A_1(\mu_S)\\
 &\ge\tfrac12F_0-\tfrac14F_0\alpha+\sum_kF_ku_k(\tau)-\sqrt{\tfrac12\Delta_j\sum_kkF_k^2},
\end{align*}
by the Cauchy--Schwarz inequality. The right side is affine in $\alpha$ on each cell. At the end points of every cell it exceeds
$0.4752$, which is larger than the right side of \eqref{eq:phaseneed}. The smallest margin is $0.0050$. So \eqref{eq:phaseneed} holds, and
Theorem~\ref{thm:window} gives $F(-\frac1{100})\ge0$. All the numerical facts of this subsection are part of the formal proof, which also carries the regularization of the moments that we have suppressed here. We have proved the following.

\begin{proposition}\label{prop:alternative}
For $N$ large, every configuration with $m\le\alpha_1n$ prescribed columns in the case with doubled interaction, and
every $S$, satisfies at least one of two statements. Either the energy bound for $\Xi_S$ gives the conclusion
of Proposition~\ref{prop:signed-goal}, or $F(-\frac1{100})\ge0$ and therefore
$|\det\Xi_S|\le499^n|\det\Xi^0_S|$.
\end{proposition}

\subsection{Gram determinants}\label{ssec:gram}
It remains to bound $\det\Xi^0_S$. Let $\mathsf D=n-1-g=53N-1$. Let $\xi_1,\ldots,\xi_{k_+}$ be the positive nodes, with
$\xi=z$ outside $S$ and $\xi=1/z$ on $S$, so that all these rows are $(\xi^{g-r})_r$. Let $x_1,\ldots,x_{k_-}$ be the
negative nodes, and put $u_b=J(x_b)$, where $J(z)=\frac12(z+1/z)$. Let $s=\mathsf D+1-k_-$ and $\omega(z)=\prod_p(z-\xi_p)$. If
$k_->\mathsf D+1$, the determinant vanishes, because the negative rows lie in a space of dimension $\mathsf D+1$. So we may assume
$s\ge0$. The Laurent polynomials with exponents in $[-\mathsf D,g]$ that vanish at all $\xi_p$ are spanned by the functions
$z^{i-\mathsf D}\omega(z)$ with $0\le i<k_-$. So elimination gives
\begin{equation}\label{eq:split}
 |\det\Xi^0_S|=\prod_p|\xi_p|^{-\mathsf D}\,|V(\xi)|\,\bigl|\det\bigl(\mathcal Fz^{i-\mathsf D}\omega\,(x_b)\bigr)_{i,b}\bigr|,
\end{equation}
where $\mathcal Ff(x)=\frac12(f(x)+f(1/x))$ is the fold. The fold sends $z^d$ to $T_{|d|}(J(z))$. So the fold of a Laurent
polynomial of degree at most $\mathsf D$ is a polynomial of degree at most $\mathsf D$ in $u=J(z)$.

We measure these functions in a finite Laurent-orthogonal measure. Let $\nu=\nu_N$ be the positive measure on $16$
concentric circles of radii $r_j=e^{-\sigma_j}$, with total masses $e^{-N\tau_j}$, which is uniform on the roots of unity
of order $2\mathsf D+1$ on each circle (scaled by $r_j$). We assume that $\sigma_0=\tau_0=0$ and $10\sigma_j\le\tau_j$ for all
$j$. The Laurent monomials with exponents in $[-\mathsf D,\mathsf D]$ are orthogonal in $L^2(\nu)$, with $\|z^c\|^2=A_c=\sum_je^{-N\tau_j}r_j^{2c}$. Let $G_0=A_0$ and $G_a=(A_a+A_{-a})/4$ for $a\ge1$, which is the squared
norm of $T_a(J(z))$.

\begin{proposition}\label{prop:gram}
Let $\ell_a$ be the leading coefficient of $T_a$, and suppose that the numbers $\theta_d\ge1$ for $-\mathsf D\le d\le g$ satisfy
$G_a\sum_{d}(A_d\theta_d)^{-1}\le1$ for each $a\le \mathsf D$, the sum running over the $d$ with $|d|=a$. Then
\begin{multline*}
 (\det\Xi^0_S)^2\prod_{a=0}^DG_a\cdot k_-!\,s!\\
 \le\Bigl(\prod_p\xi_p^{-\mathsf D}V(\xi)\Bigr)^2\Bigl(\prod_{a=0}^{\mathsf D}\ell_a\,V(u)\Bigr)^2\\
 \cdot\prod_d\theta_d\cdot Z_+Z_-,
\end{multline*}
where
\begin{gather*}
 Z_+=\int\prod_i|z_i|^{-2\mathsf D}|\omega(z_i)|^2|V(z)|^2\,d\nu^{k_-},\\
 Z_-=\int\prod_l|\pi(Jz_l)|^2|V(Jz)|^2\,d\nu^{s},
\end{gather*}
with $\pi(u)=\prod_b(u-u_b)$.
\end{proposition}

\begin{proof}
Let $\mathcal P$ be the space of polynomials of degree at most $\mathsf D$ in $u$, with the inner product for which the $T_a$ are
orthogonal of squared norm $G_a$. Consider the two families $f_i=\mathcal Fz^{i-\mathsf D}\omega$ for $0\le i<k_-$ and $\pi u^l$ for $0\le l<s$.
They have $\mathsf D+1$ members together. The second family spans the kernel of the evaluation at the $u_b$. So the
determinant of the coefficients of the joint family in the basis $(T_a)$ equals $\det(f_i(x_b))$ divided by
$\prod_a\ell_a\cdot V(u)$, up to sign. Multiplying its square by $\prod_aG_a$ gives the Gram determinant of the joint family.
Fischer's inequality bounds it by the product of the Gram determinants of the two families.

Let $f=\sum_df_dz^d$ be a combination of the functions $z^{i-\mathsf D}\omega$. The coefficient of $T_a$ in $\mathcal Ff$ is
$\sum_{|d|=a}f_d$, and the Cauchy--Schwarz inequality with the hypothesis gives $\|\mathcal Ff\|^2\le\sum_d\theta_dA_df_d^2$. Raising
the weights of a diagonal form by factors $\theta_d\ge1$ raises a Gram determinant by at most $\prod_d\theta_d$, by the
Cauchy--Binet formula. So the first Gram determinant is at most $\prod_d\theta_d$ times the Gram determinant of the
functions $z^{i-\mathsf D}\omega$ in $L^2(\nu)$. Andr\'eief's identity \cite{Forrester2018Andreief} in its finite form writes the two Gram
determinants as $Z_+/k_-!$ and $Z_-/s!$. The inequality follows with \eqref{eq:split}.
\end{proof}

\emph{The constants $\theta_d$ and $G_a$.} Put $\theta_d=1$ for $d\le0$ and $\theta_d=A_{-d}/A_d$ for $1\le d\le g$. For $a>g$ the
sum in the hypothesis has the single term $d=-a$, and $G_a/A_{-a}\le\frac12$ because $A_a\le A_{-a}$. For $1\le a\le g$
the left side of the hypothesis is $\frac12(A_a+A_{-a})/A_{-a}\le1$. Since $g=5N$ and $10\sigma_j\le\tau_j$, each term of $A_{-a}$ is at
most $1$, so $A_{-a}\le16$ for $a\le g$. Also $A_a\ge1$ by the term $j=0$. Hence $\prod_d\theta_d\le16^{5N}$. Let
$\tau^*(q)=\max_j(2q\sigma_j-\tau_j)$. Then $G_a\ge\frac14e^{N\tau^*(a/N)}$ for $a\ge1$. Since $\tau^*$ is nondecreasing,
\begin{equation}\label{eq:legendre}
 \sum_{a=0}^{\mathsf D}\log G_a\ge N^2\int_0^{53}\tau^*(q)\,dq-2\mathsf D\log2-N\tau^*(53).
\end{equation}
The factor $\prod_a\ell_a^2$ equals $2^{\mathsf D(\mathsf D-1)}$.

\emph{A tangent bound for Vandermonde determinants.} The integrals $Z_\pm$ are compared with reference measures by the
following lemma, which is the negative definiteness of the logarithmic energy on neutral measures in a form exact
for finitely many points. Let $d(z,w)^2=|z-w|^2/((1+|z|^2)(1+|w|^2))$ and $\mathsf h(z)=\frac12\log(1+|z|^2)$. For a finite positive measure $q$ on
$\C$ write $U^q(z)=\int\log|z-w|\,dq(w)$ and $I(q)=\int U^q\,dq$. Let $H_L=\sum_{j\le L}1/j$.

\begin{lemma}\label{lem:tangent}
Let $q$ be a finite positive measure on $\C$ with $\int\mathsf h\,dq<\infty$ and $q\{w:d(z,w)^2\le t^4\}\le Ct$ for all $z$ and all
$0<t\le1$. For every $\eta>0$ there is $L$ such that for every $m\ge0$, all $\nu_0\ge0$ with $\nu_0q(\C)=m$ and all
$z_1,\ldots,z_m\in\C$,
\[
 |V(z)|^2\le\exp\Bigl(2\nu_0\sum_iU^q(z_i)-\nu_0^2I(q)+2\nu_0m\eta+\tfrac m2H_L\Bigr).
\]
\end{lemma}

\begin{proof}
We may assume that the $z_i$ are distinct. Let $\mathsf n(z)$ be the stereographic image of $z$ on the unit sphere. Then
$\mathsf C(z,w)=1-d(z,w)^2=\frac12(1+\mathsf n(z)\cdot\mathsf n(w))$ is a sum of four products $b_if_i(z)f_i(w)$ with $b_i\ge0$, and so
is every power of $\mathsf C$, with more terms. Let $K_L=-\frac12\sum_{j\le L}\mathsf C^j/j$. Then $K_L=-\sum_\phi b_\phi F_\phi\otimes F_\phi$ is a finite sum with
$b_\phi\ge0$. For real $S_\phi=\sum_iF_\phi(z_i)$ and $M_\phi=\int F_\phi\,dq$ the inequality $(S_\phi-\nu_0M_\phi)^2\ge0$ gives, after summation,
\begin{equation}\label{eq:gramtangent}
 \sum_{i,j}K_L(z_i,z_j)\le2\nu_0\sum_i\int K_L(z_i,w)\,dq(w)-\nu_0^2\iint K_L\,dq\,dq .
\end{equation}
For $z\ne w$ we have $\log d(z,w)=\frac12\log(1-\mathsf C)=-\frac12\sum_{j\ge1}\mathsf C^j/j\le K_L(z,w)$, and $K_L(z,z)=-\frac12H_L$. So
$\sum_{i\ne j}\log d(z_i,z_j)$ is at most the left side of \eqref{eq:gramtangent} plus $\frac m2H_L$. Also $\iint K_L\ge\iint\log d$. For the
remaining term we show that $\int K_L(z,\cdot)\,dq\le\int\log d(z,\cdot)\,dq+\eta$ for large $L$, uniformly in $z$. Fix $h\in(0,1)$.
On $\{d^2\ge h\}$ the difference $K_L-\log d=\frac12\sum_{j>L}\mathsf C^j/j$ is at most $(1-h)^L/(2Lh)$. On $\{d^2<h\}$ it is at most $-\log d$, and
the ball hypothesis on the shells $2^{-4(i+1)}<d^2\le2^{-4i}$ bounds the integral of $-\log d$ over this set by a quantity that tends to
$0$ with $h$, uniformly in $z$. This proves the claim.

Finally $\log|z-w|=\log d(z,w)+\mathsf h(z)+\mathsf h(w)$. So $\sum_{i\ne j}\log|z_i-z_j|=\sum_{i\ne j}\log d(z_i,z_j)+2(m-1)\sum_i\mathsf h(z_i)$. Expanding
$U^q$ and $I(q)$ in the same way, and using $\nu_0q(\C)=m$, the terms in $\int\mathsf h\,dq$ cancel. What remains is
$2m\sum_i\mathsf h(z_i)$ on the right side against $2(m-1)\sum_i\mathsf h(z_i)$ on the left side. Since $\mathsf h\ge0$, this proves the lemma.
\end{proof}

Apply the lemma to the $a+c$ points consisting of $a$ fixed base points $b_i$ and $c$ points $Z_l$. If $Z_l$ runs over a finite set with
weights $\mathrm{wt}\ge0$, summing over tuples gives, with $\nu_0q(\C)=a+c$ and the same $\eta$ and $L$,
\begin{multline}\label{eq:tangent}
 |V(b)|^2\sum_{Z}\prod_l\mathrm{wt}(Z_l)\,\prod_{l,i}|Z_l-b_i|^2\,|V(Z)|^2\\
 \le\exp\Bigl(2\nu_0\sum_iU^q(b_i)-\nu_0^2I(q)+2\nu_0(a+c)\eta+\tfrac{a+c}2H_L\Bigr)\Bigl(\sum_z\mathrm{wt}(z)e^{2\nu_0U^q(z)}\Bigr)^c .
\end{multline}
The exponential on the right does not depend on the $Z_l$, and the last sum separates them. This is the reason for completing the
gas by the base points: the determinant of the base points is then bounded together with the gas.

\emph{The two gases.} Let $q_+$ be a measure of mass $1$ in the plane of $z$, and $q_-$ a measure of mass $\frac{53}{58}$ in the plane of
$u$, both satisfying the hypotheses of Lemma~\ref{lem:tangent}. Let $\kappa_+$ and $c_-$ be constants. We assume the angular bounds
\begin{equation}\label{eq:angular}
 U^{q_+}(r_je^{i\theta})+\tfrac{53}{58}\sigma_j-\tfrac{\tau_j}{116}\le\kappa_+,\qquad
 U^{q_-}\bigl(J(r_je^{i\theta})\bigr)-\tfrac{\tau_j}{116}\le c_-,
\end{equation}
for all $j$ and $\theta$. Apply \eqref{eq:tangent} to $Z_+$ with $q=q_+$, $\nu_0=n$, base points $\xi_p$ and the $k_-$ gas points $z_i$. Since
$r_j^{-2\mathsf D}\le e^{2\cdot53N\sigma_j}$, the sum in \eqref{eq:tangent} is at most $16e^{2n\kappa_+}$. Apply \eqref{eq:tangent} to $Z_-$ with $q=q_-$,
$\nu_0=n$, base points $u_b$ and the $s$ gas points $Jz_l$. The sum is at most $16e^{2nc_-}$. Insert these in
Proposition~\ref{prop:gram}, use \eqref{eq:legendre}, and take logarithms. Each positive node contributes
$\mathsf g(\xi)=-\mathsf D\log|\xi|+n(U^{q_+}(\xi)-\kappa_+)$ and each negative node $\mathsf g(x)=n(U^{q_-}(Jx)-c_-)$. We have proved the following statement.

\begin{proposition}\label{prop:gas}
For every $\eta>0$ there is $C$ such that for all $N$, all nodes and all $S$ with $\det\Xi^0_S\ne0$,
\[
 2\log|\det\Xi^0_S|\le2\sum_a\mathsf g(a)+\mathsf K_N+4n^2\eta+Cn,
\]
where $\mathsf g(a)$ is the contribution of node $a$ and
\[
 \mathsf K_N=2n^2\kappa_++2\cdot53Nn\,c_-+\mathsf D(\mathsf D+1)\log2-N^2\int_0^{53}\tau^*-n^2\bigl(I(q_+)+I(q_-)\bigr).
\]
\end{proposition}

The constant is correct because the $k_-$ points of $Z_+$ give $2nk_-\kappa_+$ and the $s$ points of $Z_-$ give $2nsc_-$. In $\mathsf K_N$ these are
written as $2n^2\kappa_+$ and $2\cdot53Nnc_-$, and the differences are absorbed by the terms $-n\kappa_+$ and $-nc_-$ of $\mathsf g$. The term $\mathsf D(\mathsf D+1)\log2$
is the sum of $\mathsf D(\mathsf D-1)\log2$ from $\prod_a\ell_a^2$ and $2\mathsf D\log2$ from \eqref{eq:legendre}.

\subsection{The final inequality}\label{ssec:budget}
We now bound the other factors of the integrand. Let $n'=n-m$, with free points $x_i$ and $s_j$ and prescribed points $y_j$, and let $\mathcal S$ be the product of the
factors of the integrand other than the row determinant and the factors $\mathcal B_\zeta$ of Lemma~\ref{lem:binomvand}. Put $\mathsf t(x)=2x/(1+x^2)$ and $t_i=\mathsf t(x_i)$. In the normalization of
\S\ref{ssec:rowdet} a free point carries the weight $|t|^{75N-1}(1-t)^{3N}$. The prescribed columns have lost the common factor $|t|^{22N}(1-t)^{3N}$ of the rows, so a prescribed point
carries $|\mathsf t(y)|^{\mathsf D}$. The Cauchy determinant of Lemma~\ref{lem:andreief} is $\pm V(t)V(s)/\prod_{i,j}(1-t_is_j)$, and the generalized Vandermonde determinant gives
$\prod_js_j^{9N}(1-s_j)^{3N}|V(s)|$. Hence
\begin{multline*}
 \mathcal S=\prod_j|\mathsf t(y_j)|^{\mathsf D}\prod_i|t_i|^{75N-1}(1-t_i)^{3N}\\
 \cdot\prod_js_j^{9N}(1-s_j)^{3N}\,\frac{|V(t)|\,V(s)^2}{\prod_{i,j}(1-t_is_j)}.
\end{multline*}
Now $1-t=(1-x)^2/(1+x^2)$, $1-ts=(1-2xs+x^2)/(1+x^2)$ and
\[
 t_i-t_{i'}=\frac{2(x_i-x_{i'})(1-x_ix_{i'})}{(1+x_i^2)(1+x_{i'}^2)}.
\]
Each $x_i$ occurs in $n'-1$ differences and in $n'$ kernels, so the exponent of $1+x_i^2$ is $-(75N-1)-3N-(n'-1)+n'=-(78N-2)$. Taking logarithms gives
\begin{multline}\label{eq:scalar}
 \log\mathcal S=\Bigl(n'(75N-1)+\binom{n'}2\Bigr)\log2+\sum_j\mathsf w(y_j)+\sum_i\mathsf W_N(x_i)\\
 +\sum_j\bigl(9N\log s_j+3N\log(1-s_j)\bigr)+n^2\mathsf Q(x,s),
\end{multline}
where $\mathsf w(y)=\mathsf D(\log2+\log|y|-\log(1+y^2))$, $\mathsf W_N(x)=(75N-1)\log|x|+6N\log(1-x)-(78N-2)\log(1+x^2)$ and
\begin{multline*}
 n^2\mathsf Q(x,s)=\sum_{i<i'}\log\bigl(|x_i-x_{i'}|(1-x_ix_{i'})\bigr)\\
 +2\sum_{j<j'}\log|s_j-s_{j'}|-\sum_{i,j}\log(1-2x_is_j+x_i^2).
\end{multline*}
This is the energy of \S\ref{ssec:energy} with the interaction among the $x_i$ at weight $1$. The second copy of that interaction is
contained in the determinant.

For sequences $a=(a_k)$ and $v=(v_k)$ put $\|a\|^2=\sum_ka_k^2/k$. Let $b$ be a real number.

\begin{lemma}\label{lem:freetangent}
Let $a,v$ be sequences with $\sum|a_k|/k$, $\sum|v_k|/k$, $\|a\|$ and $\|v\|$ finite. For every $\varepsilon>0$ there is $\nu_1$ such that
for all $n\ge\nu_1$, all $n'\le n$ and all configurations of distinct points,
\begin{multline*}
 \mathsf Q(x,s)\le\tfrac32b^2\log2+\|a\|^2+\tfrac12\|v\|^2-\tfrac32\bigl(\tfrac{n'}n-b\bigr)^2\log2\\
 +\frac1n\sum_i\mathsf R_*(x_i)+\frac1n\sum_j\mathsf C_*(s_j)\\
 -\frac1{2n^2}\sum_i\log(1-x_i)-\frac1{n^2}\sum_j\log(1-s_j)+\varepsilon,
\end{multline*}
where $\mathsf R_*(x)=-b\log2-2\mathsf T(a,x)-\mathsf S(v,x)$ and $\mathsf C_*(s)=-2b\log2+2\mathsf T(v,s)$.
\end{lemma}

\begin{proof}
Use the expansions $-\log|x-x'|=\log2+2\sum_kT_k(x)T_k(x')/k$, $-\log(1-xx')=\sum_kx^kx'^k/k$ and
$-\log(1-2xs+x^2)=2\sum_kx^kT_k(s)/k$. Let $A_k=\frac1n\sum_iT_k(x_i)$, $B_k=\frac1n\sum_ix_i^k$ and $C_k=\frac1n\sum_jT_k(s_j)$, and let $f=n'/n$.
If the diagonal terms are ignored, these expansions give
\[
 \mathsf Q=-\tfrac32f^2\log2-\|A\|^2-\tfrac12\|B-2C\|^2 .
\]
The diagonal terms diverge. OpenAI's regularization \cite[\S5]{OpenAI2026} replaces the kernels by Abel-damped series with a parameter $e$.
It gives the identity above for the damped moments, up to an error at most
$-\frac1{2n^2}\sum_i\log(1-x_i)-\frac1{n^2}\sum_j\log(1-s_j)+O_e(1/n)+14e$. For the trial sequences, $-\|A\|^2\le\|a\|^2-2\sum_ka_kA_k/k$ and
$-\frac12\|B-2C\|^2\le\frac12\|v\|^2-\sum_kv_k(B_k-2C_k)/k$. Removing the damping from the linear terms costs $O(e)$ times the $\ell^1$ norms of the trial
sequences, so we first replace $a$ and $v$ by long prefixes, with an error at most $\varepsilon/4$ uniformly in the points. Finally
$-\frac32f^2\log2=-\frac32(f-b)^2\log2+\frac32b^2\log2-3bf\log2$, and the last term is split between $\mathsf R_*$ and $\mathsf C_*$. Choosing the prefix, then $e$, then $\nu_1$
proves the lemma. The order of these choices makes the bound uniform in the configuration.
\end{proof}

We apply the lemma with $b=1$. Let $X$ be the inverse of $x\mapsto x+1/x$ with values in $[-1,1]\setminus\{0\}$. Let $R$ and $C$ be the row and column references of Appendix~\ref{app:trials},
and let $L(\mu)=\int\log|X|\,d\mu$. We take
\[
 a_k=\int T_k(X)\,dR,\qquad v_k=\int X^k\,dR-2\int T_k\,dC,
\]
and put $\mathsf I_0=\frac32\log2+\|a\|^2+\frac12\|v\|^2$. The expansions in the proof of Lemma~\ref{lem:freetangent} give
\begin{gather*}
 \mathsf R_*(x)=\int\bigl(\log|x-X|+\log(1-xX)\bigr)\,dR-\int\log(1-2xs+x^2)\,dC(s),\\
 \mathsf C_*(s)=2U^C(s)-\int\log(1-2Xs+X^2)\,dR.
\end{gather*}
For nonzero $x$ and $w$ and for $0<s<1$ we have
\begin{equation}\label{eq:jouk}
 x+\frac1x-w-\frac1w=\frac{(x-w)(xw-1)}{xw},\qquad 1-2xs+x^2=|x|\,\bigl|x+\tfrac1x-2s\bigr|.
\end{equation}
Since $R$ and $C$ have mass $1$, this gives
\[
 \mathsf R_*(x)=L(R)+U^{R-2_\#C}(x+1/x),\qquad \mathsf C_*(s)=-L(R)+2U^C(s)-U^R(2s).
\]
So $9N\log s+3N\log(1-s)+n\mathsf C_*(s)=n\mathsf C(s)$ for the column field $\mathsf C$ of the appendix.

\emph{The moved sequences.} Lemma~\ref{lem:freetangent} holds for any sequences, so we may move $a$ and $v$ with the
fraction $\alpha$. Let $\mathrm da^0,\mathrm da^1,\mathrm dv^0,\mathrm dv^1$ be fixed sequences supported on $k\le160$, with
dyadic entries of denominator $2^{40}$. Put $\alpha_e=\frac1{25}$, $\tau=\min(1,\alpha/\alpha_e)$ and
\[
 a(\tau)=a+\mathrm da^0+\tau\,\mathrm da^1,\qquad v(\tau)=v+\mathrm dv^0+\tau\,\mathrm dv^1 .
\]
The increment $(\mathrm da^0,\mathrm dv^0)$ moves the sequences already at $\alpha=0$. The fields change by finite sums:
$\mathsf R_*$ gains $h^0+\tau h^1$ with $h^i=-2\mathsf T(\mathrm da^i,\cdot)-\mathsf S(\mathrm dv^i,\cdot)$, and $\mathsf C_*$ gains
$k^0+\tau k^1$ with $k^i=2\mathsf T(\mathrm dv^i,\cdot)$. The polarization identity gives
\[
 \mathsf I_0(\tau)=\mathsf I_0+(L_0+Q_{00})+\tau(L_1+2Q_{01})+\tau^2Q_{11},
\]
with $L_i=2\langle a,\mathrm da^i\rangle+\langle v,\mathrm dv^i\rangle$ and
$Q_{ij}=\langle\mathrm da^i,\mathrm da^j\rangle+\frac12\langle\mathrm dv^i,\mathrm dv^j\rangle$. Only $L_0$ and $L_1$ involve the
references. We also include the column field $\gamma g$ of \S\ref{ssec:prescribed}, with $\gamma=\tau\alpha_e$ for $\alpha\le\alpha_e$
and $\gamma=\alpha$ for $\alpha_e\le\alpha\le\alpha_1$, so that $\gamma\ge\alpha$. The moved row field at $\tau$ is the convex
combination $(1-\tau)(F_\pm+h^0)+\tau(F_\pm+h^0+h^1)$, and the moved column field is affine in $\tau$ and $\gamma$. So five
numbers bound all of them, the bounds at the end points:
\begin{gather*}
 F_\pm+h^0\le c^{\mathsf R}_0,\qquad F_\pm+h^0+h^1\le c^{\mathsf R}_1,\\
 \mathsf C+k^0\le c^{\mathsf C}_0,\qquad \mathsf C+k^0+k^1+\alpha_eg\le c^{\mathsf C}_1,\qquad \mathsf C+k^0+k^1+\alpha_1g\le c^{\mathsf C}_2 .
\end{gather*}
For $\alpha\le\alpha_e$ the row and column caps are the chords $(1-\tau)c^{\mathsf R}_0+\tau c^{\mathsf R}_1$ and
$(1-\tau)c^{\mathsf C}_0+\tau c^{\mathsf C}_1$. For $\alpha_e\le\alpha\le\alpha_1$ they are $c^{\mathsf R}_1$ and the chord of
$c^{\mathsf C}_1$ and $c^{\mathsf C}_2$ at $\alpha$. No supremum of an increment is ever taken separately from the supremum of the base field.

We turn to the gas potentials. Let $\nu_\pm$, $A_\pm$ and $\tilde q_-$ be as in Appendix~\ref{app:trials}, so that $q_+=X_\#G_++\nu_+$ and $q_-=J_\#\tilde q_-$ with $\tilde q_-=X_\#G_-+\nu_-$.
Let $0<x<1$ and $Y=x+1/x$. By \eqref{eq:jouk}, $\log|x-X(Y')|=\log|Y-Y'|+\log x+\log|X(Y')|-\log(1-xX(Y'))$. On the support of $G_+$ we have $X>0$, so
$-\log(1-xX)=\sum_kx^kX^k/k$ has positive terms. Those with $k\le900$ integrate to $\mathsf P(x)$ and the others to at most $\mathsf E$. Hence
\begin{equation}\label{eq:qplus}
 U^{q_+}(x)\le L(G_+)+U^{G_+}(Y)+m_+\log x+\mathsf P(x)+\mathsf E+A_+(x).
\end{equation}
Let now $-1<x<0$. We have $\log|Jx-Y'/2|=\log|Y-Y'|-\log2$ on the support of $G_-$. By \eqref{eq:jouk},
$\log|Jx-Jw|=\log|x-w|+\log|1-xw|-\log2-\log|x|-\log|w|$ on the support of $\nu_-$. Put $c_M=\int\log|u|\,d\tilde q_-(u)$. Integrating, and using the mass $\frac{53}{58}$ of $\tilde q_-$, we find
\begin{multline}\label{eq:qminus}
 U^{q_-}(Jx)=L(G_-)+U^{G_-}(Y)+m_-\log|x|+A_-(x)\\
 -\tfrac{53}{58}\bigl(\log2+\log|x|\bigr)-c_M.
\end{multline}
Define $\kappa_-$ by $c_-=\kappa_--\frac{53}{58}\log2-c_M$. Since $n\cdot\frac{53}{58}=53N$, this says
\begin{equation}\label{eq:gneg}
 \mathsf g(x)=n\bigl(L(G_-)+U^{G_-}(Y)+m_-\log|x|+A_-(x)-\kappa_-\bigr)-53N\log|x|.
\end{equation}

Let $F_\pm$ be the free row fields of the appendix. Since $n\cdot\frac{11}{29}=22N$, $n\cdot\frac{39}{29}=78N$ and $n\cdot\frac3{29}=6N$, the definition reads
\begin{multline*}
 nF_\sigma=n\mathsf R_*+n\bigl(L(G_\sigma)+U^{G_\sigma}(Y)+m_\sigma\log|x|-\kappa_\sigma\\
 +[\sigma=+](\mathsf P+\mathsf E)+A_\sigma\bigr)+22N\log|x|-78N\log(1+x^2)+6N\log(1-x).
\end{multline*}
Recall that $\mathsf g(x)=-\mathsf D\log x+n(U^{q_+}(x)-\kappa_+)$ for $x>0$. Add $\mathsf W_N$ to this and use \eqref{eq:qplus} for $0<x<1$ and \eqref{eq:gneg} for $-1<x<0$. The terms in $\log(1-x)$
cancel, the exponents of $1+x^2$ leave $2$, and the coefficient of $\log|x|$ is $75N-1-22N-\mathsf D=0$ in the first case and $75N-1-22N-53N=-1$ in the second. So
\begin{gather}
 \mathsf W_N(x)+n\mathsf R_*(x)+\mathsf g(x)\le nF_+(x)+2\log(1+x^2),\label{eq:rowpos}\\
 \mathsf W_N(x)+n\mathsf R_*(x)+\mathsf g(x)=nF_-(x)+2\log(1+x^2)-\log|x|.\label{eq:rowneg}
\end{gather}
Let $\mathsf R$ and $\mathsf C_0$ be upper bounds for the moved row and column fields at the given $\alpha$. The added fields are bounded, so they do not affect the following argument, which we write for $F_\pm$ and $\mathsf C$.
The term $-\frac12\log(1-x)$ of Lemma~\ref{lem:freetangent} and the term $-\log|x|$ of \eqref{eq:rowneg} are unbounded. They are controlled by the field itself. Let $H_0$ be the supremum of $F_+-\frac3{29}\log(1-x)-\frac{11}{29}\log|x|$ on $(0,1)$ and of $F_--\frac3{29}\log(1-x)-\frac{11}{29}\log|x|$ on $(-1,0)$. It is finite.
Indeed, after the subtraction $F_\sigma$ consists of constants, the term $-\frac{39}{29}\log(1+x^2)\le0$, the polynomial $\mathsf P$, the term $A_\sigma$ and
$U^{G_\sigma+R-2_\#C}(x+1/x)+m_\sigma\log|x|$. The term $A_\sigma$ is bounded because the auxiliary measures have bounded densities on circles. The measures $G_\sigma$, $R$ and $2_\#C$
are piecewise uniform with compact support, so their potentials are continuous on the real line. This settles the limits $x\to\pm1$, where $x+1/x\to\pm2$, and every compact subset of
$(-1,1)\setminus\{0\}$. As $x\to0$ we have $|x+1/x|\to\infty$. The measure $R-2_\#C$ has mass $0$ and $G_\sigma$ has mass $m_\sigma$, so $U^{G_\sigma+R-2_\#C}(x+1/x)+m_\sigma\log|x|$ tends to $0$.
Now split $nF=(n-\frac{29}6)F+\frac{29}6F$. Since $\frac{29}6\cdot\frac3{29}=\frac12$ and $\frac{29}6\cdot\frac{11}{29}=\frac{11}6\ge1$, and since $\log|x|\le0$,
\[
 nF-\tfrac12\log(1-x)-\log|x|\le(n-\tfrac{29}6)\mathsf R+\tfrac{29}6H_0 .
\]
The columns are treated in the same way. Let $H_1$ be the supremum of $\mathsf C(s)-\frac3{58}\log(1-s)$ on $(0,1)$. It is finite, because $U^C$ and $U^R(2s)$ are continuous and
$\log s\le0$. Since $\frac{58}3\cdot\frac3{58}=1$,
\[
 n\mathsf C(s)-\log(1-s)\le(n-\tfrac{58}3)\mathsf C_0+\tfrac{58}3H_1 .
\]
So every free point contributes at most $n\mathsf R+O(1)$ and every free column at most $n\mathsf C_0+O(1)$. 

Let now $y$ be a prescribed point. All prescribed points lie in $E_x$, so $y\ge98749/10^6$ or $y\le-19501/10^5$. Let $\mathsf H_+$, $\mathsf H_-$ and $\mathsf H_S$ be the fields of the prescribed rows in
the appendix. For $y>0$ outside $S$, \eqref{eq:qplus} and $53N=\mathsf D+1$ give
\[
 \mathsf w(y)+\mathsf g(y)\le \mathsf D\log2+n\mathsf H_+(y)+\log(1+y^2).
\]
For $y<0$, \eqref{eq:gneg} gives
\[
 \mathsf w(y)+\mathsf g(y)=\mathsf D\log2+n\mathsf H_-(y)-\log|y|+\log(1+y^2).
\]
For $y\in S$ the base point is $\xi=1/y$, so $\mathsf g=\mathsf D\log y+n(U^{q_+}(1/y)-\kappa_+)$. We have $\log|1/y-w|=-\log y+\log|1-yw|$, the measure $q_+$ has mass $1$, and $\int\log(1-yX)\,dG_+\le-\mathsf P(y)$
because all the moments are nonnegative. Hence $U^{q_+}(1/y)\le-\log y-\mathsf P(y)+\int\log|1-yw|\,d\nu_+(w)$. With $2\mathsf D=106N-2$ this gives
\[
 \mathsf w(y)+\mathsf g\le \mathsf D\log2+n\mathsf H_S(y)-2\log y+\log(1+y^2).
\]
The factor $e^{-N\hat\delta(\mathsf t(y))}$ of Proposition~\ref{prop:signed-goal} adds $\hat\delta(\mathsf t(y))/58$ to these fields. Let $\mathsf M$ be an upper bound for $\mathsf H_\pm+\hat\delta\circ\mathsf t/58$ on $E_W$ and for $\mathsf H_S+\hat\delta\circ\mathsf t/58$ on $E_W\cap(0,1)$. Since $\log(1+y^2)\le\log2$ and $|y|$ is bounded below on $E_x$, every prescribed row, with its factor $e^{-N\hat\delta}$, contributes at most
$\mathsf D\log2+n\mathsf M+O(1)$.

Now add \eqref{eq:scalar}, Lemma~\ref{lem:freetangent} multiplied by $n^2$, and half of Proposition~\ref{prop:gas}. With $n'=(1-\alpha)n$ and $\eta,\varepsilon\to0$, this gives
\[
 \log\bigl(\mathcal S\,|\det\Xi^0_S|\bigr)\le n^2\bigl(\mathsf B(\alpha)+o(1)\bigr),
\]
uniformly in the configuration, with
\begin{multline}\label{eq:budget}
 \mathsf B(\alpha)=\mathsf I-\tfrac32\alpha^2\log2+(1-\alpha)(\mathsf R+\mathsf C_0)+\alpha\mathsf M\\
 +\log2\Bigl(\tfrac{75}{58}(1-\alpha)+\tfrac12(1-\alpha)^2+\tfrac{53}{58}\alpha\Bigr).
\end{multline}
Here $\mathsf I=\mathsf I_0(\tau)+\mathsf I_{\mathrm{gas}}$, and $\frac12\mathsf K_N/n^2$ tends to
\[
 \mathsf I_{\mathrm{gas}}=\kappa_++\tfrac{53}{58}(\kappa_--c_M)-\tfrac12\bigl(\tfrac{53}{58}\bigr)^2\log2-\tfrac12\bigl(I(q_+)+I(q_-)\bigr)-\frac1{2\cdot58^2}\int_0^{53}\tau^*,
\]
with $\kappa_-$ and $c_M$ as above. The three terms in $\log2$ of \eqref{eq:budget} come from the exponent
$n'(75N-1)+\binom{n'}2$ and from the term $\mathsf D\log2$ of each prescribed row. Appendix~\ref{app:trials} gives the reference measures and the bounds
\begin{equation}\label{eq:caps}
\begin{gathered}
 \mathsf I_0+\mathsf I_{\mathrm{gas}}\le2.366997,\qquad\mathsf M\le-0.976,\\
 c^{\mathsf R}_0=-2.579740254,\quad c^{\mathsf R}_1=-2.574009165,\\
 c^{\mathsf C}_0=-2.946694099,\quad c^{\mathsf C}_1=-2.888907917,\quad c^{\mathsf C}_2=-2.85421372,
\end{gathered}
\end{equation}
with $\kappa_-=5463909/10^8$ and $\kappa_+=14762657/10^8$, and the moved values
\begin{gather*}
 L_0\le-7.5421576\cdot10^{-5},\quad L_1\le-0.04943292757,\\
 Q_{00}\le1.6623\cdot10^{-7},\quad Q_{01}\le9.5232\cdot10^{-7},\quad Q_{11}\le1.002863\cdot10^{-3}.
\end{gather*}
Let $V_S(\alpha)=\mathsf B(\alpha)+U_0-\frac12\log2-c_2\alpha$, with $\mathsf I$, $\mathsf R$ and $\mathsf C_0$ the moved values above and $\log2$
replaced by an upper bound. On each of the two pieces $[0,\alpha_e]$ and $[\alpha_e,\alpha_1]$ it is a polynomial in $\alpha$ with rational
coefficients. We have $V_S(0)=-1.81\ldots\cdot10^{-5}$, and the cover of \S\ref{ssec:cases} proves $V_S(\alpha)\le\sigma\alpha$ on both pieces. The
inequality is tight near $\alpha=0.0033$, where $V_S(\alpha)/\alpha$ reaches $0.33867$.

\begin{proof}[Proof of Proposition~\ref{prop:signed-goal}]
Fix a configuration with $m\le\alpha_1n$. By \eqref{eq:expandS} it suffices to bound each of the at most $2^m$ terms
$\mathcal S\,|\det\Xi_S|$, at the cost of the factor $5^m$. By Proposition~\ref{prop:alternative} each term either satisfies the
bound directly, or is at most $499^n\mathcal S\,|\det\Xi^0_S|$. In the second case the final inequality gives
$\exp\bigl(n^2(-U_0+\frac12\log2+c\alpha+o(1))-N\sum_j\hat\delta(t_{a_j})\bigr)$, since $499^n=e^{o(n^2)}$ and $c=c_2+\sigma$.
The factor $\prod_l(1+\zeta s_l)^m$ of Lemma~\ref{lem:binomvand} is contained in the column field, and the factors $\mathcal B_\zeta$ give
$\prod_j\mathcal B_\zeta(k_j)$.
\end{proof}
\section{Proof of Theorem~\ref{thm:main}}\label{sec:proof}

OpenAI's energy method gives the decay rate $D_0=n_0^2(U-\frac12\log2)=\DecayOA\ldots$ for our parameters.
Section~\ref{sec:signed} raises it to $D=\Decay\ldots$ With $\Lambda=\LamDen\ldots$ the polynomial
$Q_N=(d_0N)!\,\cD_N\Delta_N$ has integer coefficients, degree at most $53N$, and
\[
 |Q_N(x)|\le e^{-\eta N^2}\qquad(|x-G|\le e^{-\rho N}),\qquad \eta=\EtaTot .
\]
Two further properties are needed to choose the prime in the criterion: nonvanishing of $Q_p(a/b)$ when
$p^2\nmid b$ (Proposition~\ref{prop:nonvanishing}), and divisibility of the $i$-th coefficient of $Q_N$ by
$p'^{2i}$ for every prime $N<p'\le2N$ (Lemma~\ref{lem:squares}). We now check the hypotheses of
Proposition~\ref{prop:criterion}.

Take $d=53$, and let the polynomials $Q_p$ of the proposition be the $Q_N$ with $N=p$. Let
\begin{gather*}
 \rho=\Rho,\quad \eta=\EtaTot,\quad \omega=\tfrac35,\quad \delta=2.9\cdot10^{-7},\\
 \lambda=\frac{270}{\rho}=\LamCrit\ldots
\end{gather*}
\begin{enumerate}[label=(\alph*)]
\item $\deg Q_N\le53N$ by \eqref{eq:QN} and Lemma~\ref{lem:degree}.
\item By Proposition~\ref{prop:nonvanishing}, $\Delta_p(a/b)\ne0$ for large $p$ with $p^2\nmid b$. Hence
$Q_p(a/b)\ne0$.
\item By Lemma~\ref{lem:squares}, $p'^{2i}$ divides the $i$-th coefficient of $Q_p$ for all primes
$p<p'\le2p$.
\item By Theorem~\ref{thm:real}, Proposition~\ref{prop:denominators} and $\log(53N)!=o(N^2)$, for
$|x-G|\le e^{-\rho N}$ and large $N$,
\[
 |Q_N(x)|\le\exp\bigl(-(D-\Lambda+o(1))N^2\bigr)\le e^{-\eta N^2},
\]
since $\eta=\EtaTot<D-\Lambda=\EtaExact\ldots$ In the formal proof this is the inequality $\eta<D-\Lambda$ with
$\int\kappa$ replaced by its lower bound $\ChiLean$ and $D_0-\Lambda_0$ by its lower bound $7.98775$.
\end{enumerate}
The numerical conditions hold as well. We have $\rho\lambda=270$ and $0<\eta<2d=106$. Next $\omega=\frac35<\lambda$ and
$\delta<\frac13$. The third condition reads $53\le(\lambda-\omega)(\eta+106(1-3\delta))=59.43\ldots$ Further
$2(\omega-\delta(2\lambda-\omega))=1.199999\ldots>1$. Finally
$d+4d\delta\lambda=\PayL\ldots<\eta\lambda=\PayR\ldots$ So Proposition~\ref{prop:criterion} gives
$|G-a/b|>b^{-\rho\lambda}=b^{-270}$ for all integers $a$ and all sufficiently large $b$. This proves
Theorem~\ref{thm:main}.

\section{Remarks on the method}\label{sec:remarks}

\subsection{The three ingredients}
For fixed parameters the method gives the exponent $d\rho/\eta$, and its three ingredients come from
different places. The degree $d=n-g$ comes from the construction. Of $\rho=\Rho$, the prescribed columns
contribute $n_0c\approx228.3$ through the energy, and the weight and the Vandermonde factor of one column
contribute $\psi\approx50.1$. Of $n_0c$, the part $n_0\sigma\approx19.6$ comes from the bound of
Section~\ref{sec:signed} for few prescribed columns. The number $\eta$ is a small difference of two large numbers, $D\approx6.4\cdot10^3$ and
$\Lambda\approx6.4\cdot10^3$. The corank bounds of Section~\ref{sec:content} contribute about $22$ to it, and the
bound of Section~\ref{sec:signed} about $25$. Without them $\eta$ would be about $8$. Since the exponent
is inversely proportional to $\eta$, a further gain of $0.2$ in $\eta$ would lower it by about $1$.

\subsection{The row determinant}
\begin{remark}\label{rem:L}
The energy method bounds $|\det\mathsf L|$ by $e^{O(n\log n)}$, which is invisible at the scale $e^{N^2}$.
At the extremal configuration, numerical evaluation for small $N$ suggests that $\log|\det\mathsf L|$ is about
$-28N^2$. So our bound captures most of its true size. The comparison of Lemma~\ref{lem:compare} loses only
$e^{O(N)}$, and so do the Gram inequalities of Proposition~\ref{prop:gram}. The loss is in the restriction to
Gram measures on finitely many circles, and in the requirement that the bound hold for the maximum over
all configurations. These numerical values are not part of the formal proof.
\end{remark}

\subsection{Other parameters}
All the results of Section~\ref{sec:content} are lower bounds for the corank. The closed formulas, the
counting inequality, Theorem~\ref{thm:everyU} and the families derived from them are proved for all
parameters. The rank data of Appendix~\ref{app:windows} and the reference measures belong to $(58,9,3,5,3)$, and
Sections~\ref{sec:real}--\ref{sec:signed} are written and formalized for these parameters. Theorem~\ref{thm:window} does not depend on
the parameters. A different choice of parameters changes $d$, $\rho$ and $\eta$ together, and we make no claim
that ours is optimal. Except for the pair (F4) and the intervals of Theorem~\ref{thm:everyU}, the function
$\kappa$ uses no relation between different families. So at a prime where intervals overlap the exact corank
may be larger than each bound. An upper bound for the corank is not needed. More independence between
families would increase $\eta$ without changing anything else.

\section{Formalization}\label{sec:formal}

Theorem~\ref{thm:main} and the results it depends on are formalized in Lean~4 \cite{Lean4} with Mathlib \cite{Mathlib}, on top of OpenAI's
formalization of the irrationality of $G$ \cite{OpenAI2026}. The development is available at
\begin{center}
\FormalizationLink.
\end{center}
Theorem~\ref{thm:main} is the theorem
\begin{center}
{\small\texttt{CatalanIrrationalityMeasure.catalan\_irrationalityMeasure\_270}}
\end{center}
in \lean{Solution.lean}. Its statement, in \lean{Challenge.lean}, uses only Mathlib. Catalan's constant is
written as the series $\sum_j(-1)^j/(2j+1)^2$, and the conclusion is $|G-a/b|>b^{-270}$ for all integers $a$
and all sufficiently large $b$. The modules of the proof contain no \lean{sorry} and no \lean{native\_decide}. The command
\lean{\#print axioms} shows that only the three standard axioms \lean{propext}, \lean{Classical.choice} and
\lean{Quot.sound} are used. The proof is kernel-checked: Lean's kernel accepts every declaration. It has not been run through the independent checker \lean{lake comparator} \cite{Comparator}.
The prime number theorem is taken from the PNT+ project \cite{PNTPlus}. The parts of
OpenAI's formalization that are used were converted to Lean's module system and ported to Lean
4.35.0-rc2. The repository's \texttt{README} describes the changes and the layout of the modules.

Two lemmas of Section~\ref{sec:real} on the weights of nodal polynomials come from Davis's formalization
\cite{Davis2026}. We use two of its Lean files, taken from commit \texttt{3c45153} under the Apache
License~2.0. Their declarations and proofs are unchanged. We adapted the module paths and imports, and added comments and
axiom checks. The repository's \texttt{NOTICE} file and a provenance record list these files with their source
hashes.

The proof combines conventional formal proofs with numerical certificates. These are the inverses of the base
matrices modulo $1009$, the intervals and the chord table of Section~\ref{sec:real}, the trial sequences, the
cell covers and the bounds of Appendix~\ref{app:field}, the reference measures
and the bounds of Appendix~\ref{app:trials}, and the rank data of Appendix~\ref{app:windows}. The certificates are
exact rational and integer computations checked by the Lean kernel without compiled code. The interval
arithmetic used to generate them \cite{Arb} is not trusted.

\section*{Disclosure of AI use}
The mathematics, the Lean formalization, the certificate generators and the text of this paper were produced
by AI systems at the request of the author. The author posed the problems, directed the work and is
responsible for the paper. The systems were Claude Opus~5.5, Claude Sonnet~5.5 and Claude Haiku~5.5
(Anthropic), running in Claude Code, and GPT-6.1 Sol (OpenAI), running in Codex. The underlying
irrationality proof for Catalan's constant and its Lean formalization are due to OpenAI \cite{OpenAI2026}.

\subsection*{Acknowledgements}
We thank the authors of Mathlib and of the PNT+ project, OpenAI for making its formalization available, and
Cinna Davis for making the formalization of \cite{Davis2026} available.

\appendix
\section{Counting on every subspace}\label{app:windows}

On $37$ of the $113$ intervals of Theorem~\ref{thm:profile-D} the corank bound comes from the counting
inequality applied to all constant subspaces at once. These intervals lie in $(3N,16.2N)$. No subspace is
constructed. The proof bounds the maximum over subspaces of the right side of Theorem~\ref{thm:everyU} by forced intersections of row spaces, and the
only finite data are upper bounds for ranks.

Let $V$ be a finite-dimensional vector space over $\F_p$. Let $c$ be a nonnegative integer. For $j$ in a finite
set $\mathcal J$ let $B_j\colon V\to W_j$ be a linear map and $w_j$ a nonnegative integer, and let $\lambda_{j,k}$ ($k<w_j$) be linear forms on
$\F_p^c$. Let $\mathcal C$ be the common kernel of the maps $B_j\circ(\lambda_{j,k}\otimes\mathrm{id})$ on $\F_p^c\otimes V$. A \emph{parametrization} is a linear map
$\iota$ from $\F_p^c\otimes V$ to the data $x$ of a decomposition modulo $p$ (Definition~\ref{def:slot}) with three properties on $\mathcal C$. The value
$\iota(x)$ is admissible, the map $\iota$ is injective on $\mathcal C$, and the row combination $c(\iota(x))$ satisfies the hypothesis of
Theorem~\ref{thm:count} with some polynomial of degree less than $r$. Nothing is required outside $\mathcal C$.

\begin{theorem}\label{thm:everyU}
Let $N\ge1$ and let $p$ be an odd prime with $H<p^2$. For a parametrization and every subspace $U\subset V$,
\[
 \nu_N(p)\ge c\dim U-\sum_{j\in\mathcal J}w_j\rank(B_j|_U).
\]
\end{theorem}

\begin{proof}
Since the coefficients $\lambda_{j,k}$ are scalars, $\lambda_{j,k}\otimes\mathrm{id}$ maps $\F_p^c\otimes U$ into $U$. So the restriction of $\iota$ to $\F_p^c\otimes U$ is
a parametrization with $V$ replaced by $U$ and $B_j$ by $B_j|_U$, and its common kernel is $\mathcal C\cap(\F_p^c\otimes U)$. Let $\Phi$ be the map on $\F_p^c\otimes U$ with
components $B_j\circ(\lambda_{j,k}\otimes\mathrm{id})$. These components map into $B_j(U)$, so $\rank\Phi\le\sum_jw_j\rank(B_j|_U)$. The
hypotheses of Theorem~\ref{thm:count} hold on $\ker\Phi$, with the injective map $x\mapsto c(\iota(x))$ on this kernel. The theorem gives
$\nu_N(p)\ge\dim\ker\Phi\ge c\dim U-\sum_jw_j\rank(B_j|_U)$.
\end{proof}

In the applications $c$ and the $w_j$ are affine functions of $N$ and $p$ on an interval of primes, while $V$ and
the $B_j$ are fixed. The space $\F_p^c$ is a space of polynomials of degree less than $c$ in the variable of a
decomposition modulo $p$. The forms $\lambda_{j,k}$ are Taylor coefficients at $X=-1,0,1$ and coefficients at $\infty$. The map $\iota$ places the data block by block along the rows. Its three properties are proved once for the whole family of parametrizations. On $\mathcal C$ the moment columns take the
values of a polynomial of degree less than $r$. Since the intervals lie in $(3N,16.2N)$, where $p>q$, we have $a=0$ in
Theorem~\ref{thm:count}.

\emph{Lower bounds for the maximum.} For a subspace $S$ of the dual $V^*$ let $U=S^\perp$. Let $R_j\subset V^*$ be the row space of $B_j$, and let $r_j\ge\dim R_j$ be real
numbers, in practice upper bounds for the ranks. Since $\rank(B_j|_U)=\dim R_j-\dim(R_j\cap S)$, the right side of Theorem~\ref{thm:everyU} is
\[
 E(S)=c(\dim V-\dim S)-\sum_jw_j\bigl(\dim R_j-\dim(R_j\cap S)\bigr).
\]
Let $e_0=c\dim V-\sum_jw_jr_j$ and $\delta_j=r_j-\dim R_j\ge0$. Then
\begin{equation}\label{eq:eulerid}
 E(S)=e_0+\sum_jw_j\delta_j+\sum_jw_j\dim(R_j\cap S)-c\dim S .
\end{equation}
A map may also be omitted from the sum if its weight is subtracted from $c$, because $\dim R_j-\dim(R_j\cap S)\le\dim V-\dim S$. The following three lemmas give a subspace
with $E(S)>e_0$, from relations among the row spaces. All sets of indices are subsets of $\mathcal J$, and all $w_j$ are nonnegative.

\begin{lemma}[pair]\label{lem:pairU}
Let $a\ne b$ lie in a set $I$ such that $R_a\le R_j$ or $R_b\le R_j$ for all $j\in I$. Let
$s\ge\dim(R_a+R_b)$ and $0\le h\le\min\bigl(w_a,w_b,\sum_{j\in I}w_j-c\bigr)$. Then some $S$ has
$E(S)\ge e_0+\max(0,r_a+r_b-s)\,h$.
\end{lemma}

\begin{proof}
Take $S=R_a\cap R_b$ and let $k=\dim S$. Then $S\le R_j$ for $j\in I$, so \eqref{eq:eulerid} gives
\[
 E(S)-e_0\ge w_a\delta_a+w_b\delta_b+\Bigl(\sum_{j\in I}w_j-c\Bigr)k\ge h(\delta_a+\delta_b+k).
\]
Since $\dim(R_a+R_b)+k=\dim R_a+\dim R_b$, we have $r_a+r_b-s\le\delta_a+\delta_b+k$. Also $\delta_a+\delta_b+k\ge0$.
\end{proof}

\begin{lemma}[several points]\label{lem:pointsU}
Let $c\ge0$. Let $I_1,\ldots,I_l$ be disjoint sets with $a_i\in I_i$, such that $R_{a_i}\le R_j$ for $j\in I_i$. Put $W_i=\sum_{j\in I_i}w_j$ and suppose that $W_i\le c$ for each $i$.
Let $s\ge\dim\sum_iR_{a_i}$. Let $h\ge0$ with $h\le w_{a_i}$ for each $i$ and $h\le\sum_iW_i-(l-1)c$. Then some $S$ has
$E(S)\ge e_0+\max(0,\sum_ir_{a_i}-s)\,h$.
\end{lemma}

\begin{proof}
Let $\mathsf T$ be the kernel of the sum map $\bigoplus_iR_{a_i}\to V^*$, and let $k=\dim\mathsf T$. Let $K_i\le R_{a_i}$ be the projection of $\mathsf T$ to the summand $i$, let $\kappa_i=\dim K_i$, and let
$S=\sum_iK_i$. Then $\kappa_i\le k$. The rank formula for the sum map gives $\sum_i\dim R_{a_i}\le s+k$. The kernel of the sum map on $\bigoplus_iK_i$ is again $\mathsf T$, so
$\dim S=\sum_i\kappa_i-k$.

For $j\in I_i$ we have $K_i\le R_j\cap S$, so $\dim(R_j\cap S)\ge\kappa_i$. Since the $I_i$ are disjoint, the $a_i$ are distinct, and \eqref{eq:eulerid} gives
\[
 E(S)-e_0\ge\sum_iw_{a_i}\delta_{a_i}+\sum_iW_i\kappa_i-c\sum_i\kappa_i+ck\ge h\sum_i\delta_{a_i}+\sum_i(W_i-c)\kappa_i+ck.
\]
Since $W_i-c\le0$ and $\kappa_i\le k$, the last two terms are at least $k(\sum_iW_i-lc+c)\ge hk$. So $E(S)-e_0\ge h(\sum_i\delta_{a_i}+k)$, and
$\sum_ir_{a_i}-s\le\sum_i\delta_{a_i}+k$.
\end{proof}

\begin{lemma}[kernel]\label{lem:kernelU}
Let $R\subset V^*$ be a subspace with $\dim R\le s$, and let $Z$ be a set with $R_j\le R$ for $j\in Z$. For every integer $d\ge0$
with $d+s\le\dim V$ some $S$ has
$E(S)\ge cd-\sum_{j\notin Z}w_j\min(r_j,d)$.
\end{lemma}

\begin{proof}
The space $R^\perp\subset V$ has dimension at least $d$. Choose $U\le R^\perp$ of dimension $d$ and let $S=U^\perp$. Then $B_j|_U=0$ for $j\in Z$, and $\rank(B_j|_U)\le\min(r_j,d)$ for
the other $j$.
\end{proof}

\emph{The rank data.} The space $V$ is a space of profiles of the coefficient vectors of a block, with coordinates
$T_k$. The conditions at the moment columns are linear combinations of a small number of basic rows
\[
 A[k]\mapsto T_k+T_{k+1},\qquad C[k]\mapsto T_{k+1}-T_k,
\]
together with the row of the centre coordinate. So the moment values themselves are not needed. Every row factors
through the map $G$ whose components are these basic rows. This gives $r_j\le\min(\rank F_j+\#\text{basic rows},\ \rank(G,F_j))$, where $F_j$ is the fixed part
of the map. Each rank bound is certified by an integer left inverse, checked by the Lean kernel. Integer left
inverses are valid at every prime, so no prime enters this step.

The $37$ intervals are treated as follows: $13$ by Lemma~\ref{lem:kernelU}, $12$ by Lemma~\ref{lem:pointsU}, $11$ by
Lemma~\ref{lem:pairU}, and one by $E(V)$ itself. On each interval the bound is an affine function
$\nu_N(p)\ge c_0N+c_1p-2$ with rational $c_0$ and $c_1$. The rounding of the integer orders is absorbed by
the shortening $\tau$ of the interval and the lowering of the bound by $10^{-6}$ in Section~\ref{ssec:kappa}.

\section{The bounds for the energy method}\label{app:field}

This appendix gives the trial sequences of Section~\ref{sec:energy}, the certified bounds that they satisfy, and
the bounds for the weights of Section~\ref{sec:real}.

\emph{Trial sequences.} For the constrained case we use a trial sequence with $40$ Chebyshev coefficients and a
finite number of geometric terms, with the multiplier $\lambda_1=1.68493357$. For the case with doubled
interaction we use trial sequences of the form \eqref{eq:trial}, with $\Kcoef$ Chebyshev coefficients, which are
dyadic rationals, and $\Ngeo$ geometric terms. Of these, $\Ncomplex$ have complex $z_j$ and complex $r_j$. We
place the complex charges at the $41$ equally spaced points $x_0=-0.16+0.008j$ $(0\le j\le40)$, all with
$\epsilon_0=10^{-2}$. The real terms sit at points $\pm z$ with $z$ increasing from $0.8$ to $0.999997$. The
numbers $z_j$ are rationals with at most twelve decimals. The norms of the geometric terms have closed forms:
two terms $(z_j,r_j)$ and $(z_l,r_l)$ contribute
$-\frac12\operatorname{Re}\bigl(r_jr_l\log(1-z_jz_l)+r_j\bar r_l\log(1-z_j\bar z_l)\bigr)$ to $\|u\|^2$.

\emph{The suprema of the fields.} We must bound two functions of one variable from above on their domains: the
field $X$ on $[-1,1]$ and the field $Y$ on $[0,1]$. We locate the real
critical points. The derivative is a rational function with a known denominator, and an exact identity gives its
numerator. The Descartes rule, after a M\"obius transformation of the interval, bounds the number of real roots
of the numerator. Sign evaluations at $81$-bit rational points then isolate this number of roots. One global Lipschitz
bound on each isolating interval and the values at the end points give the supremum.

\emph{The bounds.} For the base trial sequences this gives
\[
 \sup X\le\SupX,\quad \sup Y\le\SupY .
\]
So $M_X+M_Y=\FieldCap$, and the norms satisfy $2\|u\|^2+\frac12\|v\|^2\le0.7736206$. The first coefficient satisfies
$u_1\ge9/50$, which is used near $x=0$. For the constrained case, with the weights
$(\alpha_A,\beta_h,\gamma_H-\frac12,\lambda_1,2g_0/n_0)=(11/29,3/29,49/58,\lambda_1,5/29)$, the bounds
$\sup X_1\le-1.15854$ and $\sup Y\le-1.32981$ are obtained by isolating the zeros of the numerators of the
derivatives, which are polynomials with integer coefficients. The norm bound is $n_1=0.67761$.

\emph{The increments.} The increment of the case with doubled interaction has $40$ Chebyshev coefficients, and
that of the constrained case has $24$. Both are dyadic. Their fields are bounded on the whole domains by the
same methods. For the case with doubled interaction the norm bound at $\tau$ is $0.7736206-0.117647903\,\tau+0.005215061\,\tau^2$,
the field $X$ gains at most $0.01484272\,\tau$, and the column field $Y$ with $\gamma g$ is bounded by a piecewise
linear function of $\gamma$ with knots at $\alpha_1,\frac15,\frac3{10},\frac25,\frac35,\frac45,1$. The fields $F_1$ and $F_2$
contain the node term $\hat\delta/58$. They are bounded by covers by cells, with exact rational and integer arithmetic on each cell
and the maximum of $\hat\delta$ on each cell. For $F_1$ on $|x|\ge401/2500$ the cover has $150$ cells. It uses a Taylor expansion of the
polynomial part and an expansion to order $5$ of the logarithmic terms. On the positive part of $E_W$ below $401/2500$ six
further cells suffice. For $F_2$ the covers of $[98749/10^6,\,9995/10^4]$ have $17$ and $14$ cells at the two ends of the move. This gives
\[
 \sup F_1\le0.007,\qquad \sup F_2\le0.007-\log2
\]
at $\tau=0$, and $-0.0268$ and $0.0233-\log2$ at $\tau=1$. For the constrained case the norm bound at $\tau$ is
$0.67761-0.082834152\,\tau+0.004096781\,\tau^2$, and the bound $\mathcal M_1$ of Section~\ref{sec:energy} is the chord of
$0.35072$ and $0.35295$. The two numbers are maxima over the $865$ segments of the chord table at the two end values of
$\lambda_1$. With $a_0=41/58$, $a_1=76/29$ and $U_0=\Uzero$ these bounds give the inequalities of \S\ref{ssec:cases}.
Each check $V(\alpha)\le\sigma\alpha$ is a polynomial inequality on a closed interval and is proved by a cover in exact
rational arithmetic.

\emph{The bound for the weights.} Inequalities \eqref{eq:circlecap} and \eqref{eq:segcap} concern the explicit function $L$ of
\eqref{eq:Ldef}. On the circle $|z|=\frac34$ write $z=\frac34(s+i\sqrt{1-s^2})$ with $s=\cos\theta$. Then $|z-t|^2=\frac9{16}+t^2-\frac32ts$
is affine in $s$ for every real $t$. So $\log|1-z|$ is a concave function of $s$. Each $G_c(z)$ is $\operatorname{arccosh}$ of
$(|z-a_c|+|z-b_c|)/w_c$, a sum of square roots of affine functions of $s$, and therefore concave in $s$ as well. Hence the left
side of \eqref{eq:circlecap} is a concave function of $s\in[-1,1]$. Its maximum is $3.25053\ldots$ near $s=-0.1741$. A concave function
is bounded by its secants outside the points and by the chord between them, so three points at distance $\frac1{2000}$ around the
maximum suffice. At these points the value is enclosed to within $10^{-17}$ by exact rational arithmetic with integer square roots,
and the bound $D_\Gamma=\Dcirc$ holds with a margin $6\cdot10^{-5}$. On the real segments we use, for real $x$, the closed form
$e^{G_c(x)}w_c/4\le\frac12\max(w_c,|2x-a_c-b_c|)$ together with $6\log|1-x|\le6\log\frac14$ for $x\ge\frac34$ and
$6\log|1-x|\le6\log2$ for $x\le-\frac34$. This gives \eqref{eq:segcap} with values below $-9$, far below $D_\Gamma$.
The knot values $v_i$ of the chord table lie between $-1.97$ and $17.99$.

\section{The reference measures of Section~\ref{sec:signed}}\label{app:trials}

The bound of \S\ref{ssec:budget} uses reference measures $q_+$ and $q_-$, a row reference $R$ and a column reference $C$.
This appendix describes them and the certified bounds \eqref{eq:caps}.

\emph{The measures.} On the real line of the variable $Y=x+1/x$ with $|Y|\ge2$ we take measures $G_+$ and $G_-$ that
are piecewise uniform on finitely many cells in $Y$, with masses $m_+$ and $m_-$ and total mass $1$. They carry the positive and the
negative free nodes, and $R=G_++G_-$ is the row reference. The column reference $C$ has mass $1$. The measures are completed by
auxiliary measures on the $16$ inner circles. Each auxiliary measure has a cubed Poisson density: its Fourier coefficients
are $-q^k(1+Ak+Bk^2)/k$ with $q=e^{-\varepsilon/2}$ and $\varepsilon=27/250$, and the constants $A$ and $B$ are rational functions of $r=q^2$.
The auxiliary measures $\nu_+$ and $\nu_-$ sit at the phases $\pi j/93$, with $616$ and $935$ entries respectively, and have total masses $1-m_+$ and $\frac{53}{58}-m_-$. Let $X$ be the inverse of $x\mapsto x+1/x$ with $|X|\le1$. Then
$q_+=X_\#G_++\nu_+$ is the measure in the plane of $z$. In the plane of $u$ we take $q_-=J_\#\tilde q_-$ with $\tilde q_-=X_\#G_-+\nu_-$.

\emph{The fields.} Let $L(\mu)=\int\log|X|\,d\mu$. For $\sigma=\pm$ the free row field is
\begin{multline*}
 F_\sigma(x)=L(R)+L(G_\sigma)-\kappa_\sigma+U^{G_\sigma+R-2_\#C}(x+1/x)+\bigl(m_\sigma+\tfrac{11}{29}\bigr)\log|x|\\
 -\tfrac{39}{29}\log(1+x^2)+\tfrac3{29}\log(1-x)+[\sigma=+]\bigl(\mathsf P(x)+\mathsf E\bigr)+A_\sigma(x),
\end{multline*}
where $A_+=U^{\nu_+}$ and $A_-(x)=\int(\log|x-w|+\log|1-xw|)\,d\nu_-(w)$. The polynomial $\mathsf P(x)=\sum_{k\le900}m_kx^k/k$ with $m_k=\int X^k\,dG_+$ is
the truncated reflection potential of the positive measure, and $\mathsf E\ge0$ bounds its tail.
The column field is
\[
 \mathsf C(s)=-L(R)+2U^C(s)-U^R(2s)+\tfrac9{58}\log s+\tfrac3{58}\log(1-s).
\]
A prescribed row outside $S$ has the field $\mathsf H_\sigma$ obtained from $F_\sigma$ by three changes. The references $R$ and $C$ are removed, the
coefficients $\frac{11}{29}$ and $\frac3{29}$ are replaced by $0$, and $\frac{39}{29}$ is replaced by $\frac{53}{58}$. Its domain is
$x\le-19501/10^5$ and $x\ge98749/10^6$, which contains $E_x$. A prescribed row in $S$ has the field
\[
 \mathsf H_S(y)=\tfrac{48}{58}\log y-\tfrac{53}{58}\log(1+y^2)-\kappa_+-\mathsf P(y)+\int\log|1-yw|\,d\nu_+(w)
\]
for $0<y<1$.
The angular fields are the same potentials on the $16$ circles, together with the terms $(53\sigma_j-\tau_j/2)/58$.

\emph{The constant $\mathsf I$.} Let $E(\mu,\nu)=\iint\log|Y-Y'|\,d\mu\,d\nu$. Evaluating $\mathsf Q$ on the references gives $\mathsf I_0=-\frac12E(R,R)-E(C,C)+E(R,2_\#C)$. Since $R=G_++G_-$, six pair
energies of lists of cells occur: those of $G_+$, $G_-$ and $C$ with themselves, $E(G_+,G_-)$, $E(G_+,2_\#C)$ and $E(G_-,2_\#C)$. The energies $I(q_\pm)$ split into cell, cross and auxiliary parts. Their cell parts contain the same
pair energies. For piecewise uniform measures a pair energy is a bilinear form in the jumps of the densities, with the kernel $t^2\log|t|$, up to constants. We evaluate it on integer end points in fixed point arithmetic. Well separated blocks of end points are
treated by a series in the offsets from their centres, using the moments of the blocks, and the other pairs by a recursion on the blocks. This encloses each pair energy to within $10^{-7}$. The remaining terms of $I(q_\pm)$ are the logarithmic moments of the cells,
the integrals of the cells against the rings, which are sums of primitives over the end points, and the quadratic forms of the auxiliary measures, which have closed forms on the rings. All of these are enclosed to within $10^{-9}$. The reflection term $-\iint\log(1-xx')$ of $I(q_+)$
enters only through a lower bound from its truncated series. Finally $\int_0^{53}\tau^*\ge1425.3339$. Together these give $\mathsf I_0+\mathsf I_{\mathrm{gas}}\le2.366997$.

\emph{The column field.} The three bounds $c^{\mathsf C}_i$ of \eqref{eq:caps} on $(0,1)$ use a chain of $386$ units $s=(M+\delta z)/10^{12}$ with $|z|\le1$, together with two end ranges. On a unit the field is at most an integer polynomial in $z$ plus an explicit
error. The far cells of the references enter through exact integer moments of contiguous clusters, $1152$ clusters in all, and the Taylor series of the logarithmic kernel. The near cells of $C$ enter through the closed form of their potential. The two logarithms enter by
series or, on the end units, by monotonicity. The constant is $L(R)\ge-0.6905432$, from the primitive of $\log|X|$ in $Y$.

\emph{The angular bounds.} The bounds $\kappa_\pm$ are proved for both signs on each of the $16$ radii, by $32$ certificates. The auxiliary part is cut at degree $384$, and its tail is bounded. The gas is split into bins. On a bin the kernel has the form
$\mathsf Q_u(x)=\mathsf a-2\mathsf bux+\mathsf cx^2$ with $u=\cos\theta$. Jensen's inequality and an odd truncation of the logarithmic series bound $\int\frac12\log\mathsf Q_u$ by a polynomial in $u$. The negative sign has a direct and a reflected kernel. The sum over the bins
is converted exactly to a Chebyshev series. Its bound by $\kappa_\pm$ on $[-1,1]$ is certified by an identity of the same form as the one for the free rows below.

\emph{The prescribed rows.} The auxiliary series of $\mathsf H_\pm$ are cut at degree $384$. The gas potential is expanded in Chebyshev polynomials on $[0,1]$. The first $64$ terms are exact, and the tail is bounded by Bessel's inequality from a bound for the $L^2$ energy of the cell densities.
The reflection series is dropped by sign for the positive gas and kept through exact moments for the negative gas. The radial field is a sum of $16$ blocks, one for each radius. Its polynomial part goes exactly into the Chebyshev variable, and the part of degree at most $16$ is bounded by
a chord cover with the radii as nodes. This gives a polynomial majorant, and a sum of squares identity gives $\mathsf H_\pm\le-0.9783$. On $[\frac1{10},1]$ the field $\mathsf H_S$ is treated in the same way. There $\mathsf P$ is replaced by its partial sum of degree $32$, since the moments are nonnegative, the auxiliary series is cut at degree $32$ with an absolute tail, and
$\log y$ and $\log(1+y^2)$ are expanded explicitly. This gives $\mathsf H_S\le-0.997$. On $(0,\frac1{10}]$ one has directly $\mathsf H_S(y)\le\frac{24}{29}\log y+\log(1+y)-\kappa_++$(a fixed tail), which is smaller. So $\mathsf M=-0.976$ is admissible.

\emph{The free rows.} The free row fields are bounded on each sign as follows. Put $t=|x|$ and $z=2t-1$. The exact field is at most its truncation to degree $384$ in the auxiliary terms, up to a tail
below $10^{-8}$. We bound the truncated field by one polynomial in $z$. The direct potential of the own measure, $\int\log|t-X(y)|$, is replaced by its Chebyshev series,
with a tail bounded through the variation of the density. Thin pieces of nonnegative density are dropped, since $\log|t-v|\le0$ there. All other potentials
are bounded from above by finite sums of terms $w\log|1-r(\mathsf q)t|$ with rational $\mathsf q$: tangent majorants for the reflection and opposite pieces, chord majorants for the
column panels, and the same form for $\log t$ and $\log(1-t)$. Each such term has an explicit Chebyshev expansion with a geometric tail. The moments of these terms and the
direct moments are computed by exact integer arithmetic with Kronecker packing. The radial field and the term $-\frac{39}{29}\log(1+t^2)$ are bounded by a chord cover and a
Chebyshev polynomial. This gives a polynomial $\mathsf p(z)$ with $F_\sigma\le\mathsf p$ plus explicit small errors. Finally $\mathsf p(z)\le\mathsf c$ on $[-1,1]$ is proved by an exact identity
$\mathsf d(\mathsf c-\mathsf p)=\mathsf r+\mathsf q_0^2$ between Chebyshev series, with $\mathsf q_0$ of degree $16383$ and with $\mathsf r$ small in the $\ell^1$ norm of its coefficients. This gives
the bounds $c^{\mathsf R}_0$ and $c^{\mathsf R}_1$ of \eqref{eq:caps} at the two ends of the move, for both signs.

\providecommand{\bysame}{\leavevmode\hbox to3em{\hrulefill}\thinspace}
\providecommand{\MR}{\relax\ifhmode\unskip\space\fi MR }
% \MRhref is called by the amsart/book/proc definition of \MR.
\providecommand{\MRhref}[2]{%
  \href{http://www.ams.org/mathscinet-getitem?mr=#1}{#2}
}
\providecommand{\href}[2]{#2}
\begin{thebibliography}{10}

\bibitem{Apery79}
Roger Ap\'ery, \emph{Irrationalit\'e de {$\zeta(2)$} et {$\zeta(3)$}},
  Ast\'erisque \textbf{61} (1979), 11--13.

\bibitem{Beukers79}
Frits Beukers, \emph{A note on the irrationality of {$\zeta(2)$} and
  {$\zeta(3)$}}, Bull. London Math. Soc. \textbf{11} (1979), no.~3, 268--272.

\bibitem{Bojnal2026}
Hrishik~Sai Bojnal, \emph{An irrationality exponent below 2402 for {C}atalan's
  constant}, GitHub repository, October 15, 2026,
  \url{https://github.com/proof-factory/catalan-exponent}; design search and
  certification by Claude Fable~5.1 (Anthropic), 2026.

\bibitem{CDT2024}
Frank Calegari, Vesselin Dimitrov, and Yunqing Tang, \emph{The linear
  independence of {$1$}, {$\zeta(2)$}, and {$L(2,\chi_{-3})$}}, 2024,
  \href{https://arxiv.org/abs/2408.15403}{arXiv:2408.15403}.

\bibitem{CDT2025}
\bysame, \emph{Arithmetic holonomy bounds and effective {D}iophantine
  approximation}, 2025, Theorem 5.1 gives $\mu(L(2,\chi_{-3}))\le 24781$,
  \href{https://arxiv.org/abs/2510.04156}{arXiv:2510.04156}.

\bibitem{Chudnovsky26}
Gregory~V. Chudnovsky, \emph{The irrationality exponent of the lemniscate
  period: {A} determinant proof with a {L}ean 4 formalization}, 2026, Preprint,
  19 October 2026, \href{https://arxiv.org/abs/2610.26539}{arXiv:2610.26539}.

\bibitem{Carreira2026b}
Marcos Costa Santos~Carreira, \emph{Determinant constructions for the
  {D}irichlet beta value at four}, Preprint, 9 October 2026,
  \href{https://doi.org/10.5281/zenodo.23326647}{doi:10.5281/zenodo.23326647},
  2026.

\bibitem{Carreira2026}
\bysame, \emph{General designs and an irrationality exponent for {C}atalan's
  constant}, Preprint, 9 October 2026, revised 21 October 2026,
  \href{https://doi.org/10.5281/zenodo.23532239}{doi:10.5281/zenodo.23532239};
  sources at \url{https://github.com/MarcosCarreira/catalan-designs}, 2026.

\bibitem{Davis2026}
Cinna Davis, \emph{Catalan's constant: irrationality exponent research}, GitHub
  repository, \url{https://github.com/1337isnot1337/catalan-constant}; Lean
  verification at \url{https://github.com/bigospeedrun/verifier/issues/87},
  October 31, 2026.

\bibitem{Lean4}
Leonardo de~Moura and Sebastian Ullrich, \emph{The {L}ean 4 theorem prover and
  programming language}, Automated Deduction -- CADE 28, Lecture Notes in
  Computer Science, vol. 12699, Springer, 2021, pp.~625--635.

\bibitem{Eskandari2026}
Payman Eskandari, \emph{Improved rational approximations to {C}atalan's
  constant}, 2026, Preprint, 22 September 2026,
  \href{https://arxiv.org/abs/2609.26354}{arXiv:2609.26354}.

\bibitem{EMN2025}
Payman Eskandari, V.~Kumar Murty, and Yusuke Nemoto, \emph{Mixed motives and
  linear forms in the {C}atalan constant}, J. Lond. Math. Soc. (2) \textbf{114}
  (2026), no.~4, Paper No. e70708,
  \href{https://arxiv.org/abs/2510.20648}{arXiv:2510.20648}.

\bibitem{Fasten26}
Azriel Fasten, \emph{Irrationality of the {R}iemann zeta value at five}, Lean
  formalization, Palomar registry entry PALOMAR-2026-09-25-000018,
  \url{https://github.com/realazthat/zeta5-lean2}, commit \texttt{2fa4e82},
  2026.

\bibitem{Fauzan2026}
Aabir Fauzan, \emph{{$\zeta(5)$} is irrational}, Preprint v1, September 17,
  2026, Zenodo,
  \href{https://doi.org/10.5281/zenodo.22826419}{doi:10.5281/zenodo.22826419},
  2026.

\bibitem{Fischler20}
St\'ephane Fischler, \emph{Irrationality of values of {$L$}-functions of
  {D}irichlet characters}, J. Lond. Math. Soc. (2) \textbf{101} (2020), no.~2,
  857--876.

\bibitem{Forrester2018Andreief}
Peter~J. Forrester, \emph{Meet {A}ndr\'eief, {B}ordeaux 1886, and {A}ndreev,
  {K}harkov 1882--83}, 2018,
  \href{https://arxiv.org/abs/1806.10411}{arXiv:1806.10411}.

\bibitem{Giraudin26}
David~Niedbala Giraudin, \emph{A new upper bound for the irrationality exponent
  of {$\zeta(2)$}}, 2026, Preprint, 2 October 2026,
  \href{https://arxiv.org/abs/2610.02912}{arXiv:2610.02912}.

\bibitem{Giraudin26z3}
\bysame, \emph{A new upper bound for the irrationality exponent of
  {$\zeta(3)$}}, 2026, Preprint, 22 September 2026,
  \href{https://arxiv.org/abs/2609.26980}{arXiv:2609.26980}.

\bibitem{Giraudin26z5}
\bysame, \emph{A new upper bound for the irrationality exponent of
  {$\zeta(5)$}}, 2026, Preprint, 1 November 2026,
  \href{https://arxiv.org/abs/2611.01708}{arXiv:2611.01708}.

\bibitem{Hata93}
Masayoshi Hata, \emph{Rational approximations to {$\pi$} and some other
  numbers}, Acta Arith. \textbf{63} (1993), no.~4, 335--349.

\bibitem{Jha2026}
R.~K. Jha, \emph{A finite irrationality measure for {C}atalan's constant},
  Preprint v1, October 8, 2026,
  \href{https://doi.org/10.5281/zenodo.23237269}{doi:10.5281/zenodo.23237269};
  Lean sources at \url{https://github.com/rohitkrjha/catalan-measure}, 2026.

\bibitem{Arb}
Fredrik Johansson, \emph{Arb: efficient arbitrary-precision midpoint-radius
  interval arithmetic}, IEEE Trans. Comput. \textbf{66} (2017), no.~8,
  1281--1292, \href{https://arxiv.org/abs/1611.02831}{arXiv:1611.02831}.

\bibitem{KrattRivoal08}
Christian Krattenthaler and Tanguy Rivoal, \emph{On a linear form for
  {C}atalan's constant}, South East Asian J. Math. Math. Sci. \textbf{6}
  (2008), no.~2, 3--15,
  \href{https://arxiv.org/abs/0810.1927}{arXiv:0810.1927}.

\bibitem{KrattZudilin19}
Christian Krattenthaler and Wadim Zudilin, \emph{Hypergeometry inspired by
  irrationality questions}, Kyushu J. Math. \textbf{73} (2019), no.~1,
  189--203, \href{https://arxiv.org/abs/1802.08856}{arXiv:1802.08856}.

\bibitem{Krisman26}
Shawn Krisman, \emph{One of {$\beta(4),\beta(6),\ldots,\beta(36)$} is
  irrational}, Hexagon preprint hexagon:2610.00285v2, October 16, 2026,
  \url{https://hexagonmath.org/2610.00285}; Lean sources at
  \url{https://github.com/telaviv/beta-irrationality}, 2026.

\bibitem{KumarTosi26}
Veekesh Kumar and Riccardo Tosi, \emph{Irrationality of values of
  {$L$}-functions for real {D}irichlet characters}, 2026, Preprint, 30 October
  2026, \href{https://arxiv.org/abs/2610.41458}{arXiv:2610.41458}.

\bibitem{Lai2026}
Li~Lai, \emph{The linear independence of {$1$}, {$\zeta(3)$}, and
  {$\zeta(4)$}}, 2026, Preprint, 5 October 2026 (v2: 15 October 2026),
  \href{https://arxiv.org/abs/2610.06838}{arXiv:2610.06838}.

\bibitem{Comparator}
{Lean FRO}, \emph{Comparator: checking a {L}ean proof against a trusted
  statement}, \url{https://github.com/leanprover/comparator}, 2026, Software;
  accessed October 2026.

\bibitem{LBKWSK26}
Elyasheev Leibtag, Hila Barkan, Rotem Kalisch, Shachar Weinbaum, Michael
  Shalyt, and Ido Kaminer, \emph{New directions in irrationality}, 2026,
  Preprint, 20 October 2026,
  \href{https://arxiv.org/abs/2610.27799}{arXiv:2610.27799}; supplement and
  code at
  \url{https://github.com/RamanujanMachine/new-directions-in-irrationality}.

\bibitem{Litsey2026b}
Sean Litsey, \emph{Irrationality measures for {$L(2,\chi_{-3})$} and other
  differences of trigamma values}, Hexagon preprint hexagon:2610.00401,
  \url{https://hexagonmath.org/2610.00401}, 2026.

\bibitem{LiuJiangZhang26}
Jingwen Liu, Kai Jiang, and Pingwen Zhang, \emph{Irrationality exponents of
  logarithms of positive rational numbers}, 2026, Preprint, 7 October 2026,
  revised 5 November 2026 (v3),
  \href{https://arxiv.org/abs/2610.10192}{arXiv:2610.10192}.

\bibitem{Marcovecchio09}
Raffaele Marcovecchio, \emph{The {R}hin--{V}iola method for {$\log 2$}}, Acta
  Arith. \textbf{139} (2009), no.~2, 147--184.

\bibitem{MultiAgency2026}
{MultiAgency swarm}, \emph{Catalan: {$\mu(G)\le528.2315$}}, GitHub repository,
  \url{https://github.com/darkhipo/catalan-constant-gen5}; verifier submission
  \url{https://github.com/bigospeedrun/verifier/issues/103}, November 2, 2026.

\bibitem{Nesterenko16}
Yu.~V. Nesterenko, \emph{On {C}atalan's constant}, Proc. Steklov Inst. Math.
  \textbf{292} (2016), no.~1, 153--170.

\bibitem{OpenAI2026}
{OpenAI}, \emph{Catalan's constant is irrational}, Preprint dated September 24,
  2026, released October 6, 2026, with a Lean formalization, available at
  \href{https://github.com/openai/math/blob/fd4aeeb2ee4fc729c18d98444fed42fd0529eeeb/preprints/Catalans-constant-is-irrational-September-24-2026/paper.pdf}{the
  original paper PDF} and
  \href{https://github.com/openai/math/tree/fd4aeeb2ee4fc729c18d98444fed42fd0529eeeb/lean/OAI/NumberTheory/Catalan}{the
  Lean formalization} (repository commit \texttt{fd4aeeb}, October 7, 2026),
  2026.

\bibitem{OpenAIPi2026}
\bysame, \emph{The irrationality exponent of {$\pi$} is 2}, Preprint dated
  September 24, 2026, released October 6, 2026,
  \url{https://github.com/openai/math/tree/main/preprints/The-irrationality-exponent-of-pi-is-2-September-24-2026},
  2026.

\bibitem{Prellberg26}
Thomas Prellberg, \emph{Three irrationality results for the dilogarithm:
  {R}hin--{V}iola and {V}iola--{Z}udilin constructions at {$-1/4$}, {$1/5$},
  and {$-1/3$}}, 2026, Preprint, 7 September 2026,
  \href{https://arxiv.org/abs/2609.06895}{arXiv:2609.06895}.

\bibitem{RhinViola96}
Georges Rhin and Carlo Viola, \emph{On a permutation group related to
  {$\zeta(2)$}}, Acta Arith. \textbf{77} (1996), no.~1, 23--56.

\bibitem{RhinViola01}
\bysame, \emph{The group structure for {$\zeta(3)$}}, Acta Arith. \textbf{97}
  (2001), no.~3, 269--293.

\bibitem{Rivoal06}
Tanguy Rivoal, \emph{Nombres d'{E}uler, approximants de {P}ad\'e et constante
  de {C}atalan}, Ramanujan J. \textbf{11} (2006), no.~2, 199--214.

\bibitem{RivoalZudilin2003}
Tanguy Rivoal and Wadim Zudilin, \emph{Diophantine properties of numbers
  related to {C}atalan's constant}, Math. Ann. \textbf{326} (2003), no.~4,
  705--721.

\bibitem{SaffTotik}
Edward~B. Saff and Vilmos Totik, \emph{Logarithmic potentials with external
  fields}, Grundlehren der mathematischen Wissenschaften, vol. 316, Springer,
  Berlin, 1997.

\bibitem{Salikhov08}
V.~Kh. Salikhov, \emph{On the irrationality measure of the number {$\pi$}},
  Russian Math. Surveys \textbf{63} (2008), no.~3, 570--572.

\bibitem{Sun2026}
Zhi-Wei Sun, \emph{Catalan's constant is irrational}, 2026,
  \href{https://arxiv.org/abs/2609.04176}{arXiv:2609.04176}.

\bibitem{Mathlib}
{The mathlib Community}, \emph{The {L}ean mathematical library}, Proceedings of
  the 9th ACM SIGPLAN International Conference on Certified Programs and Proofs
  (CPP 2020), 2020, pp.~367--381.

\bibitem{PNTPlus}
{The PrimeNumberTheoremAnd contributors}, \emph{The {PNT+} project: a {L}ean
  formalization of the prime number theorem and related results},
  \url{https://github.com/AlexKontorovich/PrimeNumberTheoremAnd}, The vendored
  compatibility additions used here are recorded in the repository's
  \texttt{NOTICE} and \texttt{PORTING.md}.

\bibitem{Wachs2026}
Dane Wachs, \emph{A note on a recent claimed proof of the irrationality of
  {C}atalan's constant}, 2026,
  \href{https://arxiv.org/abs/2609.22339}{arXiv:2609.22339}.

\bibitem{Xu2026}
Ryan Xu, \emph{Improved {$\kappa=2$} trial: {$\mu(G)\le727.03$}}, GitHub issue,
  October 22, 2026,
  \url{https://github.com/slitsey/catalan-irrationality-measure/issues/5}; data
  and scripts at \url{https://github.com/ry-xu/slightlysupervised}, 2026.

\bibitem{ZeilbergerZudilin20}
Doron Zeilberger and Wadim Zudilin, \emph{The irrationality measure of {$\pi$}
  is at most {$7.103205334137\ldots$}}, Mosc. J. Comb. Number Theory \textbf{9}
  (2020), no.~4, 407--419,
  \href{https://arxiv.org/abs/1912.06345}{arXiv:1912.06345}.

\bibitem{Zudilin02Cheb}
Wadim Zudilin, \emph{A few remarks on linear forms involving {C}atalan's
  constant}, Chebyshevskii Sb. \textbf{3} (2002), no.~2(4), 60--70, In Russian,
  \href{https://arxiv.org/abs/math/0210423}{arXiv:math/0210423}.

\bibitem{Zudilin2003}
\bysame, \emph{An {A}p\'ery-like difference equation for {C}atalan's constant},
  Electron. J. Combin. \textbf{10} (2003), Research Paper 14,
  \href{https://arxiv.org/abs/math/0201024}{arXiv:math/0201024}.

\bibitem{Zudilin14}
\bysame, \emph{Two hypergeometric tales and a new irrationality measure of
  {$\zeta(2)$}}, Ann. Math. Qu\'ebec \textbf{38} (2014), no.~1, 101--117,
  \href{https://arxiv.org/abs/1310.1526}{arXiv:1310.1526}.

\bibitem{Zudilin}
\bysame, \emph{A determinantal approach to irrationality}, Constr. Approx.
  \textbf{45} (2017), no.~2, 301--310,
  \href{https://arxiv.org/abs/1507.05697}{arXiv:1507.05697}.

\bibitem{Zudilin19Survey}
\bysame, \emph{Arithmetic of {C}atalan's constant and its relatives}, Abh.
  Math. Semin. Univ. Hambg. \textbf{89} (2019), no.~1, 45--53,
  \href{https://arxiv.org/abs/1804.09922}{arXiv:1804.09922}.

\end{thebibliography}

\end{document}
